\documentclass[12pt, ukrainian]{article}
\usepackage[a4paper]{geometry}
\setlength\parindent{0mm}
\geometry{top=1.5cm,bottom=1.5cm,left=1.3cm,right=1.5cm}
\pagestyle{empty}
\usepackage[T2A]{fontenc}
\usepackage[utf8]{inputenc}
\usepackage[ukrainian]{babel}
\usepackage{graphicx}
\usepackage{caption}
\usepackage{caption}
\DeclareCaptionFormat{nocaption}{}
\captionsetup{format=nocaption,aboveskip=0pt,belowskip=0pt}
\usepackage[Export]{adjustbox} 
\adjustboxset{max size={0.9\linewidth}{0.9\paperheight}}
\usepackage{verbatim, listings, clrscode3e, fancyvrb}
\def\code#1{\Verb!#1!}
\begin{document}
%begin
%begin
\begin {large}

\textbf{22.11.2024 Контрольна робота, варіант  \No { 1 }}\par


{
%beginitem
1. По яких днях тижня відбувались лекції з предмету?

%enditem


\vspace{5mm}
%beginitem
2. Як зовуть (ім'я, по батькові) викладача, що вів практичні заняття?

%enditem


\vspace{5mm}
%beginitem
3. Нехай int var=13; Один потік збільшив змінну var на 3, а інший збільшив її в три рази.

гонка даних (data race), стан гонки (race condition)

Виберіть усі правильні твердження (якщо такі є).\par
\vspace{2mm}
( 1 ) Тут може не бути стану гонки.%bad

( 2 ) Тут гарантовано нема гонки даних, але може бути стан гонки.%bad

( 3 ) Тут може бути тільки стан гонки.%bad

( 4 ) Тут гарантовано є стан гонки, але гонки даних може не бути.%good
%enditem


\vspace{5mm}
%beginitem
4. Виберіть усі правильні твердження (якщо такі є).\par
\vspace{2mm}

( 1 ) Клас thread не має засобів, щоб повернути результат виконання запущеного потоку.%good

( 2 ) Розблоковувати незаблокований м’ютекс можна.%bad

( 3 ) Секція є критичною, якщо під час її виконання не можна генерувати винятки.%bad

( 4 ) Функція async запускає виконання в новому потоці.%bad
%enditem


\vspace{5mm}
%beginitem
5. Виберіть усі правильні твердження (якщо такі є).\par
\vspace{2mm}

( 1 ) За використання структурою даних рівно одного м’ютексу, взаємоблокування на її операціях неможливі.%good

( 2 ) Операція над структурою даних без блокувань завжди виконається за обмежений час.%bad

( 3 ) Реалізації різних операцій над структурою даних повинні використовувати різні м’ютекси.%bad

( 4 ) На операціях структури даних без блокувань взаємоблокувань потоків виникнути не може.%bad
\newpage

%enditem


\vspace{5mm}
%beginitem
6. Виберіть усі правильні твердження (якщо такі є).\par
\vspace{2mm}

( 1 ) Потік мови С++ не може отримати доступ до даних іншого потоку.%bad

( 2 ) Потік може мати доступ до автоматичних змінних іншого потоку, але це небезпечно.%good

( 3 ) Потік мови С++ не може отримати доступ до автоматичних змінних іншого потоку.%bad

( 4 ) Автоматичні змінні не можуть розділятися між потоками.%bad
%enditem


\vspace{5mm}
%beginitem
7. У цьому фрагменті коду
\begin{verbatim}
template <class T> void f(T a){}

void g(){
    int i=5;
    f(unique_ptr<int>(new int(i)))+(unique_ptr<int>(new int(35)));
\end{verbatim}
( 1 ) можливий виток ресурсів%good

( 2 ) є невизначена поведінка%bad

( 3 ) є гонка даних%bad

( 4 ) є стан гонки%bad
%enditem
}
%end
\end {large}

\vspace{4mm}
%end
\newpage
%begin
%begin
\begin {large}

\textbf{22.11.2024 Контрольна робота, варіант  \No { 2 }}\par


{
%beginitem
1. По яких днях тижня відбувались лекції з предмету?

%enditem


\vspace{5mm}
%beginitem
2. Як зовуть (ім'я, по батькові) викладача, що вів практичні заняття?

%enditem


\vspace{5mm}
%beginitem
3. Нехай int var=13; Один потік збільшив змінну var на 3, а інший збільшив її в три рази.

гонка даних (data race), стан гонки (race condition)

Виберіть усі правильні твердження (якщо такі є).\par
\vspace{2mm}
( 1 ) Тут гарантовано нема стану гонки, але може бути гонка даних.%bad

( 2 ) Тут гарантовано є стан гонки.%good

( 3 ) Тут гарантовано нема ані гонки даних, ані стану гонки.%bad

( 4 ) Тут гарантовано є і стан гонки, і гонка даних.%bad
%enditem


\vspace{5mm}
%beginitem
4. Виберіть усі правильні твердження (якщо такі є).\par
\vspace{2mm}

( 1 ) Текст однопотокової програми мовою С++ однозначно визначає послідовність виконуваних обчислень.%bad

( 2 ) Розблокувати м’ютекс може тільки той потік, що його заблокував.%good

( 3 ) Значення, що є результатом виконання функції, запущеної через async, зберігається безпосередньо в об’єкті ф’ючерса (future).%bad

( 4 ) Об’єкти класів promise та future можна використовувати тільки в багатопотокових програмах.%bad
%enditem


\vspace{5mm}
%beginitem
5. Виберіть усі правильні твердження (якщо такі є).\par
\vspace{2mm}

( 1 ) Одночасно використовувати атомарних змінних більше, ніж регістрів процесора, у програмі не можна.%bad

( 2 ) Реалізація структури даних з блокуваннями може використовувати атомарні змінні.%good

( 3 ) Вільних від очікувань структур даних не існує.%bad

( 4 ) Операція над вільною від блокувань структурою даних завжди виконається за обмежений час.%bad
\newpage

%enditem


\vspace{5mm}
%beginitem
6. Виберіть усі правильні твердження (якщо такі є).\par
\vspace{2mm}

( 1 ) Динамічні змінні не можуть розділятися між потоками.%bad

( 2 ) Потік має доступ тільки до своїх власних даних та до глобальних статичних даних.%bad

( 3 ) Потік може мати доступ до автоматичних змінних іншого потоку і це безпечно.%bad

( 4 ) Кожен потік програми може мати доступ до динамічних даних програми.%good
%enditem


\vspace{5mm}
%beginitem
7. У цьому фрагменті коду
\begin{verbatim}
template <class T> void f(T a1, T a2){}

void g(){
    int i=5;
    f(unique_ptr<int>(new int(i++)), unique_ptr<int>(new int(++i)));
\end{verbatim}

( 1 ) є невизначена поведінка%good

( 2 ) можливий виток ресурсів%bad

( 3 ) є стан гонки%bad

( 4 ) є гонка даних%bad
%enditem
}
%end
\end {large}

\vspace{4mm}
%end
\newpage
%begin
%begin
\begin {large}

\textbf{22.11.2024 Контрольна робота, варіант  \No { 3 }}\par


{
%beginitem
1. По яких днях тижня відбувались лекції з предмету?

%enditem


\vspace{5mm}
%beginitem
2. Як зовуть (ім'я, по батькові) викладача, що вів практичні заняття?

%enditem


\vspace{5mm}
%beginitem
3. Нехай int var=13; Один потік збільшив змінну var на 3, а інший збільшив її в три рази.

гонка даних (data race), стан гонки (race condition)

Виберіть усі правильні твердження (якщо такі є).\par
\vspace{2mm}
( 1 ) Тут гарантовано є гонка даних.%bad

( 2 ) Тут може бути тільки гонка даних.%bad

( 3 ) Тут може не бути ані гонки даних, ані стану гонки.%bad

( 4 ) Тут може не бути гонки даних.%good
%enditem


\vspace{5mm}
%beginitem
4. Виберіть усі правильні твердження (якщо такі є).\par
\vspace{2mm}

( 1 ) Виклик \code{this\_thread::sleep\_for(5s)} змушує потік проспати рівно 5 секунд.%bad

( 2 ) Чи запустить функція  async виконання в новому потоці можна задавати політикою.%good

( 3 ) Не можна блокувати заблокований м’ютекс mutex.%bad

( 4 ) Об’єкт promise можна скопіювати та роздати копії різним потокам, щоб кожен з них міг передати результат своєї роботи.%bad
%enditem


\vspace{5mm}
%beginitem
5. Виберіть усі правильні твердження (якщо такі є).\par
\vspace{2mm}

( 1 ) Структура даних є структурою даних без блокувань, якщо її реалізація не використовує м’ютекси.%bad

( 2 ) Мова С++20 має стандартну бібліотеку вільних від блокувань структур даних.%bad

( 3 ) За використання атомарних змінних структура даних є структурою даних без блокувань.%bad

( 4 ) Структура даних з блокуваннями може бути такою, що використовує атомарні змінні.%good
\newpage

%enditem


\vspace{5mm}
%beginitem
6. Виберіть усі правильні твердження (якщо такі є).\par
\vspace{2mm}

( 1 ) Потік може мати доступ до автоматичних змінних іншого потоку і це безпечно.%bad

( 2 ) Потік мови С++ не може отримати доступ до даних іншого потоку.%bad

( 3 ) Автоматичні змінні не можуть розділятися між потоками.%bad

( 4 ) Потік може мати доступ до автоматичних змінних іншого потоку, але це небезпечно.%good
%enditem


\vspace{5mm}
%beginitem
7. У цьому фрагменті коду
\begin{verbatim}
template <class T> void f(T a1, T a2){}

void g(){
    f(unique_ptr<int>(new int(3)), unique_ptr<int>(new int(35)));
\end{verbatim}

( 1 ) є стан гонки%bad

( 2 ) є гонка даних%bad

( 3 ) можливий виток ресурсів%bad

( 4 ) є невизначена поведінка%bad
%enditem
}
%end
\end {large}

\vspace{4mm}
%end
\newpage
%begin
%begin
\begin {large}

\textbf{22.11.2024 Контрольна робота, варіант  \No { 4 }}\par


{
%beginitem
1. По яких днях тижня відбувались лекції з предмету?

%enditem


\vspace{5mm}
%beginitem
2. Як зовуть (ім'я, по батькові) викладача, що вів практичні заняття?

%enditem


\vspace{5mm}
%beginitem
3. Нехай int var=13; Один потік збільшив змінну var на 3, а інший збільшив її в три рази.

гонка даних (data race), стан гонки (race condition)

Виберіть усі правильні твердження (якщо такі є).\par
\vspace{2mm}
( 1 ) Тут гарантовано нема стану гонки, але може бути гонка даних.%bad

( 2 ) Тут гарантовано є і стан гонки, і гонка даних.%bad

( 3 ) Тут може не бути ані гонки даних, ані стану гонки.%bad

( 4 ) Тут може не бути гонки даних.%good
%enditem


\vspace{5mm}
%beginitem
4. Виберіть усі правильні твердження (якщо такі є).\par
\vspace{2mm}

( 1 ) У стані гонки програма відчайдушно намагається захопити ресурси процесора.%bad

( 2 ) Клас thread має засоби, щоб повернути результат виконання запущеного потоку.%bad

( 3 ) Потік, запущений конструктором класу jthread, має обов’язково зупинитися, якщо пов’язаний об’єкт надішле команду зупинитися.%bad

( 4 ) За різних запусків програми на одних тих самих вхідних даних може виникати хронологічно різна послідовність обчислень.%good
%enditem


\vspace{5mm}
%beginitem
5. Виберіть усі правильні твердження (якщо такі є).\par
\vspace{2mm}

( 1 ) Мова С++20 має стандартну бібліотеку конкурентних структур даних.%bad

( 2 ) Кожна вільна від очікувань структура даних є структурою даних, вільною від блокувань.%good

( 3 ) Взаємоблокування потоків може виникнути тільки за використання м’ютексів.%bad

( 4 ) Стандартні контейнери stl безпечно використовувати з кількох потоків.%bad
\newpage

%enditem


\vspace{5mm}
%beginitem
6. Виберіть усі правильні твердження (якщо такі є).\par
\vspace{2mm}

( 1 ) Кожен потік програми може мати доступ до динамічних даних програми.%good

( 2 ) Потік мови С++ не може отримати доступ до автоматичних змінних іншого потоку.%bad

( 3 ) Динамічні змінні не можуть розділятися між потоками.%bad

( 4 ) Потік має доступ тільки до своїх власних даних та до глобальних статичних даних.%bad
%enditem


\vspace{5mm}
%beginitem
7. У цьому фрагменті коду
\begin{verbatim}
template <class T> void f(T a1, T a2){}

void g(){
    int i=5;
    f(unique_ptr<int>(new int(++i)), unique_ptr<int>(new int(++i)));
\end{verbatim}

( 1 ) є стан гонки%bad

( 2 ) є гонка даних%bad

( 3 ) є невизначена поведінка%good

( 4 ) можливий виток ресурсів%bad
%enditem
}
%end
\end {large}

\vspace{4mm}
%end
\newpage
%begin
%begin
\begin {large}

\textbf{22.11.2024 Контрольна робота, варіант  \No { 5 }}\par


{
%beginitem
1. По яких днях тижня відбувались лекції з предмету?

%enditem


\vspace{5mm}
%beginitem
2. Як зовуть (ім'я, по батькові) викладача, що вів практичні заняття?

%enditem


\vspace{5mm}
%beginitem
3. Нехай int var=13; Один потік збільшив змінну var на 3, а інший збільшив її в три рази.

гонка даних (data race), стан гонки (race condition)

Виберіть усі правильні твердження (якщо такі є).\par
\vspace{2mm}
( 1 ) Тут гарантовано нема ані гонки даних, ані стану гонки.%bad

( 2 ) Тут гарантовано є стан гонки, але гонки даних може не бути.%good

( 3 ) Тут гарантовано є гонка даних.%bad

( 4 ) Тут може не бути стану гонки.%bad
%enditem


\vspace{5mm}
%beginitem
4. Виберіть усі правильні твердження (якщо такі є).\par
\vspace{2mm}

( 1 ) Функція \code{sleep\_for} може змусити інший потік заснути.%bad

( 2 ) Об’єкт future можна скопіювати та роздати копії різним потокам, щоб вони змогли побачити результат.%bad

( 3 ) Щоб розблокувати заблокований \code{shared\_mutex}, його завжди достатньо розблокувати один раз.%good

( 4 ) Довільний потік може розблокувати заблокований м’ютекс.%bad
%enditem


\vspace{5mm}
%beginitem
5. Виберіть усі правильні твердження (якщо такі є).\par
\vspace{2mm}

( 1 ) Структура даних без блокувань гарантує, що жоден потік не буде очікувати доступу до виконання операцій над цією структурою даних.%bad

( 2 ) Операція над вільною від блокувань структурою даних може виконуватись скільки завгодно довго.%good

( 3 ) Для уникнення гонки даних структура даних може використовувати не більше одного м’ютекса.%bad

( 4 ) Структура даних без блокувань та вільна від блокувань структура даних --- це одне те саме.%bad
\newpage

%enditem


\vspace{5mm}
%beginitem
6. Виберіть усі правильні твердження (якщо такі є).\par
\vspace{2mm}

( 1 ) Потік має доступ тільки до своїх власних даних та до глобальних статичних даних.%bad

( 2 ) Потік мови С++ не може отримати доступ до автоматичних змінних іншого потоку.%bad

( 3 ) Потік може мати доступ до автоматичних змінних іншого потоку, але це небезпечно.%good

( 4 ) Динамічні змінні не можуть розділятися між потоками.%bad
%enditem


\vspace{5mm}
%beginitem
7. У цьому фрагменті коду
\begin{verbatim}
template <class T> void f(T a){}

void g(){
    int i=5;
    f(unique_ptr<int>(new int(i)))+(unique_ptr<int>(new int(++i)));
\end{verbatim}
( 1 ) є гонка даних%bad

( 2 ) можливий виток ресурсів%good

( 3 ) є невизначена поведінка%good

( 4 ) є стан гонки%bad
%enditem
}
%end
\end {large}

\vspace{4mm}
%end
\newpage
%begin
%begin
\begin {large}

\textbf{22.11.2024 Контрольна робота, варіант  \No { 6 }}\par


{
%beginitem
1. По яких днях тижня відбувались лекції з предмету?

%enditem


\vspace{5mm}
%beginitem
2. Як зовуть (ім'я, по батькові) викладача, що вів практичні заняття?

%enditem


\vspace{5mm}
%beginitem
3. Нехай int var=13; Один потік збільшив змінну var на 3, а інший збільшив її в три рази.

гонка даних (data race), стан гонки (race condition)

Виберіть усі правильні твердження (якщо такі є).\par
\vspace{2mm}
( 1 ) Тут може бути тільки стан гонки.%bad

( 2 ) Тут може бути тільки гонка даних.%bad

( 3 ) Тут гарантовано є стан гонки.%good

( 4 ) Тут гарантовано нема гонки даних, але може бути стан гонки.%bad
%enditem


\vspace{5mm}
%beginitem
4. Виберіть усі правильні твердження (якщо такі є).\par
\vspace{2mm}

( 1 ) Розблоковувати незаблокований м’ютекс не можна.%good

( 2 ) Досвічений програміст за однопотоковою програмою (текстом мовою С++) завжди може вказати послідовність виконуваних обчислень на конкретних вхідних даних.%bad

( 3 ) Тільки від компілятора залежить, чи запустить функція async виконання в новому потоці.%bad

( 4 ) Значення, що є результатом виконання функції, запущеної через async, зберігається безпосередньо в об’єкті проміза (promise).%bad
%enditem


\vspace{5mm}
%beginitem
5. Виберіть усі правильні твердження (якщо такі є).\par
\vspace{2mm}

( 1 ) Вільна від блокувань структура даних гарантує, що жоден потік не буде очікувати доступу до виконання операцій над цією структурою даних.%bad

( 2 ) Структура даних без блокувань гарантує завершення виконання довільної операції над цією структурою даних за обмежений час.%bad

( 3 ) Реалізація структури даних без блокувань не може використовувати м’ютекси.%good

( 4 ) Структура даних без блокувань зобов’язана використовувати атомарні змінні.%bad
\newpage

%enditem


\vspace{5mm}
%beginitem
6. Виберіть усі правильні твердження (якщо такі є).\par
\vspace{2mm}

( 1 ) Автоматичні змінні не можуть розділятися між потоками.%bad

( 2 ) Кожен потік програми може мати доступ до динамічних даних програми.%good

( 3 ) Потік мови С++ не може отримати доступ до даних іншого потоку.%bad

( 4 ) Потік може мати доступ до автоматичних змінних іншого потоку і це безпечно.%bad
%enditem


\vspace{5mm}
%beginitem
7. У цьому фрагменті коду
\begin{verbatim}
template <class T> void f(T a){}

void g(){
    f(unique_ptr<int>(new int(3)))+(unique_ptr<int>(new int(35)));
\end{verbatim}

( 1 ) є стан гонки%bad

( 2 ) є гонка даних%bad

( 3 ) є невизначена поведінка%bad

( 4 ) можливий виток ресурсів%good
%enditem
}
%end
\end {large}

\vspace{4mm}
%end
\newpage
%begin
%begin
\begin {large}

\textbf{22.11.2024 Контрольна робота, варіант  \No { 7 }}\par


{
%beginitem
1. По яких днях тижня відбувались лекції з предмету?

%enditem


\vspace{5mm}
%beginitem
2. Як зовуть (ім'я, по батькові) викладача, що вів практичні заняття?

%enditem


\vspace{5mm}
%beginitem
3. Нехай int var=13; Один потік збільшив змінну var на 3, а інший збільшив її в три рази.

гонка даних (data race), стан гонки (race condition)

Виберіть усі правильні твердження (якщо такі є).\par
\vspace{2mm}
( 1 ) Тут може не бути стану гонки.%bad

( 2 ) Тут може не бути ані гонки даних, ані стану гонки.%bad

( 3 ) Тут гарантовано є і стан гонки, і гонка даних.%bad

( 4 ) Тут гарантовано є стан гонки.%good
%enditem


\vspace{5mm}
%beginitem
4. Виберіть усі правильні твердження (якщо такі є).\par
\vspace{2mm}

( 1 ) М’ютекс має бути записаний на початку критичної секції, яку він забезпечує.%bad

( 2 ) Критичну секцію можна реалізувати за допомогою м’ютексу.%good

( 3 ) Тільки той потік, який обнулив значення семафора, може його збільшити.%bad

( 4 ) Коли програма спить, то вона не виконується.%bad
%enditem


\vspace{5mm}
%beginitem
5. Виберіть усі правильні твердження (якщо такі є).\par
\vspace{2mm}

( 1 ) Операція над вільною від блокувань структурою даних завжди виконається за обмежений час.%bad

( 2 ) Структура даних без блокувань гарантує завершення виконання довільної операції над цією структурою даних за обмежений час.%bad

( 3 ) Структура даних без блокувань може не бути вільною від блокувань.%good

( 4 ) Структура даних є структурою даних без блокувань, якщо її реалізація не використовує м’ютекси.%bad
\newpage

%enditem


\vspace{5mm}
%beginitem
6. Виберіть усі правильні твердження (якщо такі є).\par
\vspace{2mm}

( 1 ) Кожен потік програми може мати доступ до динамічних даних програми.%good

( 2 ) Потік може мати доступ до автоматичних змінних іншого потоку і це безпечно.%bad

( 3 ) Потік мови С++ не може отримати доступ до автоматичних змінних іншого потоку.%bad

( 4 ) Автоматичні змінні не можуть розділятися між потоками.%bad
%enditem


\vspace{5mm}
%beginitem
7. У цьому фрагменті коду
\begin{verbatim}
template <class T> void f(T a1, T a2){}

void g(){
    int i=5;
    f(unique_ptr<int>(new int(++i)), unique_ptr<int>(new int(++i)));
\end{verbatim}

( 1 ) є стан гонки%bad

( 2 ) є гонка даних%bad

( 3 ) є невизначена поведінка%good

( 4 ) можливий виток ресурсів%bad
%enditem
}
%end
\end {large}

\vspace{4mm}
%end
\newpage
%begin
%begin
\begin {large}

\textbf{22.11.2024 Контрольна робота, варіант  \No { 8 }}\par


{
%beginitem
1. По яких днях тижня відбувались лекції з предмету?

%enditem


\vspace{5mm}
%beginitem
2. Як зовуть (ім'я, по батькові) викладача, що вів практичні заняття?

%enditem


\vspace{5mm}
%beginitem
3. Нехай int var=13; Один потік збільшив змінну var на 3, а інший збільшив її в три рази.

гонка даних (data race), стан гонки (race condition)

Виберіть усі правильні твердження (якщо такі є).\par
\vspace{2mm}
( 1 ) Тут гарантовано є гонка даних.%bad

( 2 ) Тут гарантовано є стан гонки, але гонки даних може не бути.%good

( 3 ) Тут гарантовано нема стану гонки, але може бути гонка даних.%bad

( 4 ) Тут може бути тільки стан гонки.%bad
%enditem


\vspace{5mm}
%beginitem
4. Виберіть усі правильні твердження (якщо такі є).\par
\vspace{2mm}

( 1 ) Критична секція може складатися з кількох ділянок коду, які можуть розташовуватись де завгодно (не обов’язково зв’язний фрагмент коду).%good

( 2 ) У програмі може бути не більше однієї критичної секції.%bad

( 3 ) Критична секція це неперервна ділянка коду програми.%bad

( 4 ) Щоб розблокувати заблокований м’ютекс \code{recursive\_mutex}, його завжди достатньо розблокувати один раз.%bad
%enditem


\vspace{5mm}
%beginitem
5. Виберіть усі правильні твердження (якщо такі є).\par
\vspace{2mm}

( 1 ) Структура даних без блокувань гарантує, що жоден потік не буде очікувати доступу до виконання операцій над цією структурою даних.%bad

( 2 ) Структура даних без блокувань та вільна від блокувань структура даних --- це одне те саме.%bad

( 3 ) Структура даних може використовувати м’ютекси та атомарні змінні одночасно.%good

( 4 ) Стандартні контейнери stl безпечно використовувати з кількох потоків.%bad
\newpage

%enditem


\vspace{5mm}
%beginitem
6. Виберіть усі правильні твердження (якщо такі є).\par
\vspace{2mm}

( 1 ) Потік має доступ тільки до своїх власних даних та до глобальних статичних даних.%bad

( 2 ) Потік мови С++ не може отримати доступ до даних іншого потоку.%bad

( 3 ) Потік може мати доступ до автоматичних змінних іншого потоку, але це небезпечно.%good

( 4 ) Динамічні змінні не можуть розділятися між потоками.%bad
%enditem


\vspace{5mm}
%beginitem
7. У цьому фрагменті коду
\begin{verbatim}
template <class T> void f(T a){}

void g(){
    int i=5;
    f(unique_ptr<int>(new int(i)))+(unique_ptr<int>(new int(35)));
\end{verbatim}
( 1 ) можливий виток ресурсів%good

( 2 ) є невизначена поведінка%bad

( 3 ) є гонка даних%bad

( 4 ) є стан гонки%bad
%enditem
}
%end
\end {large}

\vspace{4mm}
%end
\newpage
%begin
%begin
\begin {large}

\textbf{22.11.2024 Контрольна робота, варіант  \No { 9 }}\par


{
%beginitem
1. По яких днях тижня відбувались лекції з предмету?

%enditem


\vspace{5mm}
%beginitem
2. Як зовуть (ім'я, по батькові) викладача, що вів практичні заняття?

%enditem


\vspace{5mm}
%beginitem
3. Нехай int var=13; Один потік збільшив змінну var на 3, а інший збільшив її в три рази.

гонка даних (data race), стан гонки (race condition)

Виберіть усі правильні твердження (якщо такі є).\par
\vspace{2mm}
( 1 ) Тут гарантовано нема ані гонки даних, ані стану гонки.%bad

( 2 ) Тут гарантовано нема гонки даних, але може бути стан гонки.%bad

( 3 ) Тут може бути тільки гонка даних.%bad

( 4 ) Тут може не бути гонки даних.%good
%enditem


\vspace{5mm}
%beginitem
4. Виберіть усі правильні твердження (якщо такі є).\par
\vspace{2mm}

( 1 ) Виклик функції \code{this\_thread::sleep\_for(5s)} змушує потік, з якого його здійснили, заснути якнайменш на 5 секунд.%good

( 2 ) Об’єкт promise можна скопіювати та роздати копії різним потокам, щоб кожен з них міг передати результат своєї роботи.%bad

( 3 ) Виклик \code{this\_thread::sleep\_for(5s)} змушує потік проспати рівно 5 секунд.%bad

( 4 ) Досвічений програміст за однопотоковою програмою (текстом мовою С++) завжди може вказати послідовність виконуваних обчислень на конкретних вхідних даних.%bad
%enditem


\vspace{5mm}
%beginitem
5. Виберіть усі правильні твердження (якщо такі є).\par
\vspace{2mm}

( 1 ) За використання атомарних змінних структура даних є структурою даних без блокувань.%bad

( 2 ) Якщо весь код структури даних утворює одну критичну секцію, то така структура даних може використовуватись в конкурентній програмі, але по-справжньо\-му конкурентною не є.%good

( 3 ) Мова С++20 має стандартну бібліотеку конкурентних структур даних.%bad

( 4 ) Одночасно використовувати атомарних змінних більше, ніж регістрів процесора, у програмі не можна.%bad
\newpage

%enditem


\vspace{5mm}
%beginitem
6. Виберіть усі правильні твердження (якщо такі є).\par
\vspace{2mm}

( 1 ) Потік може мати доступ до автоматичних змінних іншого потоку, але це небезпечно.%good

( 2 ) Потік має доступ тільки до своїх власних даних та до глобальних статичних даних.%bad

( 3 ) Потік мови С++ не може отримати доступ до автоматичних змінних іншого потоку.%bad

( 4 ) Потік мови С++ не може отримати доступ до даних іншого потоку.%bad
%enditem


\vspace{5mm}
%beginitem
7. У цьому фрагменті коду
\begin{verbatim}
template <class T> void f(T a1, T a2){}

void g(){
    int i=5;
    f(unique_ptr<int>(new int(++i)), unique_ptr<int>(new int(++i)));
\end{verbatim}

( 1 ) є стан гонки%bad

( 2 ) є невизначена поведінка%good

( 3 ) можливий виток ресурсів%bad

( 4 ) є гонка даних%bad
%enditem
}
%end
\end {large}

\vspace{4mm}
%end
\newpage
%begin
%begin
\begin {large}

\textbf{22.11.2024 Контрольна робота, варіант  \No { 10 }}\par


{
%beginitem
1. По яких днях тижня відбувались лекції з предмету?

%enditem


\vspace{5mm}
%beginitem
2. Як зовуть (ім'я, по батькові) викладача, що вів практичні заняття?

%enditem


\vspace{5mm}
%beginitem
3. Нехай int var=13; Один потік збільшив змінну var на 3, а інший збільшив її в три рази.

гонка даних (data race), стан гонки (race condition)

Виберіть усі правильні твердження (якщо такі є).\par
\vspace{2mm}
( 1 ) Тут гарантовано є стан гонки, але гонки даних може не бути.%good

( 2 ) Тут може бути тільки гонка даних.%bad

( 3 ) Тут гарантовано нема стану гонки, але може бути гонка даних.%bad

( 4 ) Тут гарантовано є гонка даних.%bad
%enditem


\vspace{5mm}
%beginitem
4. Виберіть усі правильні твердження (якщо такі є).\par
\vspace{2mm}

( 1 ) Тільки від компілятора залежить, чи запустить функція async виконання в новому потоці.%bad

( 2 ) Потік в активному очікуванні споживає ресурс процесора.%good

( 3 ) Функція \code{sleep\_for} може змусити інший потік заснути.%bad

( 4 ) Функція async запускає виконання в новому потоці.%bad
%enditem


\vspace{5mm}
%beginitem
5. Виберіть усі правильні твердження (якщо такі є).\par
\vspace{2mm}

( 1 ) Взаємоблокування потоків може виникнути тільки за використання м’ютексів.%bad

( 2 ) Структура даних без блокувань зобов’язана використовувати атомарні змінні.%bad

( 3 ) Кожна вільна від очікувань структура даних є структурою даних, вільною від блокувань.%good

( 4 ) Для уникнення гонки даних структура даних може використовувати не більше одного м’ютекса.%bad
\newpage

%enditem


\vspace{5mm}
%beginitem
6. Виберіть усі правильні твердження (якщо такі є).\par
\vspace{2mm}

( 1 ) Потік може мати доступ до автоматичних змінних іншого потоку і це безпечно.%bad

( 2 ) Автоматичні змінні не можуть розділятися між потоками.%bad

( 3 ) Кожен потік програми може мати доступ до динамічних даних програми.%good

( 4 ) Динамічні змінні не можуть розділятися між потоками.%bad
%enditem


\vspace{5mm}
%beginitem
7. У цьому фрагменті коду
\begin{verbatim}
template <class T> void f(T a){}

void g(){
    f(unique_ptr<int>(new int(3)))+(unique_ptr<int>(new int(35)));
\end{verbatim}

( 1 ) є стан гонки%bad

( 2 ) є невизначена поведінка%bad

( 3 ) є гонка даних%bad

( 4 ) можливий виток ресурсів%good
%enditem
}
%end
\end {large}

\vspace{4mm}
%end
\newpage
%begin
%begin
\begin {large}

\textbf{22.11.2024 Контрольна робота, варіант  \No { 11 }}\par


{
%beginitem
1. По яких днях тижня відбувались лекції з предмету?

%enditem


\vspace{5mm}
%beginitem
2. Як зовуть (ім'я, по батькові) викладача, що вів практичні заняття?

%enditem


\vspace{5mm}
%beginitem
3. Нехай int var=13; Один потік збільшив змінну var на 3, а інший збільшив її в три рази.

гонка даних (data race), стан гонки (race condition)

Виберіть усі правильні твердження (якщо такі є).\par
\vspace{2mm}
( 1 ) Тут гарантовано нема гонки даних, але може бути стан гонки.%bad

( 2 ) Тут може не бути ані гонки даних, ані стану гонки.%bad

( 3 ) Тут може бути тільки стан гонки.%bad

( 4 ) Тут гарантовано є стан гонки.%good
%enditem


\vspace{5mm}
%beginitem
4. Виберіть усі правильні твердження (якщо такі є).\par
\vspace{2mm}

( 1 ) У стані гонки програма відчайдушно намагається захопити ресурси процесора.%bad

( 2 ) Якщо два обчислення програми на С++ тільки читають одну ту саму ділянку пам'яті, то на них гонка даних  виникнути не може.%good

( 3 ) Критична секція це неперервна ділянка коду програми.%bad

( 4 ) Щоб розблокувати заблокований м’ютекс \code{recursive\_mutex}, його завжди достатньо розблокувати один раз.%bad
%enditem


\vspace{5mm}
%beginitem
5. Виберіть усі правильні твердження (якщо такі є).\par
\vspace{2mm}

( 1 ) Вільних від очікувань структур даних не існує.%bad

( 2 ) Реалізація структури даних з блокуваннями може використовувати атомарні змінні.%good

( 3 ) Вільна від блокувань структура даних гарантує, що жоден потік не буде очікувати доступу до виконання операцій над цією структурою даних.%bad

( 4 ) Реалізації різних операцій над структурою даних повинні використовувати різні м’ютекси.%bad
\newpage

%enditem


\vspace{5mm}
%beginitem
6. Виберіть усі правильні твердження (якщо такі є).\par
\vspace{2mm}

( 1 ) Динамічні змінні не можуть розділятися між потоками.%bad

( 2 ) Кожен потік програми може мати доступ до динамічних даних програми.%good

( 3 ) Потік мови С++ не може отримати доступ до даних іншого потоку.%bad

( 4 ) Потік має доступ тільки до своїх власних даних та до глобальних статичних даних.%bad
%enditem


\vspace{5mm}
%beginitem
7. У цьому фрагменті коду
\begin{verbatim}
template <class T> void f(T a1, T a2){}

void g(){
    int i=5;
    f(unique_ptr<int>(new int(++i)), unique_ptr<int>(new int(++i)));
\end{verbatim}

( 1 ) можливий виток ресурсів%bad

( 2 ) є невизначена поведінка%good

( 3 ) є стан гонки%bad

( 4 ) є гонка даних%bad
%enditem
}
%end
\end {large}

\vspace{4mm}
%end
\newpage
%begin
%begin
\begin {large}

\textbf{22.11.2024 Контрольна робота, варіант  \No { 12 }}\par


{
%beginitem
1. По яких днях тижня відбувались лекції з предмету?

%enditem


\vspace{5mm}
%beginitem
2. Як зовуть (ім'я, по батькові) викладача, що вів практичні заняття?

%enditem


\vspace{5mm}
%beginitem
3. Нехай int var=13; Один потік збільшив змінну var на 3, а інший збільшив її в три рази.

гонка даних (data race), стан гонки (race condition)

Виберіть усі правильні твердження (якщо такі є).\par
\vspace{2mm}
( 1 ) Тут може не бути гонки даних.%good

( 2 ) Тут гарантовано нема ані гонки даних, ані стану гонки.%bad

( 3 ) Тут гарантовано є і стан гонки, і гонка даних.%bad

( 4 ) Тут може не бути стану гонки.%bad
%enditem


\vspace{5mm}
%beginitem
4. Виберіть усі правильні твердження (якщо такі є).\par
\vspace{2mm}

( 1 ) Не можна блокувати заблокований м’ютекс mutex.%bad

( 2 ) М’ютекс має бути записаний на початку критичної секції, яку він забезпечує.%bad

( 3 ) Коли потік спить, він не споживає ресурс процесора.%good

( 4 ) Клас thread має засоби, щоб повернути результат виконання запущеного потоку.%bad
%enditem


\vspace{5mm}
%beginitem
5. Виберіть усі правильні твердження (якщо такі є).\par
\vspace{2mm}

( 1 ) Операція над вільною від блокувань структурою даних може виконуватись скільки завгодно довго.%good

( 2 ) Мова С++20 має стандартну бібліотеку вільних від блокувань структур даних.%bad

( 3 ) На операціях структури даних без блокувань взаємоблокувань потоків виникнути не може.%bad

( 4 ) Операція над структурою даних без блокувань завжди виконається за обмежений час.%bad
\newpage

%enditem


\vspace{5mm}
%beginitem
6. Виберіть усі правильні твердження (якщо такі є).\par
\vspace{2mm}

( 1 ) Потік мови С++ не може отримати доступ до автоматичних змінних іншого потоку.%bad

( 2 ) Автоматичні змінні не можуть розділятися між потоками.%bad

( 3 ) Потік може мати доступ до автоматичних змінних іншого потоку, але це небезпечно.%good

( 4 ) Потік може мати доступ до автоматичних змінних іншого потоку і це безпечно.%bad
%enditem


\vspace{5mm}
%beginitem
7. У цьому фрагменті коду
\begin{verbatim}
template <class T> void f(T a1, T a2){}

void g(){
    f(unique_ptr<int>(new int(3)), unique_ptr<int>(new int(35)));
\end{verbatim}

( 1 ) є гонка даних%bad

( 2 ) є стан гонки%bad

( 3 ) можливий виток ресурсів%bad

( 4 ) є невизначена поведінка%bad
%enditem
}
%end
\end {large}

\vspace{4mm}
%end
\newpage
%begin
%begin
\begin {large}

\textbf{22.11.2024 Контрольна робота, варіант  \No { 13 }}\par


{
%beginitem
1. По яких днях тижня відбувались лекції з предмету?

%enditem


\vspace{5mm}
%beginitem
2. Як зовуть (ім'я, по батькові) викладача, що вів практичні заняття?

%enditem


\vspace{5mm}
%beginitem
3. Нехай int var=13; Один потік збільшив змінну var на 3, а інший збільшив її в три рази.

гонка даних (data race), стан гонки (race condition)

Виберіть усі правильні твердження (якщо такі є).\par
\vspace{2mm}
( 1 ) Тут може не бути ані гонки даних, ані стану гонки.%bad

( 2 ) Тут може бути тільки стан гонки.%bad

( 3 ) Тут гарантовано є стан гонки, але гонки даних може не бути.%good

( 4 ) Тут гарантовано нема гонки даних, але може бути стан гонки.%bad
%enditem


\vspace{5mm}
%beginitem
4. Виберіть усі правильні твердження (якщо такі є).\par
\vspace{2mm}

( 1 ) Об’єкти класів promise та future можна використовувати тільки в багатопотокових програмах.%bad

( 2 ) Текст однопотокової програми мовою С++ однозначно визначає послідовність виконуваних обчислень.%bad

( 3 ) У синтаксично коректній програмі мовою С++ стан гонки (race condition) може виникати.%good

( 4 ) У програмі може бути не більше однієї критичної секції.%bad
%enditem


\vspace{5mm}
%beginitem
5. Виберіть усі правильні твердження (якщо такі є).\par
\vspace{2mm}

( 1 ) Мова С++20 має стандартну бібліотеку конкурентних структур даних.%bad

( 2 ) Взаємоблокування потоків може виникнути тільки за використання м’ютексів.%bad

( 3 ) Операція над структурою даних без блокувань завжди виконається за обмежений час.%bad

( 4 ) Структура даних з блокуваннями може бути такою, що використовує атомарні змінні.%good
\newpage

%enditem


\vspace{5mm}
%beginitem
6. Виберіть усі правильні твердження (якщо такі є).\par
\vspace{2mm}

( 1 ) Потік може мати доступ до автоматичних змінних іншого потоку і це безпечно.%bad

( 2 ) Автоматичні змінні не можуть розділятися між потоками.%bad

( 3 ) Потік має доступ тільки до своїх власних даних та до глобальних статичних даних.%bad

( 4 ) Кожен потік програми може мати доступ до динамічних даних програми.%good
%enditem


\vspace{5mm}
%beginitem
7. У цьому фрагменті коду
\begin{verbatim}
template <class T> void f(T a1, T a2){}

void g(){
    int i=5;
    f(unique_ptr<int>(new int(i++)), unique_ptr<int>(new int(++i)));
\end{verbatim}

( 1 ) є стан гонки%bad

( 2 ) є невизначена поведінка%good

( 3 ) можливий виток ресурсів%bad

( 4 ) є гонка даних%bad
%enditem
}
%end
\end {large}

\vspace{4mm}
%end
\newpage
%begin
%begin
\begin {large}

\textbf{22.11.2024 Контрольна робота, варіант  \No { 14 }}\par


{
%beginitem
1. По яких днях тижня відбувались лекції з предмету?

%enditem


\vspace{5mm}
%beginitem
2. Як зовуть (ім'я, по батькові) викладача, що вів практичні заняття?

%enditem


\vspace{5mm}
%beginitem
3. Нехай int var=13; Один потік збільшив змінну var на 3, а інший збільшив її в три рази.

гонка даних (data race), стан гонки (race condition)

Виберіть усі правильні твердження (якщо такі є).\par
\vspace{2mm}
( 1 ) Тут гарантовано нема стану гонки, але може бути гонка даних.%bad

( 2 ) Тут гарантовано нема ані гонки даних, ані стану гонки.%bad

( 3 ) Тут може не бути стану гонки.%bad

( 4 ) Тут гарантовано є стан гонки.%good
%enditem


\vspace{5mm}
%beginitem
4. Виберіть усі правильні твердження (якщо такі є).\par
\vspace{2mm}

( 1 ) Порядок обчислень у потоці для одних тих самих вхідних даних може залежати від використовуваного компілятора.%good

( 2 ) Розблоковувати незаблокований м’ютекс можна.%bad

( 3 ) Значення, що є результатом виконання функції, запущеної через async, зберігається безпосередньо в об’єкті ф’ючерса (future).%bad

( 4 ) Коли програма спить, то вона не виконується.%bad
%enditem


\vspace{5mm}
%beginitem
5. Виберіть усі правильні твердження (якщо такі є).\par
\vspace{2mm}

( 1 ) Одночасно використовувати атомарних змінних більше, ніж регістрів процесора, у програмі не можна.%bad

( 2 ) Структура даних без блокувань гарантує, що жоден потік не буде очікувати доступу до виконання операцій над цією структурою даних.%bad

( 3 ) Стандартні контейнери stl безпечно використовувати з кількох потоків.%bad

( 4 ) Реалізація структури даних без блокувань не може використовувати м’ютекси.%good
\newpage

%enditem


\vspace{5mm}
%beginitem
6. Виберіть усі правильні твердження (якщо такі є).\par
\vspace{2mm}

( 1 ) Динамічні змінні не можуть розділятися між потоками.%bad

( 2 ) Потік може мати доступ до автоматичних змінних іншого потоку, але це небезпечно.%good

( 3 ) Потік мови С++ не може отримати доступ до автоматичних змінних іншого потоку.%bad

( 4 ) Потік мови С++ не може отримати доступ до даних іншого потоку.%bad
%enditem


\vspace{5mm}
%beginitem
7. У цьому фрагменті коду
\begin{verbatim}
template <class T> void f(T a){}

void g(){
    int i=5;
    f(unique_ptr<int>(new int(i)))+(unique_ptr<int>(new int(++i)));
\end{verbatim}
( 1 ) є невизначена поведінка%good

( 2 ) є стан гонки%bad

( 3 ) є гонка даних%bad

( 4 ) можливий виток ресурсів%good
%enditem
}
%end
\end {large}

\vspace{4mm}
%end
\newpage
%begin
%begin
\begin {large}

\textbf{22.11.2024 Контрольна робота, варіант  \No { 15 }}\par


{
%beginitem
1. По яких днях тижня відбувались лекції з предмету?

%enditem


\vspace{5mm}
%beginitem
2. Як зовуть (ім'я, по батькові) викладача, що вів практичні заняття?

%enditem


\vspace{5mm}
%beginitem
3. Нехай int var=13; Один потік збільшив змінну var на 3, а інший збільшив її в три рази.

гонка даних (data race), стан гонки (race condition)

Виберіть усі правильні твердження (якщо такі є).\par
\vspace{2mm}
( 1 ) Тут може не бути гонки даних.%good

( 2 ) Тут гарантовано є гонка даних.%bad

( 3 ) Тут гарантовано є і стан гонки, і гонка даних.%bad

( 4 ) Тут може бути тільки гонка даних.%bad
%enditem


\vspace{5mm}
%beginitem
4. Виберіть усі правильні твердження (якщо такі є).\par
\vspace{2mm}

( 1 ) Блокувати заблокований м’ютекс  \code{mutex} іноді можна.%good

( 2 ) Об’єкт future можна скопіювати та роздати копії різним потокам, щоб вони змогли побачити результат.%bad

( 3 ) Потік, запущений конструктором класу jthread, має обов’язково зупинитися, якщо пов’язаний об’єкт надішле команду зупинитися.%bad

( 4 ) Довільний потік може розблокувати заблокований м’ютекс.%bad
%enditem


\vspace{5mm}
%beginitem
5. Виберіть усі правильні твердження (якщо такі є).\par
\vspace{2mm}

( 1 ) Структура даних без блокувань може не бути вільною від блокувань.%good

( 2 ) Вільних від очікувань структур даних не існує.%bad

( 3 ) Структура даних без блокувань зобов’язана використовувати атомарні змінні.%bad

( 4 ) Мова С++20 має стандартну бібліотеку вільних від блокувань структур даних.%bad
\newpage

%enditem


\vspace{5mm}
%beginitem
6. Виберіть усі правильні твердження (якщо такі є).\par
\vspace{2mm}

( 1 ) Динамічні змінні не можуть розділятися між потоками.%bad

( 2 ) Потік мови С++ не може отримати доступ до даних іншого потоку.%bad

( 3 ) Кожен потік програми може мати доступ до динамічних даних програми.%good

( 4 ) Потік мови С++ не може отримати доступ до автоматичних змінних іншого потоку.%bad
%enditem


\vspace{5mm}
%beginitem
7. У цьому фрагменті коду
\begin{verbatim}
template <class T> void f(T a1, T a2){}

void g(){
    int i=5;
    f(unique_ptr<int>(new int(++i)), unique_ptr<int>(new int(++i)));
\end{verbatim}

( 1 ) є невизначена поведінка%good

( 2 ) є стан гонки%bad

( 3 ) можливий виток ресурсів%bad

( 4 ) є гонка даних%bad
%enditem
}
%end
\end {large}

\vspace{4mm}
%end
\newpage
%begin
%begin
\begin {large}

\textbf{22.11.2024 Контрольна робота, варіант  \No { 16 }}\par


{
%beginitem
1. По яких днях тижня відбувались лекції з предмету?

%enditem


\vspace{5mm}
%beginitem
2. Як зовуть (ім'я, по батькові) викладача, що вів практичні заняття?

%enditem


\vspace{5mm}
%beginitem
3. Нехай int var=13; Один потік збільшив змінну var на 3, а інший збільшив її в три рази.

гонка даних (data race), стан гонки (race condition)

Виберіть усі правильні твердження (якщо такі є).\par
\vspace{2mm}
( 1 ) Тут може не бути ані гонки даних, ані стану гонки.%bad

( 2 ) Тут гарантовано нема стану гонки, але може бути гонка даних.%bad

( 3 ) Тут гарантовано є стан гонки.%good

( 4 ) Тут гарантовано нема гонки даних, але може бути стан гонки.%bad
%enditem


\vspace{5mm}
%beginitem
4. Виберіть усі правильні твердження (якщо такі є).\par
\vspace{2mm}

( 1 ) Значення, що є результатом виконання функції, запущеної через async, зберігається безпосередньо в об’єкті проміза (promise).%bad

( 2 ) Секція є критичною, якщо під час її виконання не можна генерувати винятки.%bad

( 3 ) Клас thread не має засобів, щоб повернути результат виконання запущеного потоку.%good

( 4 ) Тільки той потік, який обнулив значення семафора, може його збільшити.%bad
%enditem


\vspace{5mm}
%beginitem
5. Виберіть усі правильні твердження (якщо такі є).\par
\vspace{2mm}

( 1 ) Структура даних без блокувань та вільна від блокувань структура даних --- це одне те саме.%bad

( 2 ) Для уникнення гонки даних структура даних може використовувати не більше одного м’ютекса.%bad

( 3 ) На операціях структури даних без блокувань взаємоблокувань потоків виникнути не може.%bad

( 4 ) За використання структурою даних рівно одного м’ютексу, взаємоблокування на її операціях неможливі.%good
\newpage

%enditem


\vspace{5mm}
%beginitem
6. Виберіть усі правильні твердження (якщо такі є).\par
\vspace{2mm}

( 1 ) Потік має доступ тільки до своїх власних даних та до глобальних статичних даних.%bad

( 2 ) Потік може мати доступ до автоматичних змінних іншого потоку і це безпечно.%bad

( 3 ) Автоматичні змінні не можуть розділятися між потоками.%bad

( 4 ) Потік може мати доступ до автоматичних змінних іншого потоку, але це небезпечно.%good
%enditem


\vspace{5mm}
%beginitem
7. У цьому фрагменті коду
\begin{verbatim}
template <class T> void f(T a){}

void g(){
    int i=5;
    f(unique_ptr<int>(new int(i)))+(unique_ptr<int>(new int(++i)));
\end{verbatim}
( 1 ) є стан гонки%bad

( 2 ) є гонка даних%bad

( 3 ) є невизначена поведінка%good

( 4 ) можливий виток ресурсів%good
%enditem
}
%end
\end {large}

\vspace{4mm}
%end
\newpage
%begin
%begin
\begin {large}

\textbf{22.11.2024 Контрольна робота, варіант  \No { 17 }}\par


{
%beginitem
1. По яких днях тижня відбувались лекції з предмету?

%enditem


\vspace{5mm}
%beginitem
2. Як зовуть (ім'я, по батькові) викладача, що вів практичні заняття?

%enditem


\vspace{5mm}
%beginitem
3. Нехай int var=13; Один потік збільшив змінну var на 3, а інший збільшив її в три рази.

гонка даних (data race), стан гонки (race condition)

Виберіть усі правильні твердження (якщо такі є).\par
\vspace{2mm}
( 1 ) Тут гарантовано є і стан гонки, і гонка даних.%bad

( 2 ) Тут гарантовано є стан гонки, але гонки даних може не бути.%good

( 3 ) Тут гарантовано є гонка даних.%bad

( 4 ) Тут гарантовано нема ані гонки даних, ані стану гонки.%bad
%enditem


\vspace{5mm}
%beginitem
4. Виберіть усі правильні твердження (якщо такі є).\par
\vspace{2mm}

( 1 ) Тільки той потік, який обнулив значення семафора, може його збільшити.%bad

( 2 ) У стані гонки програма відчайдушно намагається захопити ресурси процесора.%bad

( 3 ) Чи запустить функція  async виконання в новому потоці можна задавати політикою.%good

( 4 ) У програмі може бути не більше однієї критичної секції.%bad
%enditem


\vspace{5mm}
%beginitem
5. Виберіть усі правильні твердження (якщо такі є).\par
\vspace{2mm}

( 1 ) Структура даних без блокувань гарантує завершення виконання довільної операції над цією структурою даних за обмежений час.%bad

( 2 ) Операція над вільною від блокувань структурою даних завжди виконається за обмежений час.%bad

( 3 ) Структура даних може використовувати м’ютекси та атомарні змінні одночасно.%good

( 4 ) Реалізації різних операцій над структурою даних повинні використовувати різні м’ютекси.%bad
\newpage

%enditem


\vspace{5mm}
%beginitem
6. Виберіть усі правильні твердження (якщо такі є).\par
\vspace{2mm}

( 1 ) Потік мови С++ не може отримати доступ до автоматичних змінних іншого потоку.%bad

( 2 ) Динамічні змінні не можуть розділятися між потоками.%bad

( 3 ) Потік може мати доступ до автоматичних змінних іншого потоку, але це небезпечно.%good

( 4 ) Потік може мати доступ до автоматичних змінних іншого потоку і це безпечно.%bad
%enditem


\vspace{5mm}
%beginitem
7. У цьому фрагменті коду
\begin{verbatim}
template <class T> void f(T a){}

void g(){
    int i=5;
    f(unique_ptr<int>(new int(i)))+(unique_ptr<int>(new int(35)));
\end{verbatim}
( 1 ) можливий виток ресурсів%good

( 2 ) є стан гонки%bad

( 3 ) є невизначена поведінка%bad

( 4 ) є гонка даних%bad
%enditem
}
%end
\end {large}

\vspace{4mm}
%end
\newpage
%begin
%begin
\begin {large}

\textbf{22.11.2024 Контрольна робота, варіант  \No { 18 }}\par


{
%beginitem
1. По яких днях тижня відбувались лекції з предмету?

%enditem


\vspace{5mm}
%beginitem
2. Як зовуть (ім'я, по батькові) викладача, що вів практичні заняття?

%enditem


\vspace{5mm}
%beginitem
3. Нехай int var=13; Один потік збільшив змінну var на 3, а інший збільшив її в три рази.

гонка даних (data race), стан гонки (race condition)

Виберіть усі правильні твердження (якщо такі є).\par
\vspace{2mm}
( 1 ) Тут може не бути гонки даних.%good

( 2 ) Тут може бути тільки гонка даних.%bad

( 3 ) Тут може бути тільки стан гонки.%bad

( 4 ) Тут може не бути стану гонки.%bad
%enditem


\vspace{5mm}
%beginitem
4. Виберіть усі правильні твердження (якщо такі є).\par
\vspace{2mm}

( 1 ) М’ютекс має бути записаний на початку критичної секції, яку він забезпечує.%bad

( 2 ) Клас thread має засоби, щоб повернути результат виконання запущеного потоку.%bad

( 3 ) Розблокувати м’ютекс може тільки той потік, що його заблокував.%good

( 4 ) Секція є критичною, якщо під час її виконання не можна генерувати винятки.%bad
%enditem


\vspace{5mm}
%beginitem
5. Виберіть усі правильні твердження (якщо такі є).\par
\vspace{2mm}

( 1 ) Якщо весь код структури даних утворює одну критичну секцію, то така структура даних може використовуватись в конкурентній програмі, але по-справжньо\-му конкурентною не є.%good

( 2 ) Структура даних є структурою даних без блокувань, якщо її реалізація не використовує м’ютекси.%bad

( 3 ) За використання атомарних змінних структура даних є структурою даних без блокувань.%bad

( 4 ) Вільна від блокувань структура даних гарантує, що жоден потік не буде очікувати доступу до виконання операцій над цією структурою даних.%bad
\newpage

%enditem


\vspace{5mm}
%beginitem
6. Виберіть усі правильні твердження (якщо такі є).\par
\vspace{2mm}

( 1 ) Автоматичні змінні не можуть розділятися між потоками.%bad

( 2 ) Потік мови С++ не може отримати доступ до даних іншого потоку.%bad

( 3 ) Кожен потік програми може мати доступ до динамічних даних програми.%good

( 4 ) Потік має доступ тільки до своїх власних даних та до глобальних статичних даних.%bad
%enditem


\vspace{5mm}
%beginitem
7. У цьому фрагменті коду
\begin{verbatim}
template <class T> void f(T a){}

void g(){
    f(unique_ptr<int>(new int(3)))+(unique_ptr<int>(new int(35)));
\end{verbatim}

( 1 ) є гонка даних%bad

( 2 ) є невизначена поведінка%bad

( 3 ) є стан гонки%bad

( 4 ) можливий виток ресурсів%good
%enditem
}
%end
\end {large}

\vspace{4mm}
%end
\newpage
%begin
%begin
\begin {large}

\textbf{22.11.2024 Контрольна робота, варіант  \No { 19 }}\par


{
%beginitem
1. По яких днях тижня відбувались лекції з предмету?

%enditem


\vspace{5mm}
%beginitem
2. Як зовуть (ім'я, по батькові) викладача, що вів практичні заняття?

%enditem


\vspace{5mm}
%beginitem
3. Нехай int var=13; Один потік збільшив змінну var на 3, а інший збільшив її в три рази.

гонка даних (data race), стан гонки (race condition)

Виберіть усі правильні твердження (якщо такі є).\par
\vspace{2mm}
( 1 ) Тут гарантовано є і стан гонки, і гонка даних.%bad

( 2 ) Тут гарантовано є стан гонки, але гонки даних може не бути.%good

( 3 ) Тут може бути тільки гонка даних.%bad

( 4 ) Тут гарантовано нема ані гонки даних, ані стану гонки.%bad
%enditem


\vspace{5mm}
%beginitem
4. Виберіть усі правильні твердження (якщо такі є).\par
\vspace{2mm}

( 1 ) Функція async запускає виконання в новому потоці.%bad

( 2 ) Об’єкт promise можна скопіювати та роздати копії різним потокам, щоб кожен з них міг передати результат своєї роботи.%bad

( 3 ) Текст однопотокової програми мовою С++ однозначно визначає послідовність виконуваних обчислень.%bad

( 4 ) Блокувати заблокований м’ютекс  \code{mutex} іноді можна.%good
%enditem


\vspace{5mm}
%beginitem
5. Виберіть усі правильні твердження (якщо такі є).\par
\vspace{2mm}

( 1 ) Структура даних може використовувати м’ютекси та атомарні змінні одночасно.%good

( 2 ) Одночасно використовувати атомарних змінних більше, ніж регістрів процесора, у програмі не можна.%bad

( 3 ) Стандартні контейнери stl безпечно використовувати з кількох потоків.%bad

( 4 ) Мова С++20 має стандартну бібліотеку конкурентних структур даних.%bad
\newpage

%enditem


\vspace{5mm}
%beginitem
6. Виберіть усі правильні твердження (якщо такі є).\par
\vspace{2mm}

( 1 ) Потік мови С++ не може отримати доступ до автоматичних змінних іншого потоку.%bad

( 2 ) Потік мови С++ не може отримати доступ до даних іншого потоку.%bad

( 3 ) Потік може мати доступ до автоматичних змінних іншого потоку і це безпечно.%bad

( 4 ) Кожен потік програми може мати доступ до динамічних даних програми.%good
%enditem


\vspace{5mm}
%beginitem
7. У цьому фрагменті коду
\begin{verbatim}
template <class T> void f(T a1, T a2){}

void g(){
    f(unique_ptr<int>(new int(3)), unique_ptr<int>(new int(35)));
\end{verbatim}

( 1 ) можливий виток ресурсів%bad

( 2 ) є стан гонки%bad

( 3 ) є гонка даних%bad

( 4 ) є невизначена поведінка%bad
%enditem
}
%end
\end {large}

\vspace{4mm}
%end
\newpage
%begin
%begin
\begin {large}

\textbf{22.11.2024 Контрольна робота, варіант  \No { 20 }}\par


{
%beginitem
1. По яких днях тижня відбувались лекції з предмету?

%enditem


\vspace{5mm}
%beginitem
2. Як зовуть (ім'я, по батькові) викладача, що вів практичні заняття?

%enditem


\vspace{5mm}
%beginitem
3. Нехай int var=13; Один потік збільшив змінну var на 3, а інший збільшив її в три рази.

гонка даних (data race), стан гонки (race condition)

Виберіть усі правильні твердження (якщо такі є).\par
\vspace{2mm}
( 1 ) Тут може не бути стану гонки.%bad

( 2 ) Тут гарантовано нема гонки даних, але може бути стан гонки.%bad

( 3 ) Тут може не бути ані гонки даних, ані стану гонки.%bad

( 4 ) Тут може не бути гонки даних.%good
%enditem


\vspace{5mm}
%beginitem
4. Виберіть усі правильні твердження (якщо такі є).\par
\vspace{2mm}

( 1 ) Коли програма спить, то вона не виконується.%bad

( 2 ) Об’єкт future можна скопіювати та роздати копії різним потокам, щоб вони змогли побачити результат.%bad

( 3 ) Не можна блокувати заблокований м’ютекс mutex.%bad

( 4 ) Критична секція може складатися з кількох ділянок коду, які можуть розташовуватись де завгодно (не обов’язково зв’язний фрагмент коду).%good
%enditem


\vspace{5mm}
%beginitem
5. Виберіть усі правильні твердження (якщо такі є).\par
\vspace{2mm}

( 1 ) Структура даних без блокувань та вільна від блокувань структура даних --- це одне те саме.%bad

( 2 ) Операція над вільною від блокувань структурою даних завжди виконається за обмежений час.%bad

( 3 ) Взаємоблокування потоків може виникнути тільки за використання м’ютексів.%bad

( 4 ) Якщо весь код структури даних утворює одну критичну секцію, то така структура даних може використовуватись в конкурентній програмі, але по-справжньо\-му конкурентною не є.%good
\newpage

%enditem


\vspace{5mm}
%beginitem
6. Виберіть усі правильні твердження (якщо такі є).\par
\vspace{2mm}

( 1 ) Потік має доступ тільки до своїх власних даних та до глобальних статичних даних.%bad

( 2 ) Автоматичні змінні не можуть розділятися між потоками.%bad

( 3 ) Потік може мати доступ до автоматичних змінних іншого потоку, але це небезпечно.%good

( 4 ) Динамічні змінні не можуть розділятися між потоками.%bad
%enditem


\vspace{5mm}
%beginitem
7. У цьому фрагменті коду
\begin{verbatim}
template <class T> void f(T a1, T a2){}

void g(){
    int i=5;
    f(unique_ptr<int>(new int(i++)), unique_ptr<int>(new int(++i)));
\end{verbatim}

( 1 ) є стан гонки%bad

( 2 ) є гонка даних%bad

( 3 ) є невизначена поведінка%good

( 4 ) можливий виток ресурсів%bad
%enditem
}
%end
\end {large}

\vspace{4mm}
%end
\newpage
%begin
%begin
\begin {large}

\textbf{22.11.2024 Контрольна робота, варіант  \No { 21 }}\par


{
%beginitem
1. По яких днях тижня відбувались лекції з предмету?

%enditem


\vspace{5mm}
%beginitem
2. Як зовуть (ім'я, по батькові) викладача, що вів практичні заняття?

%enditem


\vspace{5mm}
%beginitem
3. Нехай int var=13; Один потік збільшив змінну var на 3, а інший збільшив її в три рази.

гонка даних (data race), стан гонки (race condition)

Виберіть усі правильні твердження (якщо такі є).\par
\vspace{2mm}
( 1 ) Тут гарантовано нема стану гонки, але може бути гонка даних.%bad

( 2 ) Тут гарантовано є стан гонки.%good

( 3 ) Тут гарантовано є гонка даних.%bad

( 4 ) Тут може бути тільки стан гонки.%bad
%enditem


\vspace{5mm}
%beginitem
4. Виберіть усі правильні твердження (якщо такі є).\par
\vspace{2mm}

( 1 ) У синтаксично коректній програмі мовою С++ стан гонки (race condition) може виникати.%good

( 2 ) Значення, що є результатом виконання функції, запущеної через async, зберігається безпосередньо в об’єкті ф’ючерса (future).%bad

( 3 ) Довільний потік може розблокувати заблокований м’ютекс.%bad

( 4 ) Потік, запущений конструктором класу jthread, має обов’язково зупинитися, якщо пов’язаний об’єкт надішле команду зупинитися.%bad
%enditem


\vspace{5mm}
%beginitem
5. Виберіть усі правильні твердження (якщо такі є).\par
\vspace{2mm}

( 1 ) Вільних від очікувань структур даних не існує.%bad

( 2 ) Структура даних без блокувань гарантує завершення виконання довільної операції над цією структурою даних за обмежений час.%bad

( 3 ) Структура даних з блокуваннями може бути такою, що використовує атомарні змінні.%good

( 4 ) Структура даних без блокувань зобов’язана використовувати атомарні змінні.%bad
\newpage

%enditem


\vspace{5mm}
%beginitem
6. Виберіть усі правильні твердження (якщо такі є).\par
\vspace{2mm}

( 1 ) Кожен потік програми може мати доступ до динамічних даних програми.%good

( 2 ) Динамічні змінні не можуть розділятися між потоками.%bad

( 3 ) Потік мови С++ не може отримати доступ до даних іншого потоку.%bad

( 4 ) Потік може мати доступ до автоматичних змінних іншого потоку і це безпечно.%bad
%enditem


\vspace{5mm}
%beginitem
7. У цьому фрагменті коду
\begin{verbatim}
template <class T> void f(T a1, T a2){}

void g(){
    int i=5;
    f(unique_ptr<int>(new int(++i)), unique_ptr<int>(new int(++i)));
\end{verbatim}

( 1 ) є стан гонки%bad

( 2 ) можливий виток ресурсів%bad

( 3 ) є гонка даних%bad

( 4 ) є невизначена поведінка%good
%enditem
}
%end
\end {large}

\vspace{4mm}
%end
\newpage
%begin
%begin
\begin {large}

\textbf{22.11.2024 Контрольна робота, варіант  \No { 22 }}\par


{
%beginitem
1. По яких днях тижня відбувались лекції з предмету?

%enditem


\vspace{5mm}
%beginitem
2. Як зовуть (ім'я, по батькові) викладача, що вів практичні заняття?

%enditem


\vspace{5mm}
%beginitem
3. Нехай int var=13; Один потік збільшив змінну var на 3, а інший збільшив її в три рази.

гонка даних (data race), стан гонки (race condition)

Виберіть усі правильні твердження (якщо такі є).\par
\vspace{2mm}
( 1 ) Тут може бути тільки стан гонки.%bad

( 2 ) Тут гарантовано є стан гонки, але гонки даних може не бути.%good

( 3 ) Тут гарантовано є і стан гонки, і гонка даних.%bad

( 4 ) Тут гарантовано нема ані гонки даних, ані стану гонки.%bad
%enditem


\vspace{5mm}
%beginitem
4. Виберіть усі правильні твердження (якщо такі є).\par
\vspace{2mm}

( 1 ) Значення, що є результатом виконання функції, запущеної через async, зберігається безпосередньо в об’єкті проміза (promise).%bad

( 2 ) Критична секція це неперервна ділянка коду програми.%bad

( 3 ) Критичну секцію можна реалізувати за допомогою м’ютексу.%good

( 4 ) Об’єкти класів promise та future можна використовувати тільки в багатопотокових програмах.%bad
%enditem


\vspace{5mm}
%beginitem
5. Виберіть усі правильні твердження (якщо такі є).\par
\vspace{2mm}

( 1 ) Вільна від блокувань структура даних гарантує, що жоден потік не буде очікувати доступу до виконання операцій над цією структурою даних.%bad

( 2 ) На операціях структури даних без блокувань взаємоблокувань потоків виникнути не може.%bad

( 3 ) Операція над вільною від блокувань структурою даних може виконуватись скільки завгодно довго.%good

( 4 ) Для уникнення гонки даних структура даних може використовувати не більше одного м’ютекса.%bad
\newpage

%enditem


\vspace{5mm}
%beginitem
6. Виберіть усі правильні твердження (якщо такі є).\par
\vspace{2mm}

( 1 ) Потік мови С++ не може отримати доступ до автоматичних змінних іншого потоку.%bad

( 2 ) Потік має доступ тільки до своїх власних даних та до глобальних статичних даних.%bad

( 3 ) Автоматичні змінні не можуть розділятися між потоками.%bad

( 4 ) Потік може мати доступ до автоматичних змінних іншого потоку, але це небезпечно.%good
%enditem


\vspace{5mm}
%beginitem
7. У цьому фрагменті коду
\begin{verbatim}
template <class T> void f(T a1, T a2){}

void g(){
    int i=5;
    f(unique_ptr<int>(new int(i++)), unique_ptr<int>(new int(++i)));
\end{verbatim}

( 1 ) можливий виток ресурсів%bad

( 2 ) є невизначена поведінка%good

( 3 ) є гонка даних%bad

( 4 ) є стан гонки%bad
%enditem
}
%end
\end {large}

\vspace{4mm}
%end
\newpage
%begin
%begin
\begin {large}

\textbf{22.11.2024 Контрольна робота, варіант  \No { 23 }}\par


{
%beginitem
1. По яких днях тижня відбувались лекції з предмету?

%enditem


\vspace{5mm}
%beginitem
2. Як зовуть (ім'я, по батькові) викладача, що вів практичні заняття?

%enditem


\vspace{5mm}
%beginitem
3. Нехай int var=13; Один потік збільшив змінну var на 3, а інший збільшив її в три рази.

гонка даних (data race), стан гонки (race condition)

Виберіть усі правильні твердження (якщо такі є).\par
\vspace{2mm}
( 1 ) Тут гарантовано є стан гонки.%good

( 2 ) Тут гарантовано нема стану гонки, але може бути гонка даних.%bad

( 3 ) Тут може не бути стану гонки.%bad

( 4 ) Тут гарантовано є гонка даних.%bad
%enditem


\vspace{5mm}
%beginitem
4. Виберіть усі правильні твердження (якщо такі є).\par
\vspace{2mm}

( 1 ) Тільки від компілятора залежить, чи запустить функція async виконання в новому потоці.%bad

( 2 ) Розблоковувати незаблокований м’ютекс можна.%bad

( 3 ) Виклик \code{this\_thread::sleep\_for(5s)} змушує потік проспати рівно 5 секунд.%bad

( 4 ) Розблоковувати незаблокований м’ютекс не можна.%good
%enditem


\vspace{5mm}
%beginitem
5. Виберіть усі правильні твердження (якщо такі є).\par
\vspace{2mm}

( 1 ) Структура даних без блокувань гарантує, що жоден потік не буде очікувати доступу до виконання операцій над цією структурою даних.%bad

( 2 ) Операція над структурою даних без блокувань завжди виконається за обмежений час.%bad

( 3 ) За використання атомарних змінних структура даних є структурою даних без блокувань.%bad

( 4 ) За використання структурою даних рівно одного м’ютексу, взаємоблокування на її операціях неможливі.%good
\newpage

%enditem


\vspace{5mm}
%beginitem
6. Виберіть усі правильні твердження (якщо такі є).\par
\vspace{2mm}

( 1 ) Кожен потік програми може мати доступ до динамічних даних програми.%good

( 2 ) Потік може мати доступ до автоматичних змінних іншого потоку і це безпечно.%bad

( 3 ) Потік мови С++ не може отримати доступ до автоматичних змінних іншого потоку.%bad

( 4 ) Потік мови С++ не може отримати доступ до даних іншого потоку.%bad
%enditem


\vspace{5mm}
%beginitem
7. У цьому фрагменті коду
\begin{verbatim}
template <class T> void f(T a1, T a2){}

void g(){
    int i=5;
    f(unique_ptr<int>(new int(++i)), unique_ptr<int>(new int(++i)));
\end{verbatim}

( 1 ) є стан гонки%bad

( 2 ) можливий виток ресурсів%bad

( 3 ) є гонка даних%bad

( 4 ) є невизначена поведінка%good
%enditem
}
%end
\end {large}

\vspace{4mm}
%end
\newpage
%begin
%begin
\begin {large}

\textbf{22.11.2024 Контрольна робота, варіант  \No { 24 }}\par


{
%beginitem
1. По яких днях тижня відбувались лекції з предмету?

%enditem


\vspace{5mm}
%beginitem
2. Як зовуть (ім'я, по батькові) викладача, що вів практичні заняття?

%enditem


\vspace{5mm}
%beginitem
3. Нехай int var=13; Один потік збільшив змінну var на 3, а інший збільшив її в три рази.

гонка даних (data race), стан гонки (race condition)

Виберіть усі правильні твердження (якщо такі є).\par
\vspace{2mm}
( 1 ) Тут може не бути ані гонки даних, ані стану гонки.%bad

( 2 ) Тут може бути тільки гонка даних.%bad

( 3 ) Тут гарантовано нема гонки даних, але може бути стан гонки.%bad

( 4 ) Тут може не бути гонки даних.%good
%enditem


\vspace{5mm}
%beginitem
4. Виберіть усі правильні твердження (якщо такі є).\par
\vspace{2mm}

( 1 ) Якщо два обчислення програми на С++ тільки читають одну ту саму ділянку пам'яті, то на них гонка даних  виникнути не може.%good

( 2 ) Досвічений програміст за однопотоковою програмою (текстом мовою С++) завжди може вказати послідовність виконуваних обчислень на конкретних вхідних даних.%bad

( 3 ) Щоб розблокувати заблокований м’ютекс \code{recursive\_mutex}, його завжди достатньо розблокувати один раз.%bad

( 4 ) Функція \code{sleep\_for} може змусити інший потік заснути.%bad
%enditem


\vspace{5mm}
%beginitem
5. Виберіть усі правильні твердження (якщо такі є).\par
\vspace{2mm}

( 1 ) Реалізація структури даних без блокувань не може використовувати м’ютекси.%good

( 2 ) Реалізації різних операцій над структурою даних повинні використовувати різні м’ютекси.%bad

( 3 ) Структура даних є структурою даних без блокувань, якщо її реалізація не використовує м’ютекси.%bad

( 4 ) Мова С++20 має стандартну бібліотеку вільних від блокувань структур даних.%bad
\newpage

%enditem


\vspace{5mm}
%beginitem
6. Виберіть усі правильні твердження (якщо такі є).\par
\vspace{2mm}

( 1 ) Потік може мати доступ до автоматичних змінних іншого потоку, але це небезпечно.%good

( 2 ) Потік має доступ тільки до своїх власних даних та до глобальних статичних даних.%bad

( 3 ) Автоматичні змінні не можуть розділятися між потоками.%bad

( 4 ) Динамічні змінні не можуть розділятися між потоками.%bad
%enditem


\vspace{5mm}
%beginitem
7. У цьому фрагменті коду
\begin{verbatim}
template <class T> void f(T a){}

void g(){
    int i=5;
    f(unique_ptr<int>(new int(i)))+(unique_ptr<int>(new int(35)));
\end{verbatim}
( 1 ) можливий виток ресурсів%good

( 2 ) є невизначена поведінка%bad

( 3 ) є гонка даних%bad

( 4 ) є стан гонки%bad
%enditem
}
%end
\end {large}

\vspace{4mm}
%end
\newpage
%begin
%begin
\begin {large}

\textbf{22.11.2024 Контрольна робота, варіант  \No { 25 }}\par


{
%beginitem
1. По яких днях тижня відбувались лекції з предмету?

%enditem


\vspace{5mm}
%beginitem
2. Як зовуть (ім'я, по батькові) викладача, що вів практичні заняття?

%enditem


\vspace{5mm}
%beginitem
3. Нехай int var=13; Один потік збільшив змінну var на 3, а інший збільшив її в три рази.

гонка даних (data race), стан гонки (race condition)

Виберіть усі правильні твердження (якщо такі є).\par
\vspace{2mm}
( 1 ) Тут гарантовано є стан гонки.%good

( 2 ) Тут гарантовано є гонка даних.%bad

( 3 ) Тут гарантовано є і стан гонки, і гонка даних.%bad

( 4 ) Тут може бути тільки стан гонки.%bad
%enditem


\vspace{5mm}
%beginitem
4. Виберіть усі правильні твердження (якщо такі є).\par
\vspace{2mm}

( 1 ) Значення, що є результатом виконання функції, запущеної через async, зберігається безпосередньо в об’єкті проміза (promise).%bad

( 2 ) Коли потік спить, він не споживає ресурс процесора.%good

( 3 ) Досвічений програміст за однопотоковою програмою (текстом мовою С++) завжди може вказати послідовність виконуваних обчислень на конкретних вхідних даних.%bad

( 4 ) Критична секція це неперервна ділянка коду програми.%bad
%enditem


\vspace{5mm}
%beginitem
5. Виберіть усі правильні твердження (якщо такі є).\par
\vspace{2mm}

( 1 ) Структура даних без блокувань та вільна від блокувань структура даних --- це одне те саме.%bad

( 2 ) Стандартні контейнери stl безпечно використовувати з кількох потоків.%bad

( 3 ) Операція над вільною від блокувань структурою даних завжди виконається за обмежений час.%bad

( 4 ) Реалізація структури даних з блокуваннями може використовувати атомарні змінні.%good
\newpage

%enditem


\vspace{5mm}
%beginitem
6. Виберіть усі правильні твердження (якщо такі є).\par
\vspace{2mm}

( 1 ) Автоматичні змінні не можуть розділятися між потоками.%bad

( 2 ) Потік має доступ тільки до своїх власних даних та до глобальних статичних даних.%bad

( 3 ) Потік мови С++ не може отримати доступ до даних іншого потоку.%bad

( 4 ) Кожен потік програми може мати доступ до динамічних даних програми.%good
%enditem


\vspace{5mm}
%beginitem
7. У цьому фрагменті коду
\begin{verbatim}
template <class T> void f(T a1, T a2){}

void g(){
    int i=5;
    f(unique_ptr<int>(new int(++i)), unique_ptr<int>(new int(++i)));
\end{verbatim}

( 1 ) можливий виток ресурсів%bad

( 2 ) є невизначена поведінка%good

( 3 ) є стан гонки%bad

( 4 ) є гонка даних%bad
%enditem
}
%end
\end {large}

\vspace{4mm}
%end
\newpage
%begin
%begin
\begin {large}

\textbf{22.11.2024 Контрольна робота, варіант  \No { 26 }}\par


{
%beginitem
1. По яких днях тижня відбувались лекції з предмету?

%enditem


\vspace{5mm}
%beginitem
2. Як зовуть (ім'я, по батькові) викладача, що вів практичні заняття?

%enditem


\vspace{5mm}
%beginitem
3. Нехай int var=13; Один потік збільшив змінну var на 3, а інший збільшив її в три рази.

гонка даних (data race), стан гонки (race condition)

Виберіть усі правильні твердження (якщо такі є).\par
\vspace{2mm}
( 1 ) Тут гарантовано нема гонки даних, але може бути стан гонки.%bad

( 2 ) Тут може бути тільки гонка даних.%bad

( 3 ) Тут гарантовано є стан гонки, але гонки даних може не бути.%good

( 4 ) Тут гарантовано нема стану гонки, але може бути гонка даних.%bad
%enditem


\vspace{5mm}
%beginitem
4. Виберіть усі правильні твердження (якщо такі є).\par
\vspace{2mm}

( 1 ) Об’єкти класів promise та future можна використовувати тільки в багатопотокових програмах.%bad

( 2 ) Тільки той потік, який обнулив значення семафора, може його збільшити.%bad

( 3 ) Потік в активному очікуванні споживає ресурс процесора.%good

( 4 ) Щоб розблокувати заблокований м’ютекс \code{recursive\_mutex}, його завжди достатньо розблокувати один раз.%bad
%enditem


\vspace{5mm}
%beginitem
5. Виберіть усі правильні твердження (якщо такі є).\par
\vspace{2mm}

( 1 ) Вільних від очікувань структур даних не існує.%bad

( 2 ) Операція над структурою даних без блокувань завжди виконається за обмежений час.%bad

( 3 ) Структура даних без блокувань може не бути вільною від блокувань.%good

( 4 ) Структура даних є структурою даних без блокувань, якщо її реалізація не використовує м’ютекси.%bad
\newpage

%enditem


\vspace{5mm}
%beginitem
6. Виберіть усі правильні твердження (якщо такі є).\par
\vspace{2mm}

( 1 ) Потік може мати доступ до автоматичних змінних іншого потоку і це безпечно.%bad

( 2 ) Динамічні змінні не можуть розділятися між потоками.%bad

( 3 ) Потік мови С++ не може отримати доступ до автоматичних змінних іншого потоку.%bad

( 4 ) Потік може мати доступ до автоматичних змінних іншого потоку, але це небезпечно.%good
%enditem


\vspace{5mm}
%beginitem
7. У цьому фрагменті коду
\begin{verbatim}
template <class T> void f(T a){}

void g(){
    int i=5;
    f(unique_ptr<int>(new int(i)))+(unique_ptr<int>(new int(++i)));
\end{verbatim}
( 1 ) є гонка даних%bad

( 2 ) можливий виток ресурсів%good

( 3 ) є стан гонки%bad

( 4 ) є невизначена поведінка%good
%enditem
}
%end
\end {large}

\vspace{4mm}
%end
\newpage
%begin
%begin
\begin {large}

\textbf{22.11.2024 Контрольна робота, варіант  \No { 27 }}\par


{
%beginitem
1. По яких днях тижня відбувались лекції з предмету?

%enditem


\vspace{5mm}
%beginitem
2. Як зовуть (ім'я, по батькові) викладача, що вів практичні заняття?

%enditem


\vspace{5mm}
%beginitem
3. Нехай int var=13; Один потік збільшив змінну var на 3, а інший збільшив її в три рази.

гонка даних (data race), стан гонки (race condition)

Виберіть усі правильні твердження (якщо такі є).\par
\vspace{2mm}
( 1 ) Тут може не бути гонки даних.%good

( 2 ) Тут може не бути ані гонки даних, ані стану гонки.%bad

( 3 ) Тут може не бути стану гонки.%bad

( 4 ) Тут гарантовано нема ані гонки даних, ані стану гонки.%bad
%enditem


\vspace{5mm}
%beginitem
4. Виберіть усі правильні твердження (якщо такі є).\par
\vspace{2mm}

( 1 ) М’ютекс має бути записаний на початку критичної секції, яку він забезпечує.%bad

( 2 ) За різних запусків програми на одних тих самих вхідних даних може виникати хронологічно різна послідовність обчислень.%good

( 3 ) Функція async запускає виконання в новому потоці.%bad

( 4 ) Об’єкт promise можна скопіювати та роздати копії різним потокам, щоб кожен з них міг передати результат своєї роботи.%bad
%enditem


\vspace{5mm}
%beginitem
5. Виберіть усі правильні твердження (якщо такі є).\par
\vspace{2mm}

( 1 ) Вільна від блокувань структура даних гарантує, що жоден потік не буде очікувати доступу до виконання операцій над цією структурою даних.%bad

( 2 ) Кожна вільна від очікувань структура даних є структурою даних, вільною від блокувань.%good

( 3 ) Реалізації різних операцій над структурою даних повинні використовувати різні м’ютекси.%bad

( 4 ) Мова С++20 має стандартну бібліотеку конкурентних структур даних.%bad
\newpage

%enditem


\vspace{5mm}
%beginitem
6. Виберіть усі правильні твердження (якщо такі є).\par
\vspace{2mm}

( 1 ) Потік може мати доступ до автоматичних змінних іншого потоку і це безпечно.%bad

( 2 ) Динамічні змінні не можуть розділятися між потоками.%bad

( 3 ) Потік має доступ тільки до своїх власних даних та до глобальних статичних даних.%bad

( 4 ) Потік може мати доступ до автоматичних змінних іншого потоку, але це небезпечно.%good
%enditem


\vspace{5mm}
%beginitem
7. У цьому фрагменті коду
\begin{verbatim}
template <class T> void f(T a1, T a2){}

void g(){
    f(unique_ptr<int>(new int(3)), unique_ptr<int>(new int(35)));
\end{verbatim}

( 1 ) є гонка даних%bad

( 2 ) можливий виток ресурсів%bad

( 3 ) є невизначена поведінка%bad

( 4 ) є стан гонки%bad
%enditem
}
%end
\end {large}

\vspace{4mm}
%end
\newpage
%begin
%begin
\begin {large}

\textbf{22.11.2024 Контрольна робота, варіант  \No { 28 }}\par


{
%beginitem
1. По яких днях тижня відбувались лекції з предмету?

%enditem


\vspace{5mm}
%beginitem
2. Як зовуть (ім'я, по батькові) викладача, що вів практичні заняття?

%enditem


\vspace{5mm}
%beginitem
3. Нехай int var=13; Один потік збільшив змінну var на 3, а інший збільшив її в три рази.

гонка даних (data race), стан гонки (race condition)

Виберіть усі правильні твердження (якщо такі є).\par
\vspace{2mm}
( 1 ) Тут гарантовано нема гонки даних, але може бути стан гонки.%bad

( 2 ) Тут гарантовано нема стану гонки, але може бути гонка даних.%bad

( 3 ) Тут може бути тільки стан гонки.%bad

( 4 ) Тут гарантовано є стан гонки, але гонки даних може не бути.%good
%enditem


\vspace{5mm}
%beginitem
4. Виберіть усі правильні твердження (якщо такі є).\par
\vspace{2mm}

( 1 ) Не можна блокувати заблокований м’ютекс mutex.%bad

( 2 ) Текст однопотокової програми мовою С++ однозначно визначає послідовність виконуваних обчислень.%bad

( 3 ) Секція є критичною, якщо під час її виконання не можна генерувати винятки.%bad

( 4 ) Порядок обчислень у потоці для одних тих самих вхідних даних може залежати від використовуваного компілятора.%good
%enditem


\vspace{5mm}
%beginitem
5. Виберіть усі правильні твердження (якщо такі є).\par
\vspace{2mm}

( 1 ) Структура даних без блокувань може не бути вільною від блокувань.%good

( 2 ) Мова С++20 має стандартну бібліотеку вільних від блокувань структур даних.%bad

( 3 ) Одночасно використовувати атомарних змінних більше, ніж регістрів процесора, у програмі не можна.%bad

( 4 ) Структура даних без блокувань гарантує завершення виконання довільної операції над цією структурою даних за обмежений час.%bad
\newpage

%enditem


\vspace{5mm}
%beginitem
6. Виберіть усі правильні твердження (якщо такі є).\par
\vspace{2mm}

( 1 ) Потік мови С++ не може отримати доступ до даних іншого потоку.%bad

( 2 ) Кожен потік програми може мати доступ до динамічних даних програми.%good

( 3 ) Автоматичні змінні не можуть розділятися між потоками.%bad

( 4 ) Потік мови С++ не може отримати доступ до автоматичних змінних іншого потоку.%bad
%enditem


\vspace{5mm}
%beginitem
7. У цьому фрагменті коду
\begin{verbatim}
template <class T> void f(T a){}

void g(){
    f(unique_ptr<int>(new int(3)))+(unique_ptr<int>(new int(35)));
\end{verbatim}

( 1 ) можливий виток ресурсів%good

( 2 ) є гонка даних%bad

( 3 ) є невизначена поведінка%bad

( 4 ) є стан гонки%bad
%enditem
}
%end
\end {large}

\vspace{4mm}
%end
\newpage
%begin
%begin
\begin {large}

\textbf{22.11.2024 Контрольна робота, варіант  \No { 29 }}\par


{
%beginitem
1. По яких днях тижня відбувались лекції з предмету?

%enditem


\vspace{5mm}
%beginitem
2. Як зовуть (ім'я, по батькові) викладача, що вів практичні заняття?

%enditem


\vspace{5mm}
%beginitem
3. Нехай int var=13; Один потік збільшив змінну var на 3, а інший збільшив її в три рази.

гонка даних (data race), стан гонки (race condition)

Виберіть усі правильні твердження (якщо такі є).\par
\vspace{2mm}
( 1 ) Тут може бути тільки гонка даних.%bad

( 2 ) Тут гарантовано є стан гонки.%good

( 3 ) Тут гарантовано нема ані гонки даних, ані стану гонки.%bad

( 4 ) Тут може не бути стану гонки.%bad
%enditem


\vspace{5mm}
%beginitem
4. Виберіть усі правильні твердження (якщо такі є).\par
\vspace{2mm}

( 1 ) Виклик функції \code{this\_thread::sleep\_for(5s)} змушує потік, з якого його здійснили, заснути якнайменш на 5 секунд.%good

( 2 ) Об’єкт future можна скопіювати та роздати копії різним потокам, щоб вони змогли побачити результат.%bad

( 3 ) Розблоковувати незаблокований м’ютекс можна.%bad

( 4 ) Виклик \code{this\_thread::sleep\_for(5s)} змушує потік проспати рівно 5 секунд.%bad
%enditem


\vspace{5mm}
%beginitem
5. Виберіть усі правильні твердження (якщо такі є).\par
\vspace{2mm}

( 1 ) Для уникнення гонки даних структура даних може використовувати не більше одного м’ютекса.%bad

( 2 ) Структура даних без блокувань гарантує, що жоден потік не буде очікувати доступу до виконання операцій над цією структурою даних.%bad

( 3 ) Структура даних без блокувань зобов’язана використовувати атомарні змінні.%bad

( 4 ) Кожна вільна від очікувань структура даних є структурою даних, вільною від блокувань.%good
\newpage

%enditem


\vspace{5mm}
%beginitem
6. Виберіть усі правильні твердження (якщо такі є).\par
\vspace{2mm}

( 1 ) Потік мови С++ не може отримати доступ до даних іншого потоку.%bad

( 2 ) Кожен потік програми може мати доступ до динамічних даних програми.%good

( 3 ) Потік мови С++ не може отримати доступ до автоматичних змінних іншого потоку.%bad

( 4 ) Автоматичні змінні не можуть розділятися між потоками.%bad
%enditem


\vspace{5mm}
%beginitem
7. У цьому фрагменті коду
\begin{verbatim}
template <class T> void f(T a1, T a2){}

void g(){
    int i=5;
    f(unique_ptr<int>(new int(++i)), unique_ptr<int>(new int(++i)));
\end{verbatim}

( 1 ) є стан гонки%bad

( 2 ) є невизначена поведінка%good

( 3 ) можливий виток ресурсів%bad

( 4 ) є гонка даних%bad
%enditem
}
%end
\end {large}

\vspace{4mm}
%end
\newpage
%begin
%begin
\begin {large}

\textbf{22.11.2024 Контрольна робота, варіант  \No { 30 }}\par


{
%beginitem
1. По яких днях тижня відбувались лекції з предмету?

%enditem


\vspace{5mm}
%beginitem
2. Як зовуть (ім'я, по батькові) викладача, що вів практичні заняття?

%enditem


\vspace{5mm}
%beginitem
3. Нехай int var=13; Один потік збільшив змінну var на 3, а інший збільшив її в три рази.

гонка даних (data race), стан гонки (race condition)

Виберіть усі правильні твердження (якщо такі є).\par
\vspace{2mm}
( 1 ) Тут може не бути ані гонки даних, ані стану гонки.%bad

( 2 ) Тут гарантовано є і стан гонки, і гонка даних.%bad

( 3 ) Тут може не бути гонки даних.%good

( 4 ) Тут гарантовано є гонка даних.%bad
%enditem


\vspace{5mm}
%beginitem
4. Виберіть усі правильні твердження (якщо такі є).\par
\vspace{2mm}

( 1 ) Щоб розблокувати заблокований \code{shared\_mutex}, його завжди достатньо розблокувати один раз.%good

( 2 ) Клас thread має засоби, щоб повернути результат виконання запущеного потоку.%bad

( 3 ) У стані гонки програма відчайдушно намагається захопити ресурси процесора.%bad

( 4 ) Потік, запущений конструктором класу jthread, має обов’язково зупинитися, якщо пов’язаний об’єкт надішле команду зупинитися.%bad
%enditem


\vspace{5mm}
%beginitem
5. Виберіть усі правильні твердження (якщо такі є).\par
\vspace{2mm}

( 1 ) Взаємоблокування потоків може виникнути тільки за використання м’ютексів.%bad

( 2 ) Реалізація структури даних без блокувань не може використовувати м’ютекси.%good

( 3 ) На операціях структури даних без блокувань взаємоблокувань потоків виникнути не може.%bad

( 4 ) За використання атомарних змінних структура даних є структурою даних без блокувань.%bad
\newpage

%enditem


\vspace{5mm}
%beginitem
6. Виберіть усі правильні твердження (якщо такі є).\par
\vspace{2mm}

( 1 ) Потік має доступ тільки до своїх власних даних та до глобальних статичних даних.%bad

( 2 ) Динамічні змінні не можуть розділятися між потоками.%bad

( 3 ) Потік може мати доступ до автоматичних змінних іншого потоку, але це небезпечно.%good

( 4 ) Потік може мати доступ до автоматичних змінних іншого потоку і це безпечно.%bad
%enditem


\vspace{5mm}
%beginitem
7. У цьому фрагменті коду
\begin{verbatim}
template <class T> void f(T a){}

void g(){
    int i=5;
    f(unique_ptr<int>(new int(i)))+(unique_ptr<int>(new int(++i)));
\end{verbatim}
( 1 ) є невизначена поведінка%good

( 2 ) є гонка даних%bad

( 3 ) можливий виток ресурсів%good

( 4 ) є стан гонки%bad
%enditem
}
%end
\end {large}

\vspace{4mm}
%end
\newpage
%begin
%begin
\begin {large}

\textbf{22.11.2024 Контрольна робота, варіант  \No { 31 }}\par


{
%beginitem
1. По яких днях тижня відбувались лекції з предмету?

%enditem


\vspace{5mm}
%beginitem
2. Як зовуть (ім'я, по батькові) викладача, що вів практичні заняття?

%enditem


\vspace{5mm}
%beginitem
3. Нехай int var=13; Один потік збільшив змінну var на 3, а інший збільшив її в три рази.

гонка даних (data race), стан гонки (race condition)

Виберіть усі правильні твердження (якщо такі є).\par
\vspace{2mm}
( 1 ) Тут гарантовано нема стану гонки, але може бути гонка даних.%bad

( 2 ) Тут гарантовано є стан гонки, але гонки даних може не бути.%good

( 3 ) Тут гарантовано є і стан гонки, і гонка даних.%bad

( 4 ) Тут може бути тільки стан гонки.%bad
%enditem


\vspace{5mm}
%beginitem
4. Виберіть усі правильні твердження (якщо такі є).\par
\vspace{2mm}

( 1 ) У програмі може бути не більше однієї критичної секції.%bad

( 2 ) Функція \code{sleep\_for} може змусити інший потік заснути.%bad

( 3 ) Довільний потік може розблокувати заблокований м’ютекс.%bad

( 4 ) Коли потік спить, він не споживає ресурс процесора.%good
%enditem


\vspace{5mm}
%beginitem
5. Виберіть усі правильні твердження (якщо такі є).\par
\vspace{2mm}

( 1 ) За використання структурою даних рівно одного м’ютексу, взаємоблокування на її операціях неможливі.%good

( 2 ) Реалізації різних операцій над структурою даних повинні використовувати різні м’ютекси.%bad

( 3 ) Вільна від блокувань структура даних гарантує, що жоден потік не буде очікувати доступу до виконання операцій над цією структурою даних.%bad

( 4 ) Структура даних без блокувань гарантує, що жоден потік не буде очікувати доступу до виконання операцій над цією структурою даних.%bad
\newpage

%enditem


\vspace{5mm}
%beginitem
6. Виберіть усі правильні твердження (якщо такі є).\par
\vspace{2mm}

( 1 ) Кожен потік програми може мати доступ до динамічних даних програми.%good

( 2 ) Динамічні змінні не можуть розділятися між потоками.%bad

( 3 ) Потік може мати доступ до автоматичних змінних іншого потоку і це безпечно.%bad

( 4 ) Потік мови С++ не може отримати доступ до автоматичних змінних іншого потоку.%bad
%enditem


\vspace{5mm}
%beginitem
7. У цьому фрагменті коду
\begin{verbatim}
template <class T> void f(T a){}

void g(){
    int i=5;
    f(unique_ptr<int>(new int(i)))+(unique_ptr<int>(new int(35)));
\end{verbatim}
( 1 ) є стан гонки%bad

( 2 ) є невизначена поведінка%bad

( 3 ) є гонка даних%bad

( 4 ) можливий виток ресурсів%good
%enditem
}
%end
\end {large}

\vspace{4mm}
%end
\newpage
%begin
%begin
\begin {large}

\textbf{22.11.2024 Контрольна робота, варіант  \No { 32 }}\par


{
%beginitem
1. По яких днях тижня відбувались лекції з предмету?

%enditem


\vspace{5mm}
%beginitem
2. Як зовуть (ім'я, по батькові) викладача, що вів практичні заняття?

%enditem


\vspace{5mm}
%beginitem
3. Нехай int var=13; Один потік збільшив змінну var на 3, а інший збільшив її в три рази.

гонка даних (data race), стан гонки (race condition)

Виберіть усі правильні твердження (якщо такі є).\par
\vspace{2mm}
( 1 ) Тут гарантовано є гонка даних.%bad

( 2 ) Тут може не бути гонки даних.%good

( 3 ) Тут гарантовано нема ані гонки даних, ані стану гонки.%bad

( 4 ) Тут може не бути ані гонки даних, ані стану гонки.%bad
%enditem


\vspace{5mm}
%beginitem
4. Виберіть усі правильні твердження (якщо такі є).\par
\vspace{2mm}

( 1 ) Тільки від компілятора залежить, чи запустить функція async виконання в новому потоці.%bad

( 2 ) Щоб розблокувати заблокований \code{shared\_mutex}, його завжди достатньо розблокувати один раз.%good

( 3 ) Коли програма спить, то вона не виконується.%bad

( 4 ) Значення, що є результатом виконання функції, запущеної через async, зберігається безпосередньо в об’єкті ф’ючерса (future).%bad
%enditem


\vspace{5mm}
%beginitem
5. Виберіть усі правильні твердження (якщо такі є).\par
\vspace{2mm}

( 1 ) Мова С++20 має стандартну бібліотеку конкурентних структур даних.%bad

( 2 ) Структура даних без блокувань та вільна від блокувань структура даних --- це одне те саме.%bad

( 3 ) Операція над вільною від блокувань структурою даних завжди виконається за обмежений час.%bad

( 4 ) Операція над вільною від блокувань структурою даних може виконуватись скільки завгодно довго.%good
\newpage

%enditem


\vspace{5mm}
%beginitem
6. Виберіть усі правильні твердження (якщо такі є).\par
\vspace{2mm}

( 1 ) Потік має доступ тільки до своїх власних даних та до глобальних статичних даних.%bad

( 2 ) Автоматичні змінні не можуть розділятися між потоками.%bad

( 3 ) Потік може мати доступ до автоматичних змінних іншого потоку, але це небезпечно.%good

( 4 ) Потік мови С++ не може отримати доступ до даних іншого потоку.%bad
%enditem


\vspace{5mm}
%beginitem
7. У цьому фрагменті коду
\begin{verbatim}
template <class T> void f(T a1, T a2){}

void g(){
    int i=5;
    f(unique_ptr<int>(new int(i++)), unique_ptr<int>(new int(++i)));
\end{verbatim}

( 1 ) можливий виток ресурсів%bad

( 2 ) є гонка даних%bad

( 3 ) є невизначена поведінка%good

( 4 ) є стан гонки%bad
%enditem
}
%end
\end {large}

\vspace{4mm}
%end
\newpage
%begin
%begin
\begin {large}

\textbf{22.11.2024 Контрольна робота, варіант  \No { 33 }}\par


{
%beginitem
1. По яких днях тижня відбувались лекції з предмету?

%enditem


\vspace{5mm}
%beginitem
2. Як зовуть (ім'я, по батькові) викладача, що вів практичні заняття?

%enditem


\vspace{5mm}
%beginitem
3. Нехай int var=13; Один потік збільшив змінну var на 3, а інший збільшив її в три рази.

гонка даних (data race), стан гонки (race condition)

Виберіть усі правильні твердження (якщо такі є).\par
\vspace{2mm}
( 1 ) Тут може бути тільки гонка даних.%bad

( 2 ) Тут гарантовано нема гонки даних, але може бути стан гонки.%bad

( 3 ) Тут може не бути стану гонки.%bad

( 4 ) Тут гарантовано є стан гонки.%good
%enditem


\vspace{5mm}
%beginitem
4. Виберіть усі правильні твердження (якщо такі є).\par
\vspace{2mm}

( 1 ) У стані гонки програма відчайдушно намагається захопити ресурси процесора.%bad

( 2 ) Порядок обчислень у потоці для одних тих самих вхідних даних може залежати від використовуваного компілятора.%good

( 3 ) Об’єкт future можна скопіювати та роздати копії різним потокам, щоб вони змогли побачити результат.%bad

( 4 ) Коли програма спить, то вона не виконується.%bad
%enditem


\vspace{5mm}
%beginitem
5. Виберіть усі правильні твердження (якщо такі є).\par
\vspace{2mm}

( 1 ) Структура даних з блокуваннями може бути такою, що використовує атомарні змінні.%good

( 2 ) На операціях структури даних без блокувань взаємоблокувань потоків виникнути не може.%bad

( 3 ) Операція над структурою даних без блокувань завжди виконається за обмежений час.%bad

( 4 ) Для уникнення гонки даних структура даних може використовувати не більше одного м’ютекса.%bad
\newpage

%enditem


\vspace{5mm}
%beginitem
6. Виберіть усі правильні твердження (якщо такі є).\par
\vspace{2mm}

( 1 ) Потік може мати доступ до автоматичних змінних іншого потоку і це безпечно.%bad

( 2 ) Кожен потік програми може мати доступ до динамічних даних програми.%good

( 3 ) Потік має доступ тільки до своїх власних даних та до глобальних статичних даних.%bad

( 4 ) Потік мови С++ не може отримати доступ до даних іншого потоку.%bad
%enditem


\vspace{5mm}
%beginitem
7. У цьому фрагменті коду
\begin{verbatim}
template <class T> void f(T a){}

void g(){
    f(unique_ptr<int>(new int(3)))+(unique_ptr<int>(new int(35)));
\end{verbatim}

( 1 ) є невизначена поведінка%bad

( 2 ) є стан гонки%bad

( 3 ) можливий виток ресурсів%good

( 4 ) є гонка даних%bad
%enditem
}
%end
\end {large}

\vspace{4mm}
%end
\newpage
%begin
%begin
\begin {large}

\textbf{22.11.2024 Контрольна робота, варіант  \No { 34 }}\par


{
%beginitem
1. По яких днях тижня відбувались лекції з предмету?

%enditem


\vspace{5mm}
%beginitem
2. Як зовуть (ім'я, по батькові) викладача, що вів практичні заняття?

%enditem


\vspace{5mm}
%beginitem
3. Нехай int var=13; Один потік збільшив змінну var на 3, а інший збільшив її в три рази.

гонка даних (data race), стан гонки (race condition)

Виберіть усі правильні твердження (якщо такі є).\par
\vspace{2mm}
( 1 ) Тут може бути тільки стан гонки.%bad

( 2 ) Тут гарантовано є гонка даних.%bad

( 3 ) Тут гарантовано є стан гонки.%good

( 4 ) Тут гарантовано нема стану гонки, але може бути гонка даних.%bad
%enditem


\vspace{5mm}
%beginitem
4. Виберіть усі правильні твердження (якщо такі є).\par
\vspace{2mm}

( 1 ) Об’єкт promise можна скопіювати та роздати копії різним потокам, щоб кожен з них міг передати результат своєї роботи.%bad

( 2 ) Функція async запускає виконання в новому потоці.%bad

( 3 ) Щоб розблокувати заблокований м’ютекс \code{recursive\_mutex}, його завжди достатньо розблокувати один раз.%bad

( 4 ) За різних запусків програми на одних тих самих вхідних даних може виникати хронологічно різна послідовність обчислень.%good
%enditem


\vspace{5mm}
%beginitem
5. Виберіть усі правильні твердження (якщо такі є).\par
\vspace{2mm}

( 1 ) Структура даних може використовувати м’ютекси та атомарні змінні одночасно.%good

( 2 ) Взаємоблокування потоків може виникнути тільки за використання м’ютексів.%bad

( 3 ) Мова С++20 має стандартну бібліотеку вільних від блокувань структур даних.%bad

( 4 ) Одночасно використовувати атомарних змінних більше, ніж регістрів процесора, у програмі не можна.%bad
\newpage

%enditem


\vspace{5mm}
%beginitem
6. Виберіть усі правильні твердження (якщо такі є).\par
\vspace{2mm}

( 1 ) Потік може мати доступ до автоматичних змінних іншого потоку, але це небезпечно.%good

( 2 ) Автоматичні змінні не можуть розділятися між потоками.%bad

( 3 ) Потік мови С++ не може отримати доступ до автоматичних змінних іншого потоку.%bad

( 4 ) Динамічні змінні не можуть розділятися між потоками.%bad
%enditem


\vspace{5mm}
%beginitem
7. У цьому фрагменті коду
\begin{verbatim}
template <class T> void f(T a1, T a2){}

void g(){
    f(unique_ptr<int>(new int(3)), unique_ptr<int>(new int(35)));
\end{verbatim}

( 1 ) можливий виток ресурсів%bad

( 2 ) є стан гонки%bad

( 3 ) є невизначена поведінка%bad

( 4 ) є гонка даних%bad
%enditem
}
%end
\end {large}

\vspace{4mm}
%end
\newpage
%begin
%begin
\begin {large}

\textbf{22.11.2024 Контрольна робота, варіант  \No { 35 }}\par


{
%beginitem
1. По яких днях тижня відбувались лекції з предмету?

%enditem


\vspace{5mm}
%beginitem
2. Як зовуть (ім'я, по батькові) викладача, що вів практичні заняття?

%enditem


\vspace{5mm}
%beginitem
3. Нехай int var=13; Один потік збільшив змінну var на 3, а інший збільшив її в три рази.

гонка даних (data race), стан гонки (race condition)

Виберіть усі правильні твердження (якщо такі є).\par
\vspace{2mm}
( 1 ) Тут може не бути гонки даних.%good

( 2 ) Тут гарантовано нема гонки даних, але може бути стан гонки.%bad

( 3 ) Тут може не бути ані гонки даних, ані стану гонки.%bad

( 4 ) Тут гарантовано є і стан гонки, і гонка даних.%bad
%enditem


\vspace{5mm}
%beginitem
4. Виберіть усі правильні твердження (якщо такі є).\par
\vspace{2mm}

( 1 ) Досвічений програміст за однопотоковою програмою (текстом мовою С++) завжди може вказати послідовність виконуваних обчислень на конкретних вхідних даних.%bad

( 2 ) Виклик \code{this\_thread::sleep\_for(5s)} змушує потік проспати рівно 5 секунд.%bad

( 3 ) Розблокувати м’ютекс може тільки той потік, що його заблокував.%good

( 4 ) Функція \code{sleep\_for} може змусити інший потік заснути.%bad
%enditem


\vspace{5mm}
%beginitem
5. Виберіть усі правильні твердження (якщо такі є).\par
\vspace{2mm}

( 1 ) Структура даних без блокувань зобов’язана використовувати атомарні змінні.%bad

( 2 ) Структура даних без блокувань гарантує завершення виконання довільної операції над цією структурою даних за обмежений час.%bad

( 3 ) Реалізація структури даних з блокуваннями може використовувати атомарні змінні.%good

( 4 ) Структура даних є структурою даних без блокувань, якщо її реалізація не використовує м’ютекси.%bad
\newpage

%enditem


\vspace{5mm}
%beginitem
6. Виберіть усі правильні твердження (якщо такі є).\par
\vspace{2mm}

( 1 ) Потік має доступ тільки до своїх власних даних та до глобальних статичних даних.%bad

( 2 ) Потік мови С++ не може отримати доступ до даних іншого потоку.%bad

( 3 ) Потік може мати доступ до автоматичних змінних іншого потоку, але це небезпечно.%good

( 4 ) Динамічні змінні не можуть розділятися між потоками.%bad
%enditem


\vspace{5mm}
%beginitem
7. У цьому фрагменті коду
\begin{verbatim}
template <class T> void f(T a1, T a2){}

void g(){
    int i=5;
    f(unique_ptr<int>(new int(++i)), unique_ptr<int>(new int(++i)));
\end{verbatim}

( 1 ) є стан гонки%bad

( 2 ) є невизначена поведінка%good

( 3 ) є гонка даних%bad

( 4 ) можливий виток ресурсів%bad
%enditem
}
%end
\end {large}

\vspace{4mm}
%end
\newpage
%begin
%begin
\begin {large}

\textbf{22.11.2024 Контрольна робота, варіант  \No { 36 }}\par


{
%beginitem
1. По яких днях тижня відбувались лекції з предмету?

%enditem


\vspace{5mm}
%beginitem
2. Як зовуть (ім'я, по батькові) викладача, що вів практичні заняття?

%enditem


\vspace{5mm}
%beginitem
3. Нехай int var=13; Один потік збільшив змінну var на 3, а інший збільшив її в три рази.

гонка даних (data race), стан гонки (race condition)

Виберіть усі правильні твердження (якщо такі є).\par
\vspace{2mm}
( 1 ) Тут гарантовано є стан гонки, але гонки даних може не бути.%good

( 2 ) Тут гарантовано нема ані гонки даних, ані стану гонки.%bad

( 3 ) Тут може не бути стану гонки.%bad

( 4 ) Тут може бути тільки гонка даних.%bad
%enditem


\vspace{5mm}
%beginitem
4. Виберіть усі правильні твердження (якщо такі є).\par
\vspace{2mm}

( 1 ) Текст однопотокової програми мовою С++ однозначно визначає послідовність виконуваних обчислень.%bad

( 2 ) Не можна блокувати заблокований м’ютекс mutex.%bad

( 3 ) М’ютекс має бути записаний на початку критичної секції, яку він забезпечує.%bad

( 4 ) Критичну секцію можна реалізувати за допомогою м’ютексу.%good
%enditem


\vspace{5mm}
%beginitem
5. Виберіть усі правильні твердження (якщо такі є).\par
\vspace{2mm}

( 1 ) Стандартні контейнери stl безпечно використовувати з кількох потоків.%bad

( 2 ) Якщо весь код структури даних утворює одну критичну секцію, то така структура даних може використовуватись в конкурентній програмі, але по-справжньо\-му конкурентною не є.%good

( 3 ) За використання атомарних змінних структура даних є структурою даних без блокувань.%bad

( 4 ) Вільних від очікувань структур даних не існує.%bad
\newpage

%enditem


\vspace{5mm}
%beginitem
6. Виберіть усі правильні твердження (якщо такі є).\par
\vspace{2mm}

( 1 ) Потік мови С++ не може отримати доступ до автоматичних змінних іншого потоку.%bad

( 2 ) Автоматичні змінні не можуть розділятися між потоками.%bad

( 3 ) Потік може мати доступ до автоматичних змінних іншого потоку і це безпечно.%bad

( 4 ) Кожен потік програми може мати доступ до динамічних даних програми.%good
%enditem


\vspace{5mm}
%beginitem
7. У цьому фрагменті коду
\begin{verbatim}
template <class T> void f(T a){}

void g(){
    int i=5;
    f(unique_ptr<int>(new int(i)))+(unique_ptr<int>(new int(35)));
\end{verbatim}
( 1 ) є стан гонки%bad

( 2 ) можливий виток ресурсів%good

( 3 ) є невизначена поведінка%bad

( 4 ) є гонка даних%bad
%enditem
}
%end
\end {large}

\vspace{4mm}
%end
\newpage
%begin
%begin
\begin {large}

\textbf{22.11.2024 Контрольна робота, варіант  \No { 37 }}\par


{
%beginitem
1. По яких днях тижня відбувались лекції з предмету?

%enditem


\vspace{5mm}
%beginitem
2. Як зовуть (ім'я, по батькові) викладача, що вів практичні заняття?

%enditem


\vspace{5mm}
%beginitem
3. Нехай int var=13; Один потік збільшив змінну var на 3, а інший збільшив її в три рази.

гонка даних (data race), стан гонки (race condition)

Виберіть усі правильні твердження (якщо такі є).\par
\vspace{2mm}
( 1 ) Тут гарантовано є гонка даних.%bad

( 2 ) Тут гарантовано є стан гонки, але гонки даних може не бути.%good

( 3 ) Тут може бути тільки гонка даних.%bad

( 4 ) Тут може не бути стану гонки.%bad
%enditem


\vspace{5mm}
%beginitem
4. Виберіть усі правильні твердження (якщо такі є).\par
\vspace{2mm}

( 1 ) Потік, запущений конструктором класу jthread, має обов’язково зупинитися, якщо пов’язаний об’єкт надішле команду зупинитися.%bad

( 2 ) Значення, що є результатом виконання функції, запущеної через async, зберігається безпосередньо в об’єкті проміза (promise).%bad

( 3 ) Розблоковувати незаблокований м’ютекс не можна.%good

( 4 ) Клас thread має засоби, щоб повернути результат виконання запущеного потоку.%bad
%enditem


\vspace{5mm}
%beginitem
5. Виберіть усі правильні твердження (якщо такі є).\par
\vspace{2mm}

( 1 ) Стандартні контейнери stl безпечно використовувати з кількох потоків.%bad

( 2 ) Одночасно використовувати атомарних змінних більше, ніж регістрів процесора, у програмі не можна.%bad

( 3 ) Мова С++20 має стандартну бібліотеку конкурентних структур даних.%bad

( 4 ) Структура даних без блокувань може не бути вільною від блокувань.%good
\newpage

%enditem


\vspace{5mm}
%beginitem
6. Виберіть усі правильні твердження (якщо такі є).\par
\vspace{2mm}

( 1 ) Кожен потік програми може мати доступ до динамічних даних програми.%good

( 2 ) Автоматичні змінні не можуть розділятися між потоками.%bad

( 3 ) Потік мови С++ не може отримати доступ до автоматичних змінних іншого потоку.%bad

( 4 ) Потік може мати доступ до автоматичних змінних іншого потоку і це безпечно.%bad
%enditem


\vspace{5mm}
%beginitem
7. У цьому фрагменті коду
\begin{verbatim}
template <class T> void f(T a1, T a2){}

void g(){
    f(unique_ptr<int>(new int(3)), unique_ptr<int>(new int(35)));
\end{verbatim}

( 1 ) є гонка даних%bad

( 2 ) є невизначена поведінка%bad

( 3 ) є стан гонки%bad

( 4 ) можливий виток ресурсів%bad
%enditem
}
%end
\end {large}

\vspace{4mm}
%end
\newpage
%begin
%begin
\begin {large}

\textbf{22.11.2024 Контрольна робота, варіант  \No { 38 }}\par


{
%beginitem
1. По яких днях тижня відбувались лекції з предмету?

%enditem


\vspace{5mm}
%beginitem
2. Як зовуть (ім'я, по батькові) викладача, що вів практичні заняття?

%enditem


\vspace{5mm}
%beginitem
3. Нехай int var=13; Один потік збільшив змінну var на 3, а інший збільшив її в три рази.

гонка даних (data race), стан гонки (race condition)

Виберіть усі правильні твердження (якщо такі є).\par
\vspace{2mm}
( 1 ) Тут гарантовано нема гонки даних, але може бути стан гонки.%bad

( 2 ) Тут гарантовано є стан гонки.%good

( 3 ) Тут гарантовано нема стану гонки, але може бути гонка даних.%bad

( 4 ) Тут може бути тільки стан гонки.%bad
%enditem


\vspace{5mm}
%beginitem
4. Виберіть усі правильні твердження (якщо такі є).\par
\vspace{2mm}

( 1 ) Довільний потік може розблокувати заблокований м’ютекс.%bad

( 2 ) Клас thread не має засобів, щоб повернути результат виконання запущеного потоку.%good

( 3 ) Секція є критичною, якщо під час її виконання не можна генерувати винятки.%bad

( 4 ) У програмі може бути не більше однієї критичної секції.%bad
%enditem


\vspace{5mm}
%beginitem
5. Виберіть усі правильні твердження (якщо такі є).\par
\vspace{2mm}

( 1 ) Вільних від очікувань структур даних не існує.%bad

( 2 ) Мова С++20 має стандартну бібліотеку вільних від блокувань структур даних.%bad

( 3 ) Структура даних без блокувань та вільна від блокувань структура даних --- це одне те саме.%bad

( 4 ) Структура даних може використовувати м’ютекси та атомарні змінні одночасно.%good
\newpage

%enditem


\vspace{5mm}
%beginitem
6. Виберіть усі правильні твердження (якщо такі є).\par
\vspace{2mm}

( 1 ) Динамічні змінні не можуть розділятися між потоками.%bad

( 2 ) Потік може мати доступ до автоматичних змінних іншого потоку, але це небезпечно.%good

( 3 ) Потік має доступ тільки до своїх власних даних та до глобальних статичних даних.%bad

( 4 ) Потік мови С++ не може отримати доступ до даних іншого потоку.%bad
%enditem


\vspace{5mm}
%beginitem
7. У цьому фрагменті коду
\begin{verbatim}
template <class T> void f(T a1, T a2){}

void g(){
    int i=5;
    f(unique_ptr<int>(new int(i++)), unique_ptr<int>(new int(++i)));
\end{verbatim}

( 1 ) є невизначена поведінка%good

( 2 ) є гонка даних%bad

( 3 ) є стан гонки%bad

( 4 ) можливий виток ресурсів%bad
%enditem
}
%end
\end {large}

\vspace{4mm}
%end
\newpage
%begin
%begin
\begin {large}

\textbf{22.11.2024 Контрольна робота, варіант  \No { 39 }}\par


{
%beginitem
1. По яких днях тижня відбувались лекції з предмету?

%enditem


\vspace{5mm}
%beginitem
2. Як зовуть (ім'я, по батькові) викладача, що вів практичні заняття?

%enditem


\vspace{5mm}
%beginitem
3. Нехай int var=13; Один потік збільшив змінну var на 3, а інший збільшив її в три рази.

гонка даних (data race), стан гонки (race condition)

Виберіть усі правильні твердження (якщо такі є).\par
\vspace{2mm}
( 1 ) Тут може не бути гонки даних.%good

( 2 ) Тут гарантовано нема ані гонки даних, ані стану гонки.%bad

( 3 ) Тут може не бути ані гонки даних, ані стану гонки.%bad

( 4 ) Тут гарантовано є і стан гонки, і гонка даних.%bad
%enditem


\vspace{5mm}
%beginitem
4. Виберіть усі правильні твердження (якщо такі є).\par
\vspace{2mm}

( 1 ) Значення, що є результатом виконання функції, запущеної через async, зберігається безпосередньо в об’єкті ф’ючерса (future).%bad

( 2 ) Чи запустить функція  async виконання в новому потоці можна задавати політикою.%good

( 3 ) Тільки той потік, який обнулив значення семафора, може його збільшити.%bad

( 4 ) Об’єкти класів promise та future можна використовувати тільки в багатопотокових програмах.%bad
%enditem


\vspace{5mm}
%beginitem
5. Виберіть усі правильні твердження (якщо такі є).\par
\vspace{2mm}

( 1 ) Взаємоблокування потоків може виникнути тільки за використання м’ютексів.%bad

( 2 ) Вільна від блокувань структура даних гарантує, що жоден потік не буде очікувати доступу до виконання операцій над цією структурою даних.%bad

( 3 ) Операція над структурою даних без блокувань завжди виконається за обмежений час.%bad

( 4 ) Операція над вільною від блокувань структурою даних може виконуватись скільки завгодно довго.%good
\newpage

%enditem


\vspace{5mm}
%beginitem
6. Виберіть усі правильні твердження (якщо такі є).\par
\vspace{2mm}

( 1 ) Потік мови С++ не може отримати доступ до автоматичних змінних іншого потоку.%bad

( 2 ) Потік може мати доступ до автоматичних змінних іншого потоку, але це небезпечно.%good

( 3 ) Динамічні змінні не можуть розділятися між потоками.%bad

( 4 ) Автоматичні змінні не можуть розділятися між потоками.%bad
%enditem


\vspace{5mm}
%beginitem
7. У цьому фрагменті коду
\begin{verbatim}
template <class T> void f(T a1, T a2){}

void g(){
    int i=5;
    f(unique_ptr<int>(new int(++i)), unique_ptr<int>(new int(++i)));
\end{verbatim}

( 1 ) є гонка даних%bad

( 2 ) є невизначена поведінка%good

( 3 ) є стан гонки%bad

( 4 ) можливий виток ресурсів%bad
%enditem
}
%end
\end {large}

\vspace{4mm}
%end
\newpage
%begin
%begin
\begin {large}

\textbf{22.11.2024 Контрольна робота, варіант  \No { 40 }}\par


{
%beginitem
1. По яких днях тижня відбувались лекції з предмету?

%enditem


\vspace{5mm}
%beginitem
2. Як зовуть (ім'я, по батькові) викладача, що вів практичні заняття?

%enditem


\vspace{5mm}
%beginitem
3. Нехай int var=13; Один потік збільшив змінну var на 3, а інший збільшив її в три рази.

гонка даних (data race), стан гонки (race condition)

Виберіть усі правильні твердження (якщо такі є).\par
\vspace{2mm}
( 1 ) Тут гарантовано є гонка даних.%bad

( 2 ) Тут гарантовано нема стану гонки, але може бути гонка даних.%bad

( 3 ) Тут гарантовано є стан гонки.%good

( 4 ) Тут гарантовано є і стан гонки, і гонка даних.%bad
%enditem


\vspace{5mm}
%beginitem
4. Виберіть усі правильні твердження (якщо такі є).\par
\vspace{2mm}

( 1 ) У синтаксично коректній програмі мовою С++ стан гонки (race condition) може виникати.%good

( 2 ) Розблоковувати незаблокований м’ютекс можна.%bad

( 3 ) Критична секція це неперервна ділянка коду програми.%bad

( 4 ) Тільки від компілятора залежить, чи запустить функція async виконання в новому потоці.%bad
%enditem


\vspace{5mm}
%beginitem
5. Виберіть усі правильні твердження (якщо такі є).\par
\vspace{2mm}

( 1 ) Структура даних без блокувань зобов’язана використовувати атомарні змінні.%bad

( 2 ) Структура даних без блокувань гарантує, що жоден потік не буде очікувати доступу до виконання операцій над цією структурою даних.%bad

( 3 ) Реалізація структури даних з блокуваннями може використовувати атомарні змінні.%good

( 4 ) Операція над вільною від блокувань структурою даних завжди виконається за обмежений час.%bad
\newpage

%enditem


\vspace{5mm}
%beginitem
6. Виберіть усі правильні твердження (якщо такі є).\par
\vspace{2mm}

( 1 ) Потік може мати доступ до автоматичних змінних іншого потоку і це безпечно.%bad

( 2 ) Потік має доступ тільки до своїх власних даних та до глобальних статичних даних.%bad

( 3 ) Потік мови С++ не може отримати доступ до даних іншого потоку.%bad

( 4 ) Кожен потік програми може мати доступ до динамічних даних програми.%good
%enditem


\vspace{5mm}
%beginitem
7. У цьому фрагменті коду
\begin{verbatim}
template <class T> void f(T a){}

void g(){
    int i=5;
    f(unique_ptr<int>(new int(i)))+(unique_ptr<int>(new int(++i)));
\end{verbatim}
( 1 ) можливий виток ресурсів%good

( 2 ) є гонка даних%bad

( 3 ) є стан гонки%bad

( 4 ) є невизначена поведінка%good
%enditem
}
%end
\end {large}

\vspace{4mm}
%end
\newpage
%begin
%begin
\begin {large}

\textbf{22.11.2024 Контрольна робота, варіант  \No { 41 }}\par


{
%beginitem
1. По яких днях тижня відбувались лекції з предмету?

%enditem


\vspace{5mm}
%beginitem
2. Як зовуть (ім'я, по батькові) викладача, що вів практичні заняття?

%enditem


\vspace{5mm}
%beginitem
3. Нехай int var=13; Один потік збільшив змінну var на 3, а інший збільшив її в три рази.

гонка даних (data race), стан гонки (race condition)

Виберіть усі правильні твердження (якщо такі є).\par
\vspace{2mm}
( 1 ) Тут може бути тільки гонка даних.%bad

( 2 ) Тут може не бути гонки даних.%good

( 3 ) Тут гарантовано нема ані гонки даних, ані стану гонки.%bad

( 4 ) Тут може не бути стану гонки.%bad
%enditem


\vspace{5mm}
%beginitem
4. Виберіть усі правильні твердження (якщо такі є).\par
\vspace{2mm}

( 1 ) Виклик \code{this\_thread::sleep\_for(5s)} змушує потік проспати рівно 5 секунд.%bad

( 2 ) Критична секція може складатися з кількох ділянок коду, які можуть розташовуватись де завгодно (не обов’язково зв’язний фрагмент коду).%good

( 3 ) Текст однопотокової програми мовою С++ однозначно визначає послідовність виконуваних обчислень.%bad

( 4 ) Критична секція це неперервна ділянка коду програми.%bad
%enditem


\vspace{5mm}
%beginitem
5. Виберіть усі правильні твердження (якщо такі є).\par
\vspace{2mm}

( 1 ) Структура даних є структурою даних без блокувань, якщо її реалізація не використовує м’ютекси.%bad

( 2 ) Кожна вільна від очікувань структура даних є структурою даних, вільною від блокувань.%good

( 3 ) За використання атомарних змінних структура даних є структурою даних без блокувань.%bad

( 4 ) Для уникнення гонки даних структура даних може використовувати не більше одного м’ютекса.%bad
\newpage

%enditem


\vspace{5mm}
%beginitem
6. Виберіть усі правильні твердження (якщо такі є).\par
\vspace{2mm}

( 1 ) Кожен потік програми може мати доступ до динамічних даних програми.%good

( 2 ) Потік мови С++ не може отримати доступ до автоматичних змінних іншого потоку.%bad

( 3 ) Потік може мати доступ до автоматичних змінних іншого потоку і це безпечно.%bad

( 4 ) Динамічні змінні не можуть розділятися між потоками.%bad
%enditem


\vspace{5mm}
%beginitem
7. У цьому фрагменті коду
\begin{verbatim}
template <class T> void f(T a){}

void g(){
    f(unique_ptr<int>(new int(3)))+(unique_ptr<int>(new int(35)));
\end{verbatim}

( 1 ) є стан гонки%bad

( 2 ) є гонка даних%bad

( 3 ) є невизначена поведінка%bad

( 4 ) можливий виток ресурсів%good
%enditem
}
%end
\end {large}

\vspace{4mm}
%end
\newpage
%begin
%begin
\begin {large}

\textbf{22.11.2024 Контрольна робота, варіант  \No { 42 }}\par


{
%beginitem
1. По яких днях тижня відбувались лекції з предмету?

%enditem


\vspace{5mm}
%beginitem
2. Як зовуть (ім'я, по батькові) викладача, що вів практичні заняття?

%enditem


\vspace{5mm}
%beginitem
3. Нехай int var=13; Один потік збільшив змінну var на 3, а інший збільшив її в три рази.

гонка даних (data race), стан гонки (race condition)

Виберіть усі правильні твердження (якщо такі є).\par
\vspace{2mm}
( 1 ) Тут гарантовано є стан гонки, але гонки даних може не бути.%good

( 2 ) Тут може не бути ані гонки даних, ані стану гонки.%bad

( 3 ) Тут може бути тільки стан гонки.%bad

( 4 ) Тут гарантовано нема гонки даних, але може бути стан гонки.%bad
%enditem


\vspace{5mm}
%beginitem
4. Виберіть усі правильні твердження (якщо такі є).\par
\vspace{2mm}

( 1 ) Якщо два обчислення програми на С++ тільки читають одну ту саму ділянку пам'яті, то на них гонка даних  виникнути не може.%good

( 2 ) Тільки той потік, який обнулив значення семафора, може його збільшити.%bad

( 3 ) Об’єкти класів promise та future можна використовувати тільки в багатопотокових програмах.%bad

( 4 ) Значення, що є результатом виконання функції, запущеної через async, зберігається безпосередньо в об’єкті ф’ючерса (future).%bad
%enditem


\vspace{5mm}
%beginitem
5. Виберіть усі правильні твердження (якщо такі є).\par
\vspace{2mm}

( 1 ) Реалізації різних операцій над структурою даних повинні використовувати різні м’ютекси.%bad

( 2 ) За використання структурою даних рівно одного м’ютексу, взаємоблокування на її операціях неможливі.%good

( 3 ) На операціях структури даних без блокувань взаємоблокувань потоків виникнути не може.%bad

( 4 ) Структура даних без блокувань гарантує завершення виконання довільної операції над цією структурою даних за обмежений час.%bad
\newpage

%enditem


\vspace{5mm}
%beginitem
6. Виберіть усі правильні твердження (якщо такі є).\par
\vspace{2mm}

( 1 ) Автоматичні змінні не можуть розділятися між потоками.%bad

( 2 ) Потік мови С++ не може отримати доступ до даних іншого потоку.%bad

( 3 ) Потік може мати доступ до автоматичних змінних іншого потоку, але це небезпечно.%good

( 4 ) Потік має доступ тільки до своїх власних даних та до глобальних статичних даних.%bad
%enditem


\vspace{5mm}
%beginitem
7. У цьому фрагменті коду
\begin{verbatim}
template <class T> void f(T a1, T a2){}

void g(){
    int i=5;
    f(unique_ptr<int>(new int(++i)), unique_ptr<int>(new int(++i)));
\end{verbatim}

( 1 ) є стан гонки%bad

( 2 ) можливий виток ресурсів%bad

( 3 ) є невизначена поведінка%good

( 4 ) є гонка даних%bad
%enditem
}
%end
\end {large}

\vspace{4mm}
%end
\newpage
%begin
%begin
\begin {large}

\textbf{22.11.2024 Контрольна робота, варіант  \No { 43 }}\par


{
%beginitem
1. По яких днях тижня відбувались лекції з предмету?

%enditem


\vspace{5mm}
%beginitem
2. Як зовуть (ім'я, по батькові) викладача, що вів практичні заняття?

%enditem


\vspace{5mm}
%beginitem
3. Нехай int var=13; Один потік збільшив змінну var на 3, а інший збільшив її в три рази.

гонка даних (data race), стан гонки (race condition)

Виберіть усі правильні твердження (якщо такі є).\par
\vspace{2mm}
( 1 ) Тут гарантовано є стан гонки.%good

( 2 ) Тут гарантовано є і стан гонки, і гонка даних.%bad

( 3 ) Тут може не бути стану гонки.%bad

( 4 ) Тут може не бути ані гонки даних, ані стану гонки.%bad
%enditem


\vspace{5mm}
%beginitem
4. Виберіть усі правильні твердження (якщо такі є).\par
\vspace{2mm}

( 1 ) Виклик функції \code{this\_thread::sleep\_for(5s)} змушує потік, з якого його здійснили, заснути якнайменш на 5 секунд.%good

( 2 ) М’ютекс має бути записаний на початку критичної секції, яку він забезпечує.%bad

( 3 ) Функція \code{sleep\_for} може змусити інший потік заснути.%bad

( 4 ) Потік, запущений конструктором класу jthread, має обов’язково зупинитися, якщо пов’язаний об’єкт надішле команду зупинитися.%bad
%enditem


\vspace{5mm}
%beginitem
5. Виберіть усі правильні твердження (якщо такі є).\par
\vspace{2mm}

( 1 ) Вільна від блокувань структура даних гарантує, що жоден потік не буде очікувати доступу до виконання операцій над цією структурою даних.%bad

( 2 ) Реалізація структури даних без блокувань не може використовувати м’ютекси.%good

( 3 ) Операція над структурою даних без блокувань завжди виконається за обмежений час.%bad

( 4 ) Структура даних без блокувань гарантує, що жоден потік не буде очікувати доступу до виконання операцій над цією структурою даних.%bad
\newpage

%enditem


\vspace{5mm}
%beginitem
6. Виберіть усі правильні твердження (якщо такі є).\par
\vspace{2mm}

( 1 ) Потік може мати доступ до автоматичних змінних іншого потоку і це безпечно.%bad

( 2 ) Потік має доступ тільки до своїх власних даних та до глобальних статичних даних.%bad

( 3 ) Потік може мати доступ до автоматичних змінних іншого потоку, але це небезпечно.%good

( 4 ) Потік мови С++ не може отримати доступ до даних іншого потоку.%bad
%enditem


\vspace{5mm}
%beginitem
7. У цьому фрагменті коду
\begin{verbatim}
template <class T> void f(T a1, T a2){}

void g(){
    int i=5;
    f(unique_ptr<int>(new int(++i)), unique_ptr<int>(new int(++i)));
\end{verbatim}

( 1 ) є гонка даних%bad

( 2 ) можливий виток ресурсів%bad

( 3 ) є стан гонки%bad

( 4 ) є невизначена поведінка%good
%enditem
}
%end
\end {large}

\vspace{4mm}
%end
\newpage
%begin
%begin
\begin {large}

\textbf{22.11.2024 Контрольна робота, варіант  \No { 44 }}\par


{
%beginitem
1. По яких днях тижня відбувались лекції з предмету?

%enditem


\vspace{5mm}
%beginitem
2. Як зовуть (ім'я, по батькові) викладача, що вів практичні заняття?

%enditem


\vspace{5mm}
%beginitem
3. Нехай int var=13; Один потік збільшив змінну var на 3, а інший збільшив її в три рази.

гонка даних (data race), стан гонки (race condition)

Виберіть усі правильні твердження (якщо такі є).\par
\vspace{2mm}
( 1 ) Тут гарантовано є стан гонки, але гонки даних може не бути.%good

( 2 ) Тут гарантовано нема гонки даних, але може бути стан гонки.%bad

( 3 ) Тут гарантовано є гонка даних.%bad

( 4 ) Тут гарантовано нема ані гонки даних, ані стану гонки.%bad
%enditem


\vspace{5mm}
%beginitem
4. Виберіть усі правильні твердження (якщо такі є).\par
\vspace{2mm}

( 1 ) Не можна блокувати заблокований м’ютекс mutex.%bad

( 2 ) Потік в активному очікуванні споживає ресурс процесора.%good

( 3 ) Об’єкт future можна скопіювати та роздати копії різним потокам, щоб вони змогли побачити результат.%bad

( 4 ) У стані гонки програма відчайдушно намагається захопити ресурси процесора.%bad
%enditem


\vspace{5mm}
%beginitem
5. Виберіть усі правильні твердження (якщо такі є).\par
\vspace{2mm}

( 1 ) Взаємоблокування потоків може виникнути тільки за використання м’ютексів.%bad

( 2 ) Мова С++20 має стандартну бібліотеку конкурентних структур даних.%bad

( 3 ) Якщо весь код структури даних утворює одну критичну секцію, то така структура даних може використовуватись в конкурентній програмі, але по-справжньо\-му конкурентною не є.%good

( 4 ) Стандартні контейнери stl безпечно використовувати з кількох потоків.%bad
\newpage

%enditem


\vspace{5mm}
%beginitem
6. Виберіть усі правильні твердження (якщо такі є).\par
\vspace{2mm}

( 1 ) Автоматичні змінні не можуть розділятися між потоками.%bad

( 2 ) Динамічні змінні не можуть розділятися між потоками.%bad

( 3 ) Потік мови С++ не може отримати доступ до автоматичних змінних іншого потоку.%bad

( 4 ) Кожен потік програми може мати доступ до динамічних даних програми.%good
%enditem


\vspace{5mm}
%beginitem
7. У цьому фрагменті коду
\begin{verbatim}
template <class T> void f(T a1, T a2){}

void g(){
    int i=5;
    f(unique_ptr<int>(new int(++i)), unique_ptr<int>(new int(++i)));
\end{verbatim}

( 1 ) можливий виток ресурсів%bad

( 2 ) є невизначена поведінка%good

( 3 ) є гонка даних%bad

( 4 ) є стан гонки%bad
%enditem
}
%end
\end {large}

\vspace{4mm}
%end
\newpage
%begin
%begin
\begin {large}

\textbf{22.11.2024 Контрольна робота, варіант  \No { 45 }}\par


{
%beginitem
1. По яких днях тижня відбувались лекції з предмету?

%enditem


\vspace{5mm}
%beginitem
2. Як зовуть (ім'я, по батькові) викладача, що вів практичні заняття?

%enditem


\vspace{5mm}
%beginitem
3. Нехай int var=13; Один потік збільшив змінну var на 3, а інший збільшив її в три рази.

гонка даних (data race), стан гонки (race condition)

Виберіть усі правильні твердження (якщо такі є).\par
\vspace{2mm}
( 1 ) Тут може не бути гонки даних.%good

( 2 ) Тут може бути тільки стан гонки.%bad

( 3 ) Тут гарантовано нема стану гонки, але може бути гонка даних.%bad

( 4 ) Тут може бути тільки гонка даних.%bad
%enditem


\vspace{5mm}
%beginitem
4. Виберіть усі правильні твердження (якщо такі є).\par
\vspace{2mm}

( 1 ) Досвічений програміст за однопотоковою програмою (текстом мовою С++) завжди може вказати послідовність виконуваних обчислень на конкретних вхідних даних.%bad

( 2 ) Блокувати заблокований м’ютекс  \code{mutex} іноді можна.%good

( 3 ) Значення, що є результатом виконання функції, запущеної через async, зберігається безпосередньо в об’єкті проміза (promise).%bad

( 4 ) Довільний потік може розблокувати заблокований м’ютекс.%bad
%enditem


\vspace{5mm}
%beginitem
5. Виберіть усі правильні твердження (якщо такі є).\par
\vspace{2mm}

( 1 ) Структура даних з блокуваннями може бути такою, що використовує атомарні змінні.%good

( 2 ) Структура даних без блокувань та вільна від блокувань структура даних --- це одне те саме.%bad

( 3 ) Операція над вільною від блокувань структурою даних завжди виконається за обмежений час.%bad

( 4 ) Реалізації різних операцій над структурою даних повинні використовувати різні м’ютекси.%bad
\newpage

%enditem


\vspace{5mm}
%beginitem
6. Виберіть усі правильні твердження (якщо такі є).\par
\vspace{2mm}

( 1 ) Автоматичні змінні не можуть розділятися між потоками.%bad

( 2 ) Потік має доступ тільки до своїх власних даних та до глобальних статичних даних.%bad

( 3 ) Кожен потік програми може мати доступ до динамічних даних програми.%good

( 4 ) Динамічні змінні не можуть розділятися між потоками.%bad
%enditem


\vspace{5mm}
%beginitem
7. У цьому фрагменті коду
\begin{verbatim}
template <class T> void f(T a){}

void g(){
    f(unique_ptr<int>(new int(3)))+(unique_ptr<int>(new int(35)));
\end{verbatim}

( 1 ) є невизначена поведінка%bad

( 2 ) можливий виток ресурсів%good

( 3 ) є гонка даних%bad

( 4 ) є стан гонки%bad
%enditem
}
%end
\end {large}

\vspace{4mm}
%end
\newpage
%begin
%begin
\begin {large}

\textbf{22.11.2024 Контрольна робота, варіант  \No { 46 }}\par


{
%beginitem
1. По яких днях тижня відбувались лекції з предмету?

%enditem


\vspace{5mm}
%beginitem
2. Як зовуть (ім'я, по батькові) викладача, що вів практичні заняття?

%enditem


\vspace{5mm}
%beginitem
3. Нехай int var=13; Один потік збільшив змінну var на 3, а інший збільшив її в три рази.

гонка даних (data race), стан гонки (race condition)

Виберіть усі правильні твердження (якщо такі є).\par
\vspace{2mm}
( 1 ) Тут гарантовано нема ані гонки даних, ані стану гонки.%bad

( 2 ) Тут гарантовано є стан гонки.%good

( 3 ) Тут гарантовано є і стан гонки, і гонка даних.%bad

( 4 ) Тут може не бути ані гонки даних, ані стану гонки.%bad
%enditem


\vspace{5mm}
%beginitem
4. Виберіть усі правильні твердження (якщо такі є).\par
\vspace{2mm}

( 1 ) Чи запустить функція  async виконання в новому потоці можна задавати політикою.%good

( 2 ) Клас thread має засоби, щоб повернути результат виконання запущеного потоку.%bad

( 3 ) Тільки від компілятора залежить, чи запустить функція async виконання в новому потоці.%bad

( 4 ) Об’єкт promise можна скопіювати та роздати копії різним потокам, щоб кожен з них міг передати результат своєї роботи.%bad
%enditem


\vspace{5mm}
%beginitem
5. Виберіть усі правильні твердження (якщо такі є).\par
\vspace{2mm}

( 1 ) За використання структурою даних рівно одного м’ютексу, взаємоблокування на її операціях неможливі.%good

( 2 ) Структура даних є структурою даних без блокувань, якщо її реалізація не використовує м’ютекси.%bad

( 3 ) Для уникнення гонки даних структура даних може використовувати не більше одного м’ютекса.%bad

( 4 ) Одночасно використовувати атомарних змінних більше, ніж регістрів процесора, у програмі не можна.%bad
\newpage

%enditem


\vspace{5mm}
%beginitem
6. Виберіть усі правильні твердження (якщо такі є).\par
\vspace{2mm}

( 1 ) Потік мови С++ не може отримати доступ до автоматичних змінних іншого потоку.%bad

( 2 ) Потік може мати доступ до автоматичних змінних іншого потоку і це безпечно.%bad

( 3 ) Потік мови С++ не може отримати доступ до даних іншого потоку.%bad

( 4 ) Потік може мати доступ до автоматичних змінних іншого потоку, але це небезпечно.%good
%enditem


\vspace{5mm}
%beginitem
7. У цьому фрагменті коду
\begin{verbatim}
template <class T> void f(T a){}

void g(){
    int i=5;
    f(unique_ptr<int>(new int(i)))+(unique_ptr<int>(new int(++i)));
\end{verbatim}
( 1 ) є гонка даних%bad

( 2 ) можливий виток ресурсів%good

( 3 ) є стан гонки%bad

( 4 ) є невизначена поведінка%good
%enditem
}
%end
\end {large}

\vspace{4mm}
%end
\newpage
%begin
%begin
\begin {large}

\textbf{22.11.2024 Контрольна робота, варіант  \No { 47 }}\par


{
%beginitem
1. По яких днях тижня відбувались лекції з предмету?

%enditem


\vspace{5mm}
%beginitem
2. Як зовуть (ім'я, по батькові) викладача, що вів практичні заняття?

%enditem


\vspace{5mm}
%beginitem
3. Нехай int var=13; Один потік збільшив змінну var на 3, а інший збільшив її в три рази.

гонка даних (data race), стан гонки (race condition)

Виберіть усі правильні твердження (якщо такі є).\par
\vspace{2mm}
( 1 ) Тут гарантовано є стан гонки, але гонки даних може не бути.%good

( 2 ) Тут може бути тільки гонка даних.%bad

( 3 ) Тут гарантовано є гонка даних.%bad

( 4 ) Тут гарантовано нема стану гонки, але може бути гонка даних.%bad
%enditem


\vspace{5mm}
%beginitem
4. Виберіть усі правильні твердження (якщо такі є).\par
\vspace{2mm}

( 1 ) Блокувати заблокований м’ютекс  \code{mutex} іноді можна.%good

( 2 ) Секція є критичною, якщо під час її виконання не можна генерувати винятки.%bad

( 3 ) Розблоковувати незаблокований м’ютекс можна.%bad

( 4 ) Функція async запускає виконання в новому потоці.%bad
%enditem


\vspace{5mm}
%beginitem
5. Виберіть усі правильні твердження (якщо такі є).\par
\vspace{2mm}

( 1 ) На операціях структури даних без блокувань взаємоблокувань потоків виникнути не може.%bad

( 2 ) Структура даних без блокувань гарантує завершення виконання довільної операції над цією структурою даних за обмежений час.%bad

( 3 ) Структура даних без блокувань зобов’язана використовувати атомарні змінні.%bad

( 4 ) Структура даних з блокуваннями може бути такою, що використовує атомарні змінні.%good
\newpage

%enditem


\vspace{5mm}
%beginitem
6. Виберіть усі правильні твердження (якщо такі є).\par
\vspace{2mm}

( 1 ) Автоматичні змінні не можуть розділятися між потоками.%bad

( 2 ) Потік мови С++ не може отримати доступ до даних іншого потоку.%bad

( 3 ) Потік може мати доступ до автоматичних змінних іншого потоку, але це небезпечно.%good

( 4 ) Потік мови С++ не може отримати доступ до автоматичних змінних іншого потоку.%bad
%enditem


\vspace{5mm}
%beginitem
7. У цьому фрагменті коду
\begin{verbatim}
template <class T> void f(T a1, T a2){}

void g(){
    f(unique_ptr<int>(new int(3)), unique_ptr<int>(new int(35)));
\end{verbatim}

( 1 ) є гонка даних%bad

( 2 ) можливий виток ресурсів%bad

( 3 ) є невизначена поведінка%bad

( 4 ) є стан гонки%bad
%enditem
}
%end
\end {large}

\vspace{4mm}
%end
\newpage
%begin
%begin
\begin {large}

\textbf{22.11.2024 Контрольна робота, варіант  \No { 48 }}\par


{
%beginitem
1. По яких днях тижня відбувались лекції з предмету?

%enditem


\vspace{5mm}
%beginitem
2. Як зовуть (ім'я, по батькові) викладача, що вів практичні заняття?

%enditem


\vspace{5mm}
%beginitem
3. Нехай int var=13; Один потік збільшив змінну var на 3, а інший збільшив її в три рази.

гонка даних (data race), стан гонки (race condition)

Виберіть усі правильні твердження (якщо такі є).\par
\vspace{2mm}
( 1 ) Тут може не бути гонки даних.%good

( 2 ) Тут може не бути стану гонки.%bad

( 3 ) Тут гарантовано нема гонки даних, але може бути стан гонки.%bad

( 4 ) Тут може бути тільки стан гонки.%bad
%enditem


\vspace{5mm}
%beginitem
4. Виберіть усі правильні твердження (якщо такі є).\par
\vspace{2mm}

( 1 ) Коли програма спить, то вона не виконується.%bad

( 2 ) У програмі може бути не більше однієї критичної секції.%bad

( 3 ) Виклик функції \code{this\_thread::sleep\_for(5s)} змушує потік, з якого його здійснили, заснути якнайменш на 5 секунд.%good

( 4 ) Щоб розблокувати заблокований м’ютекс \code{recursive\_mutex}, його завжди достатньо розблокувати один раз.%bad
%enditem


\vspace{5mm}
%beginitem
5. Виберіть усі правильні твердження (якщо такі є).\par
\vspace{2mm}

( 1 ) Кожна вільна від очікувань структура даних є структурою даних, вільною від блокувань.%good

( 2 ) Мова С++20 має стандартну бібліотеку вільних від блокувань структур даних.%bad

( 3 ) Вільних від очікувань структур даних не існує.%bad

( 4 ) За використання атомарних змінних структура даних є структурою даних без блокувань.%bad
\newpage

%enditem


\vspace{5mm}
%beginitem
6. Виберіть усі правильні твердження (якщо такі є).\par
\vspace{2mm}

( 1 ) Динамічні змінні не можуть розділятися між потоками.%bad

( 2 ) Кожен потік програми може мати доступ до динамічних даних програми.%good

( 3 ) Потік може мати доступ до автоматичних змінних іншого потоку і це безпечно.%bad

( 4 ) Потік має доступ тільки до своїх власних даних та до глобальних статичних даних.%bad
%enditem


\vspace{5mm}
%beginitem
7. У цьому фрагменті коду
\begin{verbatim}
template <class T> void f(T a){}

void g(){
    int i=5;
    f(unique_ptr<int>(new int(i)))+(unique_ptr<int>(new int(35)));
\end{verbatim}
( 1 ) є гонка даних%bad

( 2 ) можливий виток ресурсів%good

( 3 ) є невизначена поведінка%bad

( 4 ) є стан гонки%bad
%enditem
}
%end
\end {large}

\vspace{4mm}
%end
\newpage
%begin
%begin
\begin {large}

\textbf{22.11.2024 Контрольна робота, варіант  \No { 49 }}\par


{
%beginitem
1. По яких днях тижня відбувались лекції з предмету?

%enditem


\vspace{5mm}
%beginitem
2. Як зовуть (ім'я, по батькові) викладача, що вів практичні заняття?

%enditem


\vspace{5mm}
%beginitem
3. Нехай int var=13; Один потік збільшив змінну var на 3, а інший збільшив її в три рази.

гонка даних (data race), стан гонки (race condition)

Виберіть усі правильні твердження (якщо такі є).\par
\vspace{2mm}
( 1 ) Тут гарантовано нема гонки даних, але може бути стан гонки.%bad

( 2 ) Тут може не бути ані гонки даних, ані стану гонки.%bad

( 3 ) Тут може не бути гонки даних.%good

( 4 ) Тут може бути тільки стан гонки.%bad
%enditem


\vspace{5mm}
%beginitem
4. Виберіть усі правильні твердження (якщо такі є).\par
\vspace{2mm}

( 1 ) Потік в активному очікуванні споживає ресурс процесора.%good

( 2 ) Об’єкт promise можна скопіювати та роздати копії різним потокам, щоб кожен з них міг передати результат своєї роботи.%bad

( 3 ) Тільки від компілятора залежить, чи запустить функція async виконання в новому потоці.%bad

( 4 ) М’ютекс має бути записаний на початку критичної секції, яку він забезпечує.%bad
%enditem


\vspace{5mm}
%beginitem
5. Виберіть усі правильні твердження (якщо такі є).\par
\vspace{2mm}

( 1 ) За використання атомарних змінних структура даних є структурою даних без блокувань.%bad

( 2 ) Взаємоблокування потоків може виникнути тільки за використання м’ютексів.%bad

( 3 ) На операціях структури даних без блокувань взаємоблокувань потоків виникнути не може.%bad

( 4 ) Структура даних без блокувань може не бути вільною від блокувань.%good
\newpage

%enditem


\vspace{5mm}
%beginitem
6. Виберіть усі правильні твердження (якщо такі є).\par
\vspace{2mm}

( 1 ) Потік мови С++ не може отримати доступ до автоматичних змінних іншого потоку.%bad

( 2 ) Потік може мати доступ до автоматичних змінних іншого потоку і це безпечно.%bad

( 3 ) Динамічні змінні не можуть розділятися між потоками.%bad

( 4 ) Потік може мати доступ до автоматичних змінних іншого потоку, але це небезпечно.%good
%enditem


\vspace{5mm}
%beginitem
7. У цьому фрагменті коду
\begin{verbatim}
template <class T> void f(T a1, T a2){}

void g(){
    int i=5;
    f(unique_ptr<int>(new int(i++)), unique_ptr<int>(new int(++i)));
\end{verbatim}

( 1 ) є гонка даних%bad

( 2 ) є невизначена поведінка%good

( 3 ) можливий виток ресурсів%bad

( 4 ) є стан гонки%bad
%enditem
}
%end
\end {large}

\vspace{4mm}
%end
\newpage
%begin
%begin
\begin {large}

\textbf{22.11.2024 Контрольна робота, варіант  \No { 50 }}\par


{
%beginitem
1. По яких днях тижня відбувались лекції з предмету?

%enditem


\vspace{5mm}
%beginitem
2. Як зовуть (ім'я, по батькові) викладача, що вів практичні заняття?

%enditem


\vspace{5mm}
%beginitem
3. Нехай int var=13; Один потік збільшив змінну var на 3, а інший збільшив її в три рази.

гонка даних (data race), стан гонки (race condition)

Виберіть усі правильні твердження (якщо такі є).\par
\vspace{2mm}
( 1 ) Тут гарантовано нема стану гонки, але може бути гонка даних.%bad

( 2 ) Тут гарантовано є стан гонки.%good

( 3 ) Тут гарантовано нема ані гонки даних, ані стану гонки.%bad

( 4 ) Тут гарантовано є і стан гонки, і гонка даних.%bad
%enditem


\vspace{5mm}
%beginitem
4. Виберіть усі правильні твердження (якщо такі є).\par
\vspace{2mm}

( 1 ) Тільки той потік, який обнулив значення семафора, може його збільшити.%bad

( 2 ) Розблоковувати незаблокований м’ютекс можна.%bad

( 3 ) Функція async запускає виконання в новому потоці.%bad

( 4 ) Критичну секцію можна реалізувати за допомогою м’ютексу.%good
%enditem


\vspace{5mm}
%beginitem
5. Виберіть усі правильні твердження (якщо такі є).\par
\vspace{2mm}

( 1 ) Реалізації різних операцій над структурою даних повинні використовувати різні м’ютекси.%bad

( 2 ) Вільних від очікувань структур даних не існує.%bad

( 3 ) Реалізація структури даних з блокуваннями може використовувати атомарні змінні.%good

( 4 ) Операція над структурою даних без блокувань завжди виконається за обмежений час.%bad
\newpage

%enditem


\vspace{5mm}
%beginitem
6. Виберіть усі правильні твердження (якщо такі є).\par
\vspace{2mm}

( 1 ) Потік мови С++ не може отримати доступ до даних іншого потоку.%bad

( 2 ) Кожен потік програми може мати доступ до динамічних даних програми.%good

( 3 ) Потік має доступ тільки до своїх власних даних та до глобальних статичних даних.%bad

( 4 ) Автоматичні змінні не можуть розділятися між потоками.%bad
%enditem


\vspace{5mm}
%beginitem
7. У цьому фрагменті коду
\begin{verbatim}
template <class T> void f(T a1, T a2){}

void g(){
    int i=5;
    f(unique_ptr<int>(new int(++i)), unique_ptr<int>(new int(++i)));
\end{verbatim}

( 1 ) можливий виток ресурсів%bad

( 2 ) є гонка даних%bad

( 3 ) є стан гонки%bad

( 4 ) є невизначена поведінка%good
%enditem
}
%end
\end {large}

\vspace{4mm}
%end
\newpage
%begin
%begin
\begin {large}

\textbf{22.11.2024 Контрольна робота, варіант  \No { 51 }}\par


{
%beginitem
1. По яких днях тижня відбувались лекції з предмету?

%enditem


\vspace{5mm}
%beginitem
2. Як зовуть (ім'я, по батькові) викладача, що вів практичні заняття?

%enditem


\vspace{5mm}
%beginitem
3. Нехай int var=13; Один потік збільшив змінну var на 3, а інший збільшив її в три рази.

гонка даних (data race), стан гонки (race condition)

Виберіть усі правильні твердження (якщо такі є).\par
\vspace{2mm}
( 1 ) Тут гарантовано є стан гонки, але гонки даних може не бути.%good

( 2 ) Тут може бути тільки гонка даних.%bad

( 3 ) Тут гарантовано є гонка даних.%bad

( 4 ) Тут може не бути стану гонки.%bad
%enditem


\vspace{5mm}
%beginitem
4. Виберіть усі правильні твердження (якщо такі є).\par
\vspace{2mm}

( 1 ) Значення, що є результатом виконання функції, запущеної через async, зберігається безпосередньо в об’єкті проміза (promise).%bad

( 2 ) Порядок обчислень у потоці для одних тих самих вхідних даних може залежати від використовуваного компілятора.%good

( 3 ) Текст однопотокової програми мовою С++ однозначно визначає послідовність виконуваних обчислень.%bad

( 4 ) Досвічений програміст за однопотоковою програмою (текстом мовою С++) завжди може вказати послідовність виконуваних обчислень на конкретних вхідних даних.%bad
%enditem


\vspace{5mm}
%beginitem
5. Виберіть усі правильні твердження (якщо такі є).\par
\vspace{2mm}

( 1 ) Структура даних може використовувати м’ютекси та атомарні змінні одночасно.%good

( 2 ) Структура даних без блокувань та вільна від блокувань структура даних --- це одне те саме.%bad

( 3 ) Структура даних без блокувань гарантує завершення виконання довільної операції над цією структурою даних за обмежений час.%bad

( 4 ) Структура даних без блокувань зобов’язана використовувати атомарні змінні.%bad
\newpage

%enditem


\vspace{5mm}
%beginitem
6. Виберіть усі правильні твердження (якщо такі є).\par
\vspace{2mm}

( 1 ) Динамічні змінні не можуть розділятися між потоками.%bad

( 2 ) Потік може мати доступ до автоматичних змінних іншого потоку і це безпечно.%bad

( 3 ) Кожен потік програми може мати доступ до динамічних даних програми.%good

( 4 ) Автоматичні змінні не можуть розділятися між потоками.%bad
%enditem


\vspace{5mm}
%beginitem
7. У цьому фрагменті коду
\begin{verbatim}
template <class T> void f(T a1, T a2){}

void g(){
    f(unique_ptr<int>(new int(3)), unique_ptr<int>(new int(35)));
\end{verbatim}

( 1 ) є стан гонки%bad

( 2 ) можливий виток ресурсів%bad

( 3 ) є невизначена поведінка%bad

( 4 ) є гонка даних%bad
%enditem
}
%end
\end {large}

\vspace{4mm}
%end
\newpage
%begin
%begin
\begin {large}

\textbf{22.11.2024 Контрольна робота, варіант  \No { 52 }}\par


{
%beginitem
1. По яких днях тижня відбувались лекції з предмету?

%enditem


\vspace{5mm}
%beginitem
2. Як зовуть (ім'я, по батькові) викладача, що вів практичні заняття?

%enditem


\vspace{5mm}
%beginitem
3. Нехай int var=13; Один потік збільшив змінну var на 3, а інший збільшив її в три рази.

гонка даних (data race), стан гонки (race condition)

Виберіть усі правильні твердження (якщо такі є).\par
\vspace{2mm}
( 1 ) Тут гарантовано нема стану гонки, але може бути гонка даних.%bad

( 2 ) Тут може не бути гонки даних.%good

( 3 ) Тут може не бути ані гонки даних, ані стану гонки.%bad

( 4 ) Тут гарантовано є і стан гонки, і гонка даних.%bad
%enditem


\vspace{5mm}
%beginitem
4. Виберіть усі правильні твердження (якщо такі є).\par
\vspace{2mm}

( 1 ) Об’єкт future можна скопіювати та роздати копії різним потокам, щоб вони змогли побачити результат.%bad

( 2 ) Якщо два обчислення програми на С++ тільки читають одну ту саму ділянку пам'яті, то на них гонка даних  виникнути не може.%good

( 3 ) Секція є критичною, якщо під час її виконання не можна генерувати винятки.%bad

( 4 ) Не можна блокувати заблокований м’ютекс mutex.%bad
%enditem


\vspace{5mm}
%beginitem
5. Виберіть усі правильні твердження (якщо такі є).\par
\vspace{2mm}

( 1 ) Структура даних без блокувань гарантує, що жоден потік не буде очікувати доступу до виконання операцій над цією структурою даних.%bad

( 2 ) Мова С++20 має стандартну бібліотеку конкурентних структур даних.%bad

( 3 ) Якщо весь код структури даних утворює одну критичну секцію, то така структура даних може використовуватись в конкурентній програмі, але по-справжньо\-му конкурентною не є.%good

( 4 ) Одночасно використовувати атомарних змінних більше, ніж регістрів процесора, у програмі не можна.%bad
\newpage

%enditem


\vspace{5mm}
%beginitem
6. Виберіть усі правильні твердження (якщо такі є).\par
\vspace{2mm}

( 1 ) Потік мови С++ не може отримати доступ до автоматичних змінних іншого потоку.%bad

( 2 ) Потік має доступ тільки до своїх власних даних та до глобальних статичних даних.%bad

( 3 ) Потік може мати доступ до автоматичних змінних іншого потоку, але це небезпечно.%good

( 4 ) Потік мови С++ не може отримати доступ до даних іншого потоку.%bad
%enditem


\vspace{5mm}
%beginitem
7. У цьому фрагменті коду
\begin{verbatim}
template <class T> void f(T a1, T a2){}

void g(){
    int i=5;
    f(unique_ptr<int>(new int(i++)), unique_ptr<int>(new int(++i)));
\end{verbatim}

( 1 ) є стан гонки%bad

( 2 ) є невизначена поведінка%good

( 3 ) можливий виток ресурсів%bad

( 4 ) є гонка даних%bad
%enditem
}
%end
\end {large}

\vspace{4mm}
%end
\newpage
%begin
%begin
\begin {large}

\textbf{22.11.2024 Контрольна робота, варіант  \No { 53 }}\par


{
%beginitem
1. По яких днях тижня відбувались лекції з предмету?

%enditem


\vspace{5mm}
%beginitem
2. Як зовуть (ім'я, по батькові) викладача, що вів практичні заняття?

%enditem


\vspace{5mm}
%beginitem
3. Нехай int var=13; Один потік збільшив змінну var на 3, а інший збільшив її в три рази.

гонка даних (data race), стан гонки (race condition)

Виберіть усі правильні твердження (якщо такі є).\par
\vspace{2mm}
( 1 ) Тут може бути тільки гонка даних.%bad

( 2 ) Тут гарантовано є гонка даних.%bad

( 3 ) Тут гарантовано є стан гонки, але гонки даних може не бути.%good

( 4 ) Тут гарантовано нема ані гонки даних, ані стану гонки.%bad
%enditem


\vspace{5mm}
%beginitem
4. Виберіть усі правильні твердження (якщо такі є).\par
\vspace{2mm}

( 1 ) За різних запусків програми на одних тих самих вхідних даних може виникати хронологічно різна послідовність обчислень.%good

( 2 ) У стані гонки програма відчайдушно намагається захопити ресурси процесора.%bad

( 3 ) Потік, запущений конструктором класу jthread, має обов’язково зупинитися, якщо пов’язаний об’єкт надішле команду зупинитися.%bad

( 4 ) Значення, що є результатом виконання функції, запущеної через async, зберігається безпосередньо в об’єкті ф’ючерса (future).%bad
%enditem


\vspace{5mm}
%beginitem
5. Виберіть усі правильні твердження (якщо такі є).\par
\vspace{2mm}

( 1 ) Структура даних є структурою даних без блокувань, якщо її реалізація не використовує м’ютекси.%bad

( 2 ) Реалізація структури даних без блокувань не може використовувати м’ютекси.%good

( 3 ) Для уникнення гонки даних структура даних може використовувати не більше одного м’ютекса.%bad

( 4 ) Стандартні контейнери stl безпечно використовувати з кількох потоків.%bad
\newpage

%enditem


\vspace{5mm}
%beginitem
6. Виберіть усі правильні твердження (якщо такі є).\par
\vspace{2mm}

( 1 ) Кожен потік програми може мати доступ до динамічних даних програми.%good

( 2 ) Потік мови С++ не може отримати доступ до даних іншого потоку.%bad

( 3 ) Потік мови С++ не може отримати доступ до автоматичних змінних іншого потоку.%bad

( 4 ) Автоматичні змінні не можуть розділятися між потоками.%bad
%enditem


\vspace{5mm}
%beginitem
7. У цьому фрагменті коду
\begin{verbatim}
template <class T> void f(T a){}

void g(){
    int i=5;
    f(unique_ptr<int>(new int(i)))+(unique_ptr<int>(new int(++i)));
\end{verbatim}
( 1 ) можливий виток ресурсів%good

( 2 ) є стан гонки%bad

( 3 ) є невизначена поведінка%good

( 4 ) є гонка даних%bad
%enditem
}
%end
\end {large}

\vspace{4mm}
%end
\newpage
%begin
%begin
\begin {large}

\textbf{22.11.2024 Контрольна робота, варіант  \No { 54 }}\par


{
%beginitem
1. По яких днях тижня відбувались лекції з предмету?

%enditem


\vspace{5mm}
%beginitem
2. Як зовуть (ім'я, по батькові) викладача, що вів практичні заняття?

%enditem


\vspace{5mm}
%beginitem
3. Нехай int var=13; Один потік збільшив змінну var на 3, а інший збільшив її в три рази.

гонка даних (data race), стан гонки (race condition)

Виберіть усі правильні твердження (якщо такі є).\par
\vspace{2mm}
( 1 ) Тут гарантовано нема гонки даних, але може бути стан гонки.%bad

( 2 ) Тут може бути тільки стан гонки.%bad

( 3 ) Тут може не бути стану гонки.%bad

( 4 ) Тут гарантовано є стан гонки.%good
%enditem


\vspace{5mm}
%beginitem
4. Виберіть усі правильні твердження (якщо такі є).\par
\vspace{2mm}

( 1 ) Коли потік спить, він не споживає ресурс процесора.%good

( 2 ) Довільний потік може розблокувати заблокований м’ютекс.%bad

( 3 ) Виклик \code{this\_thread::sleep\_for(5s)} змушує потік проспати рівно 5 секунд.%bad

( 4 ) Об’єкти класів promise та future можна використовувати тільки в багатопотокових програмах.%bad
%enditem


\vspace{5mm}
%beginitem
5. Виберіть усі правильні твердження (якщо такі є).\par
\vspace{2mm}

( 1 ) Мова С++20 має стандартну бібліотеку вільних від блокувань структур даних.%bad

( 2 ) Операція над вільною від блокувань структурою даних може виконуватись скільки завгодно довго.%good

( 3 ) Операція над вільною від блокувань структурою даних завжди виконається за обмежений час.%bad

( 4 ) Вільна від блокувань структура даних гарантує, що жоден потік не буде очікувати доступу до виконання операцій над цією структурою даних.%bad
\newpage

%enditem


\vspace{5mm}
%beginitem
6. Виберіть усі правильні твердження (якщо такі є).\par
\vspace{2mm}

( 1 ) Потік може мати доступ до автоматичних змінних іншого потоку і це безпечно.%bad

( 2 ) Динамічні змінні не можуть розділятися між потоками.%bad

( 3 ) Потік має доступ тільки до своїх власних даних та до глобальних статичних даних.%bad

( 4 ) Потік може мати доступ до автоматичних змінних іншого потоку, але це небезпечно.%good
%enditem


\vspace{5mm}
%beginitem
7. У цьому фрагменті коду
\begin{verbatim}
template <class T> void f(T a){}

void g(){
    int i=5;
    f(unique_ptr<int>(new int(i)))+(unique_ptr<int>(new int(35)));
\end{verbatim}
( 1 ) є невизначена поведінка%bad

( 2 ) є стан гонки%bad

( 3 ) є гонка даних%bad

( 4 ) можливий виток ресурсів%good
%enditem
}
%end
\end {large}

\vspace{4mm}
%end
\newpage
%begin
%begin
\begin {large}

\textbf{22.11.2024 Контрольна робота, варіант  \No { 55 }}\par


{
%beginitem
1. По яких днях тижня відбувались лекції з предмету?

%enditem


\vspace{5mm}
%beginitem
2. Як зовуть (ім'я, по батькові) викладача, що вів практичні заняття?

%enditem


\vspace{5mm}
%beginitem
3. Нехай int var=13; Один потік збільшив змінну var на 3, а інший збільшив її в три рази.

гонка даних (data race), стан гонки (race condition)

Виберіть усі правильні твердження (якщо такі є).\par
\vspace{2mm}
( 1 ) Тут гарантовано нема стану гонки, але може бути гонка даних.%bad

( 2 ) Тут гарантовано є стан гонки, але гонки даних може не бути.%good

( 3 ) Тут може бути тільки гонка даних.%bad

( 4 ) Тут може бути тільки стан гонки.%bad
%enditem


\vspace{5mm}
%beginitem
4. Виберіть усі правильні твердження (якщо такі є).\par
\vspace{2mm}

( 1 ) Щоб розблокувати заблокований м’ютекс \code{recursive\_mutex}, його завжди достатньо розблокувати один раз.%bad

( 2 ) Клас thread має засоби, щоб повернути результат виконання запущеного потоку.%bad

( 3 ) Коли програма спить, то вона не виконується.%bad

( 4 ) Розблокувати м’ютекс може тільки той потік, що його заблокував.%good
%enditem


\vspace{5mm}
%beginitem
5. Виберіть усі правильні твердження (якщо такі є).\par
\vspace{2mm}

( 1 ) Структура даних без блокувань зобов’язана використовувати атомарні змінні.%bad

( 2 ) Якщо весь код структури даних утворює одну критичну секцію, то така структура даних може використовуватись в конкурентній програмі, але по-справжньо\-му конкурентною не є.%good

( 3 ) Структура даних без блокувань та вільна від блокувань структура даних --- це одне те саме.%bad

( 4 ) Структура даних без блокувань гарантує завершення виконання довільної операції над цією структурою даних за обмежений час.%bad
\newpage

%enditem


\vspace{5mm}
%beginitem
6. Виберіть усі правильні твердження (якщо такі є).\par
\vspace{2mm}

( 1 ) Потік може мати доступ до автоматичних змінних іншого потоку і це безпечно.%bad

( 2 ) Автоматичні змінні не можуть розділятися між потоками.%bad

( 3 ) Потік мови С++ не може отримати доступ до автоматичних змінних іншого потоку.%bad

( 4 ) Потік може мати доступ до автоматичних змінних іншого потоку, але це небезпечно.%good
%enditem


\vspace{5mm}
%beginitem
7. У цьому фрагменті коду
\begin{verbatim}
template <class T> void f(T a){}

void g(){
    f(unique_ptr<int>(new int(3)))+(unique_ptr<int>(new int(35)));
\end{verbatim}

( 1 ) можливий виток ресурсів%good

( 2 ) є невизначена поведінка%bad

( 3 ) є стан гонки%bad

( 4 ) є гонка даних%bad
%enditem
}
%end
\end {large}

\vspace{4mm}
%end
\newpage
%begin
%begin
\begin {large}

\textbf{22.11.2024 Контрольна робота, варіант  \No { 56 }}\par


{
%beginitem
1. По яких днях тижня відбувались лекції з предмету?

%enditem


\vspace{5mm}
%beginitem
2. Як зовуть (ім'я, по батькові) викладача, що вів практичні заняття?

%enditem


\vspace{5mm}
%beginitem
3. Нехай int var=13; Один потік збільшив змінну var на 3, а інший збільшив її в три рази.

гонка даних (data race), стан гонки (race condition)

Виберіть усі правильні твердження (якщо такі є).\par
\vspace{2mm}
( 1 ) Тут гарантовано є стан гонки.%good

( 2 ) Тут може не бути стану гонки.%bad

( 3 ) Тут гарантовано нема ані гонки даних, ані стану гонки.%bad

( 4 ) Тут може не бути ані гонки даних, ані стану гонки.%bad
%enditem


\vspace{5mm}
%beginitem
4. Виберіть усі правильні твердження (якщо такі є).\par
\vspace{2mm}

( 1 ) Функція \code{sleep\_for} може змусити інший потік заснути.%bad

( 2 ) У програмі може бути не більше однієї критичної секції.%bad

( 3 ) Критична секція це неперервна ділянка коду програми.%bad

( 4 ) Клас thread не має засобів, щоб повернути результат виконання запущеного потоку.%good
%enditem


\vspace{5mm}
%beginitem
5. Виберіть усі правильні твердження (якщо такі є).\par
\vspace{2mm}

( 1 ) Взаємоблокування потоків може виникнути тільки за використання м’ютексів.%bad

( 2 ) Операція над структурою даних без блокувань завжди виконається за обмежений час.%bad

( 3 ) Структура даних може використовувати м’ютекси та атомарні змінні одночасно.%good

( 4 ) Для уникнення гонки даних структура даних може використовувати не більше одного м’ютекса.%bad
\newpage

%enditem


\vspace{5mm}
%beginitem
6. Виберіть усі правильні твердження (якщо такі є).\par
\vspace{2mm}

( 1 ) Потік має доступ тільки до своїх власних даних та до глобальних статичних даних.%bad

( 2 ) Кожен потік програми може мати доступ до динамічних даних програми.%good

( 3 ) Динамічні змінні не можуть розділятися між потоками.%bad

( 4 ) Потік мови С++ не може отримати доступ до даних іншого потоку.%bad
%enditem


\vspace{5mm}
%beginitem
7. У цьому фрагменті коду
\begin{verbatim}
template <class T> void f(T a1, T a2){}

void g(){
    int i=5;
    f(unique_ptr<int>(new int(++i)), unique_ptr<int>(new int(++i)));
\end{verbatim}

( 1 ) є гонка даних%bad

( 2 ) є стан гонки%bad

( 3 ) можливий виток ресурсів%bad

( 4 ) є невизначена поведінка%good
%enditem
}
%end
\end {large}

\vspace{4mm}
%end
\newpage
%begin
%begin
\begin {large}

\textbf{22.11.2024 Контрольна робота, варіант  \No { 57 }}\par


{
%beginitem
1. По яких днях тижня відбувались лекції з предмету?

%enditem


\vspace{5mm}
%beginitem
2. Як зовуть (ім'я, по батькові) викладача, що вів практичні заняття?

%enditem


\vspace{5mm}
%beginitem
3. Нехай int var=13; Один потік збільшив змінну var на 3, а інший збільшив її в три рази.

гонка даних (data race), стан гонки (race condition)

Виберіть усі правильні твердження (якщо такі є).\par
\vspace{2mm}
( 1 ) Тут гарантовано є і стан гонки, і гонка даних.%bad

( 2 ) Тут гарантовано є гонка даних.%bad

( 3 ) Тут гарантовано нема гонки даних, але може бути стан гонки.%bad

( 4 ) Тут може не бути гонки даних.%good
%enditem


\vspace{5mm}
%beginitem
4. Виберіть усі правильні твердження (якщо такі є).\par
\vspace{2mm}

( 1 ) Щоб розблокувати заблокований м’ютекс \code{recursive\_mutex}, його завжди достатньо розблокувати один раз.%bad

( 2 ) Тільки від компілятора залежить, чи запустить функція async виконання в новому потоці.%bad

( 3 ) Критична секція може складатися з кількох ділянок коду, які можуть розташовуватись де завгодно (не обов’язково зв’язний фрагмент коду).%good

( 4 ) Потік, запущений конструктором класу jthread, має обов’язково зупинитися, якщо пов’язаний об’єкт надішле команду зупинитися.%bad
%enditem


\vspace{5mm}
%beginitem
5. Виберіть усі правильні твердження (якщо такі є).\par
\vspace{2mm}

( 1 ) Мова С++20 має стандартну бібліотеку конкурентних структур даних.%bad

( 2 ) Операція над вільною від блокувань структурою даних може виконуватись скільки завгодно довго.%good

( 3 ) Операція над вільною від блокувань структурою даних завжди виконається за обмежений час.%bad

( 4 ) Одночасно використовувати атомарних змінних більше, ніж регістрів процесора, у програмі не можна.%bad
\newpage

%enditem


\vspace{5mm}
%beginitem
6. Виберіть усі правильні твердження (якщо такі є).\par
\vspace{2mm}

( 1 ) Динамічні змінні не можуть розділятися між потоками.%bad

( 2 ) Потік має доступ тільки до своїх власних даних та до глобальних статичних даних.%bad

( 3 ) Автоматичні змінні не можуть розділятися між потоками.%bad

( 4 ) Кожен потік програми може мати доступ до динамічних даних програми.%good
%enditem


\vspace{5mm}
%beginitem
7. У цьому фрагменті коду
\begin{verbatim}
template <class T> void f(T a){}

void g(){
    f(unique_ptr<int>(new int(3)))+(unique_ptr<int>(new int(35)));
\end{verbatim}

( 1 ) є гонка даних%bad

( 2 ) можливий виток ресурсів%good

( 3 ) є стан гонки%bad

( 4 ) є невизначена поведінка%bad
%enditem
}
%end
\end {large}

\vspace{4mm}
%end
\newpage
%begin
%begin
\begin {large}

\textbf{22.11.2024 Контрольна робота, варіант  \No { 58 }}\par


{
%beginitem
1. По яких днях тижня відбувались лекції з предмету?

%enditem


\vspace{5mm}
%beginitem
2. Як зовуть (ім'я, по батькові) викладача, що вів практичні заняття?

%enditem


\vspace{5mm}
%beginitem
3. Нехай int var=13; Один потік збільшив змінну var на 3, а інший збільшив її в три рази.

гонка даних (data race), стан гонки (race condition)

Виберіть усі правильні твердження (якщо такі є).\par
\vspace{2mm}
( 1 ) Тут гарантовано є стан гонки, але гонки даних може не бути.%good

( 2 ) Тут може не бути ані гонки даних, ані стану гонки.%bad

( 3 ) Тут гарантовано нема стану гонки, але може бути гонка даних.%bad

( 4 ) Тут гарантовано є гонка даних.%bad
%enditem


\vspace{5mm}
%beginitem
4. Виберіть усі правильні твердження (якщо такі є).\par
\vspace{2mm}

( 1 ) У синтаксично коректній програмі мовою С++ стан гонки (race condition) може виникати.%good

( 2 ) Значення, що є результатом виконання функції, запущеної через async, зберігається безпосередньо в об’єкті проміза (promise).%bad

( 3 ) Критична секція це неперервна ділянка коду програми.%bad

( 4 ) Текст однопотокової програми мовою С++ однозначно визначає послідовність виконуваних обчислень.%bad
%enditem


\vspace{5mm}
%beginitem
5. Виберіть усі правильні твердження (якщо такі є).\par
\vspace{2mm}

( 1 ) За використання атомарних змінних структура даних є структурою даних без блокувань.%bad

( 2 ) Структура даних є структурою даних без блокувань, якщо її реалізація не використовує м’ютекси.%bad

( 3 ) Реалізація структури даних без блокувань не може використовувати м’ютекси.%good

( 4 ) Вільна від блокувань структура даних гарантує, що жоден потік не буде очікувати доступу до виконання операцій над цією структурою даних.%bad
\newpage

%enditem


\vspace{5mm}
%beginitem
6. Виберіть усі правильні твердження (якщо такі є).\par
\vspace{2mm}

( 1 ) Потік мови С++ не може отримати доступ до даних іншого потоку.%bad

( 2 ) Потік може мати доступ до автоматичних змінних іншого потоку і це безпечно.%bad

( 3 ) Потік мови С++ не може отримати доступ до автоматичних змінних іншого потоку.%bad

( 4 ) Потік може мати доступ до автоматичних змінних іншого потоку, але це небезпечно.%good
%enditem


\vspace{5mm}
%beginitem
7. У цьому фрагменті коду
\begin{verbatim}
template <class T> void f(T a1, T a2){}

void g(){
    int i=5;
    f(unique_ptr<int>(new int(++i)), unique_ptr<int>(new int(++i)));
\end{verbatim}

( 1 ) можливий виток ресурсів%bad

( 2 ) є невизначена поведінка%good

( 3 ) є гонка даних%bad

( 4 ) є стан гонки%bad
%enditem
}
%end
\end {large}

\vspace{4mm}
%end
\newpage
%begin
%begin
\begin {large}

\textbf{22.11.2024 Контрольна робота, варіант  \No { 59 }}\par


{
%beginitem
1. По яких днях тижня відбувались лекції з предмету?

%enditem


\vspace{5mm}
%beginitem
2. Як зовуть (ім'я, по батькові) викладача, що вів практичні заняття?

%enditem


\vspace{5mm}
%beginitem
3. Нехай int var=13; Один потік збільшив змінну var на 3, а інший збільшив її в три рази.

гонка даних (data race), стан гонки (race condition)

Виберіть усі правильні твердження (якщо такі є).\par
\vspace{2mm}
( 1 ) Тут гарантовано нема ані гонки даних, ані стану гонки.%bad

( 2 ) Тут гарантовано є стан гонки.%good

( 3 ) Тут може не бути стану гонки.%bad

( 4 ) Тут може бути тільки гонка даних.%bad
%enditem


\vspace{5mm}
%beginitem
4. Виберіть усі правильні твердження (якщо такі є).\par
\vspace{2mm}

( 1 ) У програмі може бути не більше однієї критичної секції.%bad

( 2 ) Клас thread має засоби, щоб повернути результат виконання запущеного потоку.%bad

( 3 ) Виклик \code{this\_thread::sleep\_for(5s)} змушує потік проспати рівно 5 секунд.%bad

( 4 ) Щоб розблокувати заблокований \code{shared\_mutex}, його завжди достатньо розблокувати один раз.%good
%enditem


\vspace{5mm}
%beginitem
5. Виберіть усі правильні твердження (якщо такі є).\par
\vspace{2mm}

( 1 ) На операціях структури даних без блокувань взаємоблокувань потоків виникнути не може.%bad

( 2 ) Структура даних без блокувань гарантує, що жоден потік не буде очікувати доступу до виконання операцій над цією структурою даних.%bad

( 3 ) Структура даних з блокуваннями може бути такою, що використовує атомарні змінні.%good

( 4 ) Реалізації різних операцій над структурою даних повинні використовувати різні м’ютекси.%bad
\newpage

%enditem


\vspace{5mm}
%beginitem
6. Виберіть усі правильні твердження (якщо такі є).\par
\vspace{2mm}

( 1 ) Потік має доступ тільки до своїх власних даних та до глобальних статичних даних.%bad

( 2 ) Автоматичні змінні не можуть розділятися між потоками.%bad

( 3 ) Кожен потік програми може мати доступ до динамічних даних програми.%good

( 4 ) Потік мови С++ не може отримати доступ до даних іншого потоку.%bad
%enditem


\vspace{5mm}
%beginitem
7. У цьому фрагменті коду
\begin{verbatim}
template <class T> void f(T a){}

void g(){
    int i=5;
    f(unique_ptr<int>(new int(i)))+(unique_ptr<int>(new int(++i)));
\end{verbatim}
( 1 ) є гонка даних%bad

( 2 ) можливий виток ресурсів%good

( 3 ) є стан гонки%bad

( 4 ) є невизначена поведінка%good
%enditem
}
%end
\end {large}

\vspace{4mm}
%end
\newpage
%begin
%begin
\begin {large}

\textbf{22.11.2024 Контрольна робота, варіант  \No { 60 }}\par


{
%beginitem
1. По яких днях тижня відбувались лекції з предмету?

%enditem


\vspace{5mm}
%beginitem
2. Як зовуть (ім'я, по батькові) викладача, що вів практичні заняття?

%enditem


\vspace{5mm}
%beginitem
3. Нехай int var=13; Один потік збільшив змінну var на 3, а інший збільшив її в три рази.

гонка даних (data race), стан гонки (race condition)

Виберіть усі правильні твердження (якщо такі є).\par
\vspace{2mm}
( 1 ) Тут може не бути гонки даних.%good

( 2 ) Тут гарантовано нема гонки даних, але може бути стан гонки.%bad

( 3 ) Тут гарантовано є і стан гонки, і гонка даних.%bad

( 4 ) Тут може бути тільки стан гонки.%bad
%enditem


\vspace{5mm}
%beginitem
4. Виберіть усі правильні твердження (якщо такі є).\par
\vspace{2mm}

( 1 ) Розблоковувати незаблокований м’ютекс можна.%bad

( 2 ) Функція \code{sleep\_for} може змусити інший потік заснути.%bad

( 3 ) Не можна блокувати заблокований м’ютекс mutex.%bad

( 4 ) Розблоковувати незаблокований м’ютекс не можна.%good
%enditem


\vspace{5mm}
%beginitem
5. Виберіть усі правильні твердження (якщо такі є).\par
\vspace{2mm}

( 1 ) Мова С++20 має стандартну бібліотеку вільних від блокувань структур даних.%bad

( 2 ) Стандартні контейнери stl безпечно використовувати з кількох потоків.%bad

( 3 ) Структура даних без блокувань може не бути вільною від блокувань.%good

( 4 ) Вільних від очікувань структур даних не існує.%bad
\newpage

%enditem


\vspace{5mm}
%beginitem
6. Виберіть усі правильні твердження (якщо такі є).\par
\vspace{2mm}

( 1 ) Динамічні змінні не можуть розділятися між потоками.%bad

( 2 ) Потік мови С++ не може отримати доступ до автоматичних змінних іншого потоку.%bad

( 3 ) Потік може мати доступ до автоматичних змінних іншого потоку, але це небезпечно.%good

( 4 ) Потік може мати доступ до автоматичних змінних іншого потоку і це безпечно.%bad
%enditem


\vspace{5mm}
%beginitem
7. У цьому фрагменті коду
\begin{verbatim}
template <class T> void f(T a1, T a2){}

void g(){
    f(unique_ptr<int>(new int(3)), unique_ptr<int>(new int(35)));
\end{verbatim}

( 1 ) є гонка даних%bad

( 2 ) можливий виток ресурсів%bad

( 3 ) є стан гонки%bad

( 4 ) є невизначена поведінка%bad
%enditem
}
%end
\end {large}

\vspace{4mm}
%end
\newpage
%begin
%begin
\begin {large}

\textbf{22.11.2024 Контрольна робота, варіант  \No { 61 }}\par


{
%beginitem
1. По яких днях тижня відбувались лекції з предмету?

%enditem


\vspace{5mm}
%beginitem
2. Як зовуть (ім'я, по батькові) викладача, що вів практичні заняття?

%enditem


\vspace{5mm}
%beginitem
3. Нехай int var=13; Один потік збільшив змінну var на 3, а інший збільшив її в три рази.

гонка даних (data race), стан гонки (race condition)

Виберіть усі правильні твердження (якщо такі є).\par
\vspace{2mm}
( 1 ) Тут гарантовано є стан гонки.%good

( 2 ) Тут гарантовано нема ані гонки даних, ані стану гонки.%bad

( 3 ) Тут може не бути стану гонки.%bad

( 4 ) Тут може не бути ані гонки даних, ані стану гонки.%bad
%enditem


\vspace{5mm}
%beginitem
4. Виберіть усі правильні твердження (якщо такі є).\par
\vspace{2mm}

( 1 ) За різних запусків програми на одних тих самих вхідних даних може виникати хронологічно різна послідовність обчислень.%good

( 2 ) Тільки той потік, який обнулив значення семафора, може його збільшити.%bad

( 3 ) У стані гонки програма відчайдушно намагається захопити ресурси процесора.%bad

( 4 ) Об’єкти класів promise та future можна використовувати тільки в багатопотокових програмах.%bad
%enditem


\vspace{5mm}
%beginitem
5. Виберіть усі правильні твердження (якщо такі є).\par
\vspace{2mm}

( 1 ) Вільних від очікувань структур даних не існує.%bad

( 2 ) Взаємоблокування потоків може виникнути тільки за використання м’ютексів.%bad

( 3 ) Кожна вільна від очікувань структура даних є структурою даних, вільною від блокувань.%good

( 4 ) Структура даних без блокувань та вільна від блокувань структура даних --- це одне те саме.%bad
\newpage

%enditem


\vspace{5mm}
%beginitem
6. Виберіть усі правильні твердження (якщо такі є).\par
\vspace{2mm}

( 1 ) Динамічні змінні не можуть розділятися між потоками.%bad

( 2 ) Потік мови С++ не може отримати доступ до автоматичних змінних іншого потоку.%bad

( 3 ) Потік може мати доступ до автоматичних змінних іншого потоку, але це небезпечно.%good

( 4 ) Потік має доступ тільки до своїх власних даних та до глобальних статичних даних.%bad
%enditem


\vspace{5mm}
%beginitem
7. У цьому фрагменті коду
\begin{verbatim}
template <class T> void f(T a1, T a2){}

void g(){
    int i=5;
    f(unique_ptr<int>(new int(i++)), unique_ptr<int>(new int(++i)));
\end{verbatim}

( 1 ) є стан гонки%bad

( 2 ) є гонка даних%bad

( 3 ) можливий виток ресурсів%bad

( 4 ) є невизначена поведінка%good
%enditem
}
%end
\end {large}

\vspace{4mm}
%end
\newpage
%begin
%begin
\begin {large}

\textbf{22.11.2024 Контрольна робота, варіант  \No { 62 }}\par


{
%beginitem
1. По яких днях тижня відбувались лекції з предмету?

%enditem


\vspace{5mm}
%beginitem
2. Як зовуть (ім'я, по батькові) викладача, що вів практичні заняття?

%enditem


\vspace{5mm}
%beginitem
3. Нехай int var=13; Один потік збільшив змінну var на 3, а інший збільшив її в три рази.

гонка даних (data race), стан гонки (race condition)

Виберіть усі правильні твердження (якщо такі є).\par
\vspace{2mm}
( 1 ) Тут гарантовано є і стан гонки, і гонка даних.%bad

( 2 ) Тут гарантовано є гонка даних.%bad

( 3 ) Тут гарантовано нема стану гонки, але може бути гонка даних.%bad

( 4 ) Тут може не бути гонки даних.%good
%enditem


\vspace{5mm}
%beginitem
4. Виберіть усі правильні твердження (якщо такі є).\par
\vspace{2mm}

( 1 ) Досвічений програміст за однопотоковою програмою (текстом мовою С++) завжди може вказати послідовність виконуваних обчислень на конкретних вхідних даних.%bad

( 2 ) Функція async запускає виконання в новому потоці.%bad

( 3 ) Якщо два обчислення програми на С++ тільки читають одну ту саму ділянку пам'яті, то на них гонка даних  виникнути не може.%good

( 4 ) Секція є критичною, якщо під час її виконання не можна генерувати винятки.%bad
%enditem


\vspace{5mm}
%beginitem
5. Виберіть усі правильні твердження (якщо такі є).\par
\vspace{2mm}

( 1 ) Стандартні контейнери stl безпечно використовувати з кількох потоків.%bad

( 2 ) Структура даних є структурою даних без блокувань, якщо її реалізація не використовує м’ютекси.%bad

( 3 ) За використання структурою даних рівно одного м’ютексу, взаємоблокування на її операціях неможливі.%good

( 4 ) Структура даних без блокувань гарантує завершення виконання довільної операції над цією структурою даних за обмежений час.%bad
\newpage

%enditem


\vspace{5mm}
%beginitem
6. Виберіть усі правильні твердження (якщо такі є).\par
\vspace{2mm}

( 1 ) Автоматичні змінні не можуть розділятися між потоками.%bad

( 2 ) Потік мови С++ не може отримати доступ до даних іншого потоку.%bad

( 3 ) Потік може мати доступ до автоматичних змінних іншого потоку і це безпечно.%bad

( 4 ) Кожен потік програми може мати доступ до динамічних даних програми.%good
%enditem


\vspace{5mm}
%beginitem
7. У цьому фрагменті коду
\begin{verbatim}
template <class T> void f(T a1, T a2){}

void g(){
    int i=5;
    f(unique_ptr<int>(new int(++i)), unique_ptr<int>(new int(++i)));
\end{verbatim}

( 1 ) можливий виток ресурсів%bad

( 2 ) є невизначена поведінка%good

( 3 ) є стан гонки%bad

( 4 ) є гонка даних%bad
%enditem
}
%end
\end {large}

\vspace{4mm}
%end
\newpage
%begin
%begin
\begin {large}

\textbf{22.11.2024 Контрольна робота, варіант  \No { 63 }}\par


{
%beginitem
1. По яких днях тижня відбувались лекції з предмету?

%enditem


\vspace{5mm}
%beginitem
2. Як зовуть (ім'я, по батькові) викладача, що вів практичні заняття?

%enditem


\vspace{5mm}
%beginitem
3. Нехай int var=13; Один потік збільшив змінну var на 3, а інший збільшив її в три рази.

гонка даних (data race), стан гонки (race condition)

Виберіть усі правильні твердження (якщо такі є).\par
\vspace{2mm}
( 1 ) Тут гарантовано є стан гонки, але гонки даних може не бути.%good

( 2 ) Тут може бути тільки стан гонки.%bad

( 3 ) Тут гарантовано нема гонки даних, але може бути стан гонки.%bad

( 4 ) Тут може бути тільки гонка даних.%bad
%enditem


\vspace{5mm}
%beginitem
4. Виберіть усі правильні твердження (якщо такі є).\par
\vspace{2mm}

( 1 ) Розблоковувати незаблокований м’ютекс не можна.%good

( 2 ) Об’єкт future можна скопіювати та роздати копії різним потокам, щоб вони змогли побачити результат.%bad

( 3 ) Значення, що є результатом виконання функції, запущеної через async, зберігається безпосередньо в об’єкті ф’ючерса (future).%bad

( 4 ) Коли програма спить, то вона не виконується.%bad
%enditem


\vspace{5mm}
%beginitem
5. Виберіть усі правильні твердження (якщо такі є).\par
\vspace{2mm}

( 1 ) Операція над вільною від блокувань структурою даних завжди виконається за обмежений час.%bad

( 2 ) Мова С++20 має стандартну бібліотеку вільних від блокувань структур даних.%bad

( 3 ) Реалізація структури даних з блокуваннями може використовувати атомарні змінні.%good

( 4 ) Вільна від блокувань структура даних гарантує, що жоден потік не буде очікувати доступу до виконання операцій над цією структурою даних.%bad
\newpage

%enditem


\vspace{5mm}
%beginitem
6. Виберіть усі правильні твердження (якщо такі є).\par
\vspace{2mm}

( 1 ) Потік мови С++ не може отримати доступ до даних іншого потоку.%bad

( 2 ) Автоматичні змінні не можуть розділятися між потоками.%bad

( 3 ) Кожен потік програми може мати доступ до динамічних даних програми.%good

( 4 ) Потік мови С++ не може отримати доступ до автоматичних змінних іншого потоку.%bad
%enditem


\vspace{5mm}
%beginitem
7. У цьому фрагменті коду
\begin{verbatim}
template <class T> void f(T a){}

void g(){
    int i=5;
    f(unique_ptr<int>(new int(i)))+(unique_ptr<int>(new int(35)));
\end{verbatim}
( 1 ) є невизначена поведінка%bad

( 2 ) можливий виток ресурсів%good

( 3 ) є гонка даних%bad

( 4 ) є стан гонки%bad
%enditem
}
%end
\end {large}

\vspace{4mm}
%end
\newpage
%begin
%begin
\begin {large}

\textbf{22.11.2024 Контрольна робота, варіант  \No { 64 }}\par


{
%beginitem
1. По яких днях тижня відбувались лекції з предмету?

%enditem


\vspace{5mm}
%beginitem
2. Як зовуть (ім'я, по батькові) викладача, що вів практичні заняття?

%enditem


\vspace{5mm}
%beginitem
3. Нехай int var=13; Один потік збільшив змінну var на 3, а інший збільшив її в три рази.

гонка даних (data race), стан гонки (race condition)

Виберіть усі правильні твердження (якщо такі є).\par
\vspace{2mm}
( 1 ) Тут гарантовано нема гонки даних, але може бути стан гонки.%bad

( 2 ) Тут може не бути гонки даних.%good

( 3 ) Тут гарантовано нема стану гонки, але може бути гонка даних.%bad

( 4 ) Тут може бути тільки гонка даних.%bad
%enditem


\vspace{5mm}
%beginitem
4. Виберіть усі правильні твердження (якщо такі є).\par
\vspace{2mm}

( 1 ) Клас thread не має засобів, щоб повернути результат виконання запущеного потоку.%good

( 2 ) Об’єкт promise можна скопіювати та роздати копії різним потокам, щоб кожен з них міг передати результат своєї роботи.%bad

( 3 ) М’ютекс має бути записаний на початку критичної секції, яку він забезпечує.%bad

( 4 ) Довільний потік може розблокувати заблокований м’ютекс.%bad
%enditem


\vspace{5mm}
%beginitem
5. Виберіть усі правильні твердження (якщо такі є).\par
\vspace{2mm}

( 1 ) На операціях структури даних без блокувань взаємоблокувань потоків виникнути не може.%bad

( 2 ) Структура даних без блокувань зобов’язана використовувати атомарні змінні.%bad

( 3 ) За використання атомарних змінних структура даних є структурою даних без блокувань.%bad

( 4 ) Операція над вільною від блокувань структурою даних може виконуватись скільки завгодно довго.%good
\newpage

%enditem


\vspace{5mm}
%beginitem
6. Виберіть усі правильні твердження (якщо такі є).\par
\vspace{2mm}

( 1 ) Динамічні змінні не можуть розділятися між потоками.%bad

( 2 ) Потік може мати доступ до автоматичних змінних іншого потоку і це безпечно.%bad

( 3 ) Потік має доступ тільки до своїх власних даних та до глобальних статичних даних.%bad

( 4 ) Потік може мати доступ до автоматичних змінних іншого потоку, але це небезпечно.%good
%enditem


\vspace{5mm}
%beginitem
7. У цьому фрагменті коду
\begin{verbatim}
template <class T> void f(T a1, T a2){}

void g(){
    int i=5;
    f(unique_ptr<int>(new int(++i)), unique_ptr<int>(new int(++i)));
\end{verbatim}

( 1 ) є невизначена поведінка%good

( 2 ) є гонка даних%bad

( 3 ) є стан гонки%bad

( 4 ) можливий виток ресурсів%bad
%enditem
}
%end
\end {large}

\vspace{4mm}
%end
\newpage
%begin
%begin
\begin {large}

\textbf{22.11.2024 Контрольна робота, варіант  \No { 65 }}\par


{
%beginitem
1. По яких днях тижня відбувались лекції з предмету?

%enditem


\vspace{5mm}
%beginitem
2. Як зовуть (ім'я, по батькові) викладача, що вів практичні заняття?

%enditem


\vspace{5mm}
%beginitem
3. Нехай int var=13; Один потік збільшив змінну var на 3, а інший збільшив її в три рази.

гонка даних (data race), стан гонки (race condition)

Виберіть усі правильні твердження (якщо такі є).\par
\vspace{2mm}
( 1 ) Тут гарантовано нема ані гонки даних, ані стану гонки.%bad

( 2 ) Тут гарантовано є стан гонки, але гонки даних може не бути.%good

( 3 ) Тут може не бути стану гонки.%bad

( 4 ) Тут може не бути ані гонки даних, ані стану гонки.%bad
%enditem


\vspace{5mm}
%beginitem
4. Виберіть усі правильні твердження (якщо такі є).\par
\vspace{2mm}

( 1 ) Щоб розблокувати заблокований м’ютекс \code{recursive\_mutex}, його завжди достатньо розблокувати один раз.%bad

( 2 ) Виклик \code{this\_thread::sleep\_for(5s)} змушує потік проспати рівно 5 секунд.%bad

( 3 ) Розблокувати м’ютекс може тільки той потік, що його заблокував.%good

( 4 ) Критична секція це неперервна ділянка коду програми.%bad
%enditem


\vspace{5mm}
%beginitem
5. Виберіть усі правильні твердження (якщо такі є).\par
\vspace{2mm}

( 1 ) Структура даних без блокувань гарантує, що жоден потік не буде очікувати доступу до виконання операцій над цією структурою даних.%bad

( 2 ) Для уникнення гонки даних структура даних може використовувати не більше одного м’ютекса.%bad

( 3 ) Реалізація структури даних без блокувань не може використовувати м’ютекси.%good

( 4 ) Операція над структурою даних без блокувань завжди виконається за обмежений час.%bad
\newpage

%enditem


\vspace{5mm}
%beginitem
6. Виберіть усі правильні твердження (якщо такі є).\par
\vspace{2mm}

( 1 ) Потік може мати доступ до автоматичних змінних іншого потоку, але це небезпечно.%good

( 2 ) Автоматичні змінні не можуть розділятися між потоками.%bad

( 3 ) Потік може мати доступ до автоматичних змінних іншого потоку і це безпечно.%bad

( 4 ) Потік мови С++ не може отримати доступ до автоматичних змінних іншого потоку.%bad
%enditem


\vspace{5mm}
%beginitem
7. У цьому фрагменті коду
\begin{verbatim}
template <class T> void f(T a1, T a2){}

void g(){
    int i=5;
    f(unique_ptr<int>(new int(i++)), unique_ptr<int>(new int(++i)));
\end{verbatim}

( 1 ) є гонка даних%bad

( 2 ) є невизначена поведінка%good

( 3 ) є стан гонки%bad

( 4 ) можливий виток ресурсів%bad
%enditem
}
%end
\end {large}

\vspace{4mm}
%end
\newpage
%begin
%begin
\begin {large}

\textbf{22.11.2024 Контрольна робота, варіант  \No { 66 }}\par


{
%beginitem
1. По яких днях тижня відбувались лекції з предмету?

%enditem


\vspace{5mm}
%beginitem
2. Як зовуть (ім'я, по батькові) викладача, що вів практичні заняття?

%enditem


\vspace{5mm}
%beginitem
3. Нехай int var=13; Один потік збільшив змінну var на 3, а інший збільшив її в три рази.

гонка даних (data race), стан гонки (race condition)

Виберіть усі правильні твердження (якщо такі є).\par
\vspace{2mm}
( 1 ) Тут гарантовано є і стан гонки, і гонка даних.%bad

( 2 ) Тут гарантовано є гонка даних.%bad

( 3 ) Тут гарантовано є стан гонки.%good

( 4 ) Тут може бути тільки стан гонки.%bad
%enditem


\vspace{5mm}
%beginitem
4. Виберіть усі правильні твердження (якщо такі є).\par
\vspace{2mm}

( 1 ) Секція є критичною, якщо під час її виконання не можна генерувати винятки.%bad

( 2 ) У програмі може бути не більше однієї критичної секції.%bad

( 3 ) Функція \code{sleep\_for} може змусити інший потік заснути.%bad

( 4 ) У синтаксично коректній програмі мовою С++ стан гонки (race condition) може виникати.%good
%enditem


\vspace{5mm}
%beginitem
5. Виберіть усі правильні твердження (якщо такі є).\par
\vspace{2mm}

( 1 ) Кожна вільна від очікувань структура даних є структурою даних, вільною від блокувань.%good

( 2 ) Мова С++20 має стандартну бібліотеку конкурентних структур даних.%bad

( 3 ) Одночасно використовувати атомарних змінних більше, ніж регістрів процесора, у програмі не можна.%bad

( 4 ) Реалізації різних операцій над структурою даних повинні використовувати різні м’ютекси.%bad
\newpage

%enditem


\vspace{5mm}
%beginitem
6. Виберіть усі правильні твердження (якщо такі є).\par
\vspace{2mm}

( 1 ) Потік має доступ тільки до своїх власних даних та до глобальних статичних даних.%bad

( 2 ) Потік мови С++ не може отримати доступ до даних іншого потоку.%bad

( 3 ) Кожен потік програми може мати доступ до динамічних даних програми.%good

( 4 ) Динамічні змінні не можуть розділятися між потоками.%bad
%enditem


\vspace{5mm}
%beginitem
7. У цьому фрагменті коду
\begin{verbatim}
template <class T> void f(T a1, T a2){}

void g(){
    f(unique_ptr<int>(new int(3)), unique_ptr<int>(new int(35)));
\end{verbatim}

( 1 ) є гонка даних%bad

( 2 ) є стан гонки%bad

( 3 ) можливий виток ресурсів%bad

( 4 ) є невизначена поведінка%bad
%enditem
}
%end
\end {large}

\vspace{4mm}
%end
\newpage
%begin
%begin
\begin {large}

\textbf{22.11.2024 Контрольна робота, варіант  \No { 67 }}\par


{
%beginitem
1. По яких днях тижня відбувались лекції з предмету?

%enditem


\vspace{5mm}
%beginitem
2. Як зовуть (ім'я, по батькові) викладача, що вів практичні заняття?

%enditem


\vspace{5mm}
%beginitem
3. Нехай int var=13; Один потік збільшив змінну var на 3, а інший збільшив її в три рази.

гонка даних (data race), стан гонки (race condition)

Виберіть усі правильні твердження (якщо такі є).\par
\vspace{2mm}
( 1 ) Тут гарантовано є гонка даних.%bad

( 2 ) Тут гарантовано нема ані гонки даних, ані стану гонки.%bad

( 3 ) Тут гарантовано нема стану гонки, але може бути гонка даних.%bad

( 4 ) Тут може не бути гонки даних.%good
%enditem


\vspace{5mm}
%beginitem
4. Виберіть усі правильні твердження (якщо такі є).\par
\vspace{2mm}

( 1 ) Текст однопотокової програми мовою С++ однозначно визначає послідовність виконуваних обчислень.%bad

( 2 ) Потік в активному очікуванні споживає ресурс процесора.%good

( 3 ) У стані гонки програма відчайдушно намагається захопити ресурси процесора.%bad

( 4 ) Клас thread має засоби, щоб повернути результат виконання запущеного потоку.%bad
%enditem


\vspace{5mm}
%beginitem
5. Виберіть усі правильні твердження (якщо такі є).\par
\vspace{2mm}

( 1 ) Структура даних без блокувань може не бути вільною від блокувань.%good

( 2 ) Вільних від очікувань структур даних не існує.%bad

( 3 ) Реалізації різних операцій над структурою даних повинні використовувати різні м’ютекси.%bad

( 4 ) Стандартні контейнери stl безпечно використовувати з кількох потоків.%bad
\newpage

%enditem


\vspace{5mm}
%beginitem
6. Виберіть усі правильні твердження (якщо такі є).\par
\vspace{2mm}

( 1 ) Динамічні змінні не можуть розділятися між потоками.%bad

( 2 ) Потік може мати доступ до автоматичних змінних іншого потоку і це безпечно.%bad

( 3 ) Кожен потік програми може мати доступ до динамічних даних програми.%good

( 4 ) Потік мови С++ не може отримати доступ до автоматичних змінних іншого потоку.%bad
%enditem


\vspace{5mm}
%beginitem
7. У цьому фрагменті коду
\begin{verbatim}
template <class T> void f(T a){}

void g(){
    int i=5;
    f(unique_ptr<int>(new int(i)))+(unique_ptr<int>(new int(++i)));
\end{verbatim}
( 1 ) є невизначена поведінка%good

( 2 ) є стан гонки%bad

( 3 ) можливий виток ресурсів%good

( 4 ) є гонка даних%bad
%enditem
}
%end
\end {large}

\vspace{4mm}
%end
\newpage
%begin
%begin
\begin {large}

\textbf{22.11.2024 Контрольна робота, варіант  \No { 68 }}\par


{
%beginitem
1. По яких днях тижня відбувались лекції з предмету?

%enditem


\vspace{5mm}
%beginitem
2. Як зовуть (ім'я, по батькові) викладача, що вів практичні заняття?

%enditem


\vspace{5mm}
%beginitem
3. Нехай int var=13; Один потік збільшив змінну var на 3, а інший збільшив її в три рази.

гонка даних (data race), стан гонки (race condition)

Виберіть усі правильні твердження (якщо такі є).\par
\vspace{2mm}
( 1 ) Тут може не бути стану гонки.%bad

( 2 ) Тут гарантовано є стан гонки, але гонки даних може не бути.%good

( 3 ) Тут гарантовано є і стан гонки, і гонка даних.%bad

( 4 ) Тут може не бути ані гонки даних, ані стану гонки.%bad
%enditem


\vspace{5mm}
%beginitem
4. Виберіть усі правильні твердження (якщо такі є).\par
\vspace{2mm}

( 1 ) Розблоковувати незаблокований м’ютекс можна.%bad

( 2 ) Функція async запускає виконання в новому потоці.%bad

( 3 ) Чи запустить функція  async виконання в новому потоці можна задавати політикою.%good

( 4 ) М’ютекс має бути записаний на початку критичної секції, яку він забезпечує.%bad
%enditem


\vspace{5mm}
%beginitem
5. Виберіть усі правильні твердження (якщо такі є).\par
\vspace{2mm}

( 1 ) Вільна від блокувань структура даних гарантує, що жоден потік не буде очікувати доступу до виконання операцій над цією структурою даних.%bad

( 2 ) Структура даних без блокувань зобов’язана використовувати атомарні змінні.%bad

( 3 ) Реалізація структури даних з блокуваннями може використовувати атомарні змінні.%good

( 4 ) Мова С++20 має стандартну бібліотеку вільних від блокувань структур даних.%bad
\newpage

%enditem


\vspace{5mm}
%beginitem
6. Виберіть усі правильні твердження (якщо такі є).\par
\vspace{2mm}

( 1 ) Потік має доступ тільки до своїх власних даних та до глобальних статичних даних.%bad

( 2 ) Потік може мати доступ до автоматичних змінних іншого потоку, але це небезпечно.%good

( 3 ) Автоматичні змінні не можуть розділятися між потоками.%bad

( 4 ) Потік мови С++ не може отримати доступ до даних іншого потоку.%bad
%enditem


\vspace{5mm}
%beginitem
7. У цьому фрагменті коду
\begin{verbatim}
template <class T> void f(T a){}

void g(){
    f(unique_ptr<int>(new int(3)))+(unique_ptr<int>(new int(35)));
\end{verbatim}

( 1 ) є гонка даних%bad

( 2 ) є невизначена поведінка%bad

( 3 ) можливий виток ресурсів%good

( 4 ) є стан гонки%bad
%enditem
}
%end
\end {large}

\vspace{4mm}
%end
\newpage
%begin
%begin
\begin {large}

\textbf{22.11.2024 Контрольна робота, варіант  \No { 69 }}\par


{
%beginitem
1. По яких днях тижня відбувались лекції з предмету?

%enditem


\vspace{5mm}
%beginitem
2. Як зовуть (ім'я, по батькові) викладача, що вів практичні заняття?

%enditem


\vspace{5mm}
%beginitem
3. Нехай int var=13; Один потік збільшив змінну var на 3, а інший збільшив її в три рази.

гонка даних (data race), стан гонки (race condition)

Виберіть усі правильні твердження (якщо такі є).\par
\vspace{2mm}
( 1 ) Тут може бути тільки гонка даних.%bad

( 2 ) Тут гарантовано є стан гонки.%good

( 3 ) Тут гарантовано нема гонки даних, але може бути стан гонки.%bad

( 4 ) Тут може бути тільки стан гонки.%bad
%enditem


\vspace{5mm}
%beginitem
4. Виберіть усі правильні твердження (якщо такі є).\par
\vspace{2mm}

( 1 ) Значення, що є результатом виконання функції, запущеної через async, зберігається безпосередньо в об’єкті ф’ючерса (future).%bad

( 2 ) Об’єкт future можна скопіювати та роздати копії різним потокам, щоб вони змогли побачити результат.%bad

( 3 ) Критична секція може складатися з кількох ділянок коду, які можуть розташовуватись де завгодно (не обов’язково зв’язний фрагмент коду).%good

( 4 ) Об’єкти класів promise та future можна використовувати тільки в багатопотокових програмах.%bad
%enditem


\vspace{5mm}
%beginitem
5. Виберіть усі правильні твердження (якщо такі є).\par
\vspace{2mm}

( 1 ) Якщо весь код структури даних утворює одну критичну секцію, то така структура даних може використовуватись в конкурентній програмі, але по-справжньо\-му конкурентною не є.%good

( 2 ) На операціях структури даних без блокувань взаємоблокувань потоків виникнути не може.%bad

( 3 ) Мова С++20 має стандартну бібліотеку конкурентних структур даних.%bad

( 4 ) Операція над вільною від блокувань структурою даних завжди виконається за обмежений час.%bad
\newpage

%enditem


\vspace{5mm}
%beginitem
6. Виберіть усі правильні твердження (якщо такі є).\par
\vspace{2mm}

( 1 ) Кожен потік програми може мати доступ до динамічних даних програми.%good

( 2 ) Потік мови С++ не може отримати доступ до даних іншого потоку.%bad

( 3 ) Динамічні змінні не можуть розділятися між потоками.%bad

( 4 ) Потік має доступ тільки до своїх власних даних та до глобальних статичних даних.%bad
%enditem


\vspace{5mm}
%beginitem
7. У цьому фрагменті коду
\begin{verbatim}
template <class T> void f(T a1, T a2){}

void g(){
    int i=5;
    f(unique_ptr<int>(new int(++i)), unique_ptr<int>(new int(++i)));
\end{verbatim}

( 1 ) є стан гонки%bad

( 2 ) є невизначена поведінка%good

( 3 ) є гонка даних%bad

( 4 ) можливий виток ресурсів%bad
%enditem
}
%end
\end {large}

\vspace{4mm}
%end
\newpage
%begin
%begin
\begin {large}

\textbf{22.11.2024 Контрольна робота, варіант  \No { 70 }}\par


{
%beginitem
1. По яких днях тижня відбувались лекції з предмету?

%enditem


\vspace{5mm}
%beginitem
2. Як зовуть (ім'я, по батькові) викладача, що вів практичні заняття?

%enditem


\vspace{5mm}
%beginitem
3. Нехай int var=13; Один потік збільшив змінну var на 3, а інший збільшив її в три рази.

гонка даних (data race), стан гонки (race condition)

Виберіть усі правильні твердження (якщо такі є).\par
\vspace{2mm}
( 1 ) Тут може не бути стану гонки.%bad

( 2 ) Тут може не бути ані гонки даних, ані стану гонки.%bad

( 3 ) Тут може не бути гонки даних.%good

( 4 ) Тут гарантовано є гонка даних.%bad
%enditem


\vspace{5mm}
%beginitem
4. Виберіть усі правильні твердження (якщо такі є).\par
\vspace{2mm}

( 1 ) Досвічений програміст за однопотоковою програмою (текстом мовою С++) завжди може вказати послідовність виконуваних обчислень на конкретних вхідних даних.%bad

( 2 ) Об’єкт promise можна скопіювати та роздати копії різним потокам, щоб кожен з них міг передати результат своєї роботи.%bad

( 3 ) Не можна блокувати заблокований м’ютекс mutex.%bad

( 4 ) Критичну секцію можна реалізувати за допомогою м’ютексу.%good
%enditem


\vspace{5mm}
%beginitem
5. Виберіть усі правильні твердження (якщо такі є).\par
\vspace{2mm}

( 1 ) Структура даних є структурою даних без блокувань, якщо її реалізація не використовує м’ютекси.%bad

( 2 ) Структура даних може використовувати м’ютекси та атомарні змінні одночасно.%good

( 3 ) За використання атомарних змінних структура даних є структурою даних без блокувань.%bad

( 4 ) Структура даних без блокувань гарантує завершення виконання довільної операції над цією структурою даних за обмежений час.%bad
\newpage

%enditem


\vspace{5mm}
%beginitem
6. Виберіть усі правильні твердження (якщо такі є).\par
\vspace{2mm}

( 1 ) Автоматичні змінні не можуть розділятися між потоками.%bad

( 2 ) Потік може мати доступ до автоматичних змінних іншого потоку, але це небезпечно.%good

( 3 ) Потік може мати доступ до автоматичних змінних іншого потоку і це безпечно.%bad

( 4 ) Потік мови С++ не може отримати доступ до автоматичних змінних іншого потоку.%bad
%enditem


\vspace{5mm}
%beginitem
7. У цьому фрагменті коду
\begin{verbatim}
template <class T> void f(T a){}

void g(){
    int i=5;
    f(unique_ptr<int>(new int(i)))+(unique_ptr<int>(new int(35)));
\end{verbatim}
( 1 ) є гонка даних%bad

( 2 ) є стан гонки%bad

( 3 ) є невизначена поведінка%bad

( 4 ) можливий виток ресурсів%good
%enditem
}
%end
\end {large}

\vspace{4mm}
%end
\newpage
%begin
%begin
\begin {large}

\textbf{22.11.2024 Контрольна робота, варіант  \No { 71 }}\par


{
%beginitem
1. По яких днях тижня відбувались лекції з предмету?

%enditem


\vspace{5mm}
%beginitem
2. Як зовуть (ім'я, по батькові) викладача, що вів практичні заняття?

%enditem


\vspace{5mm}
%beginitem
3. Нехай int var=13; Один потік збільшив змінну var на 3, а інший збільшив її в три рази.

гонка даних (data race), стан гонки (race condition)

Виберіть усі правильні твердження (якщо такі є).\par
\vspace{2mm}
( 1 ) Тут гарантовано нема ані гонки даних, ані стану гонки.%bad

( 2 ) Тут гарантовано є стан гонки.%good

( 3 ) Тут гарантовано нема гонки даних, але може бути стан гонки.%bad

( 4 ) Тут може бути тільки гонка даних.%bad
%enditem


\vspace{5mm}
%beginitem
4. Виберіть усі правильні твердження (якщо такі є).\par
\vspace{2mm}

( 1 ) Коли програма спить, то вона не виконується.%bad

( 2 ) Тільки той потік, який обнулив значення семафора, може його збільшити.%bad

( 3 ) Порядок обчислень у потоці для одних тих самих вхідних даних може залежати від використовуваного компілятора.%good

( 4 ) Потік, запущений конструктором класу jthread, має обов’язково зупинитися, якщо пов’язаний об’єкт надішле команду зупинитися.%bad
%enditem


\vspace{5mm}
%beginitem
5. Виберіть усі правильні твердження (якщо такі є).\par
\vspace{2mm}

( 1 ) Структура даних без блокувань та вільна від блокувань структура даних --- це одне те саме.%bad

( 2 ) Структура даних без блокувань гарантує, що жоден потік не буде очікувати доступу до виконання операцій над цією структурою даних.%bad

( 3 ) За використання структурою даних рівно одного м’ютексу, взаємоблокування на її операціях неможливі.%good

( 4 ) Для уникнення гонки даних структура даних може використовувати не більше одного м’ютекса.%bad
\newpage

%enditem


\vspace{5mm}
%beginitem
6. Виберіть усі правильні твердження (якщо такі є).\par
\vspace{2mm}

( 1 ) Потік може мати доступ до автоматичних змінних іншого потоку, але це небезпечно.%good

( 2 ) Динамічні змінні не можуть розділятися між потоками.%bad

( 3 ) Потік може мати доступ до автоматичних змінних іншого потоку і це безпечно.%bad

( 4 ) Потік мови С++ не може отримати доступ до автоматичних змінних іншого потоку.%bad
%enditem


\vspace{5mm}
%beginitem
7. У цьому фрагменті коду
\begin{verbatim}
template <class T> void f(T a1, T a2){}

void g(){
    int i=5;
    f(unique_ptr<int>(new int(++i)), unique_ptr<int>(new int(++i)));
\end{verbatim}

( 1 ) є невизначена поведінка%good

( 2 ) є гонка даних%bad

( 3 ) можливий виток ресурсів%bad

( 4 ) є стан гонки%bad
%enditem
}
%end
\end {large}

\vspace{4mm}
%end
\newpage
%begin
%begin
\begin {large}

\textbf{22.11.2024 Контрольна робота, варіант  \No { 72 }}\par


{
%beginitem
1. По яких днях тижня відбувались лекції з предмету?

%enditem


\vspace{5mm}
%beginitem
2. Як зовуть (ім'я, по батькові) викладача, що вів практичні заняття?

%enditem


\vspace{5mm}
%beginitem
3. Нехай int var=13; Один потік збільшив змінну var на 3, а інший збільшив її в три рази.

гонка даних (data race), стан гонки (race condition)

Виберіть усі правильні твердження (якщо такі є).\par
\vspace{2mm}
( 1 ) Тут гарантовано є і стан гонки, і гонка даних.%bad

( 2 ) Тут може бути тільки стан гонки.%bad

( 3 ) Тут гарантовано нема стану гонки, але може бути гонка даних.%bad

( 4 ) Тут гарантовано є стан гонки, але гонки даних може не бути.%good
%enditem


\vspace{5mm}
%beginitem
4. Виберіть усі правильні твердження (якщо такі є).\par
\vspace{2mm}

( 1 ) Блокувати заблокований м’ютекс  \code{mutex} іноді можна.%good

( 2 ) Значення, що є результатом виконання функції, запущеної через async, зберігається безпосередньо в об’єкті проміза (promise).%bad

( 3 ) Довільний потік може розблокувати заблокований м’ютекс.%bad

( 4 ) Тільки від компілятора залежить, чи запустить функція async виконання в новому потоці.%bad
%enditem


\vspace{5mm}
%beginitem
5. Виберіть усі правильні твердження (якщо такі є).\par
\vspace{2mm}

( 1 ) Структура даних з блокуваннями може бути такою, що використовує атомарні змінні.%good

( 2 ) Одночасно використовувати атомарних змінних більше, ніж регістрів процесора, у програмі не можна.%bad

( 3 ) Взаємоблокування потоків може виникнути тільки за використання м’ютексів.%bad

( 4 ) Операція над структурою даних без блокувань завжди виконається за обмежений час.%bad
\newpage

%enditem


\vspace{5mm}
%beginitem
6. Виберіть усі правильні твердження (якщо такі є).\par
\vspace{2mm}

( 1 ) Кожен потік програми може мати доступ до динамічних даних програми.%good

( 2 ) Автоматичні змінні не можуть розділятися між потоками.%bad

( 3 ) Потік має доступ тільки до своїх власних даних та до глобальних статичних даних.%bad

( 4 ) Потік мови С++ не може отримати доступ до даних іншого потоку.%bad
%enditem


\vspace{5mm}
%beginitem
7. У цьому фрагменті коду
\begin{verbatim}
template <class T> void f(T a){}

void g(){
    int i=5;
    f(unique_ptr<int>(new int(i)))+(unique_ptr<int>(new int(35)));
\end{verbatim}
( 1 ) є стан гонки%bad

( 2 ) є гонка даних%bad

( 3 ) є невизначена поведінка%bad

( 4 ) можливий виток ресурсів%good
%enditem
}
%end
\end {large}

\vspace{4mm}
%end
\newpage
%begin
%begin
\begin {large}

\textbf{22.11.2024 Контрольна робота, варіант  \No { 73 }}\par


{
%beginitem
1. По яких днях тижня відбувались лекції з предмету?

%enditem


\vspace{5mm}
%beginitem
2. Як зовуть (ім'я, по батькові) викладача, що вів практичні заняття?

%enditem


\vspace{5mm}
%beginitem
3. Нехай int var=13; Один потік збільшив змінну var на 3, а інший збільшив її в три рази.

гонка даних (data race), стан гонки (race condition)

Виберіть усі правильні твердження (якщо такі є).\par
\vspace{2mm}
( 1 ) Тут гарантовано нема стану гонки, але може бути гонка даних.%bad

( 2 ) Тут може бути тільки стан гонки.%bad

( 3 ) Тут гарантовано є стан гонки, але гонки даних може не бути.%good

( 4 ) Тут гарантовано нема ані гонки даних, ані стану гонки.%bad
%enditem


\vspace{5mm}
%beginitem
4. Виберіть усі правильні твердження (якщо такі є).\par
\vspace{2mm}

( 1 ) Не можна блокувати заблокований м’ютекс mutex.%bad

( 2 ) Коли потік спить, він не споживає ресурс процесора.%good

( 3 ) Функція async запускає виконання в новому потоці.%bad

( 4 ) Критична секція це неперервна ділянка коду програми.%bad
%enditem


\vspace{5mm}
%beginitem
5. Виберіть усі правильні твердження (якщо такі є).\par
\vspace{2mm}

( 1 ) На операціях структури даних без блокувань взаємоблокувань потоків виникнути не може.%bad

( 2 ) Реалізація структури даних з блокуваннями може використовувати атомарні змінні.%good

( 3 ) Для уникнення гонки даних структура даних може використовувати не більше одного м’ютекса.%bad

( 4 ) Вільна від блокувань структура даних гарантує, що жоден потік не буде очікувати доступу до виконання операцій над цією структурою даних.%bad
\newpage

%enditem


\vspace{5mm}
%beginitem
6. Виберіть усі правильні твердження (якщо такі є).\par
\vspace{2mm}

( 1 ) Потік має доступ тільки до своїх власних даних та до глобальних статичних даних.%bad

( 2 ) Динамічні змінні не можуть розділятися між потоками.%bad

( 3 ) Автоматичні змінні не можуть розділятися між потоками.%bad

( 4 ) Кожен потік програми може мати доступ до динамічних даних програми.%good
%enditem


\vspace{5mm}
%beginitem
7. У цьому фрагменті коду
\begin{verbatim}
template <class T> void f(T a1, T a2){}

void g(){
    int i=5;
    f(unique_ptr<int>(new int(++i)), unique_ptr<int>(new int(++i)));
\end{verbatim}

( 1 ) можливий виток ресурсів%bad

( 2 ) є невизначена поведінка%good

( 3 ) є гонка даних%bad

( 4 ) є стан гонки%bad
%enditem
}
%end
\end {large}

\vspace{4mm}
%end
\newpage
%begin
%begin
\begin {large}

\textbf{22.11.2024 Контрольна робота, варіант  \No { 74 }}\par


{
%beginitem
1. По яких днях тижня відбувались лекції з предмету?

%enditem


\vspace{5mm}
%beginitem
2. Як зовуть (ім'я, по батькові) викладача, що вів практичні заняття?

%enditem


\vspace{5mm}
%beginitem
3. Нехай int var=13; Один потік збільшив змінну var на 3, а інший збільшив її в три рази.

гонка даних (data race), стан гонки (race condition)

Виберіть усі правильні твердження (якщо такі є).\par
\vspace{2mm}
( 1 ) Тут гарантовано нема гонки даних, але може бути стан гонки.%bad

( 2 ) Тут гарантовано є стан гонки.%good

( 3 ) Тут гарантовано є гонка даних.%bad

( 4 ) Тут може бути тільки гонка даних.%bad
%enditem


\vspace{5mm}
%beginitem
4. Виберіть усі правильні твердження (якщо такі є).\par
\vspace{2mm}

( 1 ) М’ютекс має бути записаний на початку критичної секції, яку він забезпечує.%bad

( 2 ) Щоб розблокувати заблокований \code{shared\_mutex}, його завжди достатньо розблокувати один раз.%good

( 3 ) Розблоковувати незаблокований м’ютекс можна.%bad

( 4 ) Значення, що є результатом виконання функції, запущеної через async, зберігається безпосередньо в об’єкті ф’ючерса (future).%bad
%enditem


\vspace{5mm}
%beginitem
5. Виберіть усі правильні твердження (якщо такі є).\par
\vspace{2mm}

( 1 ) Мова С++20 має стандартну бібліотеку вільних від блокувань структур даних.%bad

( 2 ) Мова С++20 має стандартну бібліотеку конкурентних структур даних.%bad

( 3 ) Взаємоблокування потоків може виникнути тільки за використання м’ютексів.%bad

( 4 ) Реалізація структури даних без блокувань не може використовувати м’ютекси.%good
\newpage

%enditem


\vspace{5mm}
%beginitem
6. Виберіть усі правильні твердження (якщо такі є).\par
\vspace{2mm}

( 1 ) Потік може мати доступ до автоматичних змінних іншого потоку, але це небезпечно.%good

( 2 ) Потік мови С++ не може отримати доступ до даних іншого потоку.%bad

( 3 ) Потік може мати доступ до автоматичних змінних іншого потоку і це безпечно.%bad

( 4 ) Потік мови С++ не може отримати доступ до автоматичних змінних іншого потоку.%bad
%enditem


\vspace{5mm}
%beginitem
7. У цьому фрагменті коду
\begin{verbatim}
template <class T> void f(T a){}

void g(){
    int i=5;
    f(unique_ptr<int>(new int(i)))+(unique_ptr<int>(new int(++i)));
\end{verbatim}
( 1 ) є стан гонки%bad

( 2 ) є гонка даних%bad

( 3 ) є невизначена поведінка%good

( 4 ) можливий виток ресурсів%good
%enditem
}
%end
\end {large}

\vspace{4mm}
%end
\newpage
%begin
%begin
\begin {large}

\textbf{22.11.2024 Контрольна робота, варіант  \No { 75 }}\par


{
%beginitem
1. По яких днях тижня відбувались лекції з предмету?

%enditem


\vspace{5mm}
%beginitem
2. Як зовуть (ім'я, по батькові) викладача, що вів практичні заняття?

%enditem


\vspace{5mm}
%beginitem
3. Нехай int var=13; Один потік збільшив змінну var на 3, а інший збільшив її в три рази.

гонка даних (data race), стан гонки (race condition)

Виберіть усі правильні твердження (якщо такі є).\par
\vspace{2mm}
( 1 ) Тут може не бути гонки даних.%good

( 2 ) Тут гарантовано є і стан гонки, і гонка даних.%bad

( 3 ) Тут може не бути ані гонки даних, ані стану гонки.%bad

( 4 ) Тут може не бути стану гонки.%bad
%enditem


\vspace{5mm}
%beginitem
4. Виберіть усі правильні твердження (якщо такі є).\par
\vspace{2mm}

( 1 ) Виклик функції \code{this\_thread::sleep\_for(5s)} змушує потік, з якого його здійснили, заснути якнайменш на 5 секунд.%good

( 2 ) Довільний потік може розблокувати заблокований м’ютекс.%bad

( 3 ) Тільки від компілятора залежить, чи запустить функція async виконання в новому потоці.%bad

( 4 ) Потік, запущений конструктором класу jthread, має обов’язково зупинитися, якщо пов’язаний об’єкт надішле команду зупинитися.%bad
%enditem


\vspace{5mm}
%beginitem
5. Виберіть усі правильні твердження (якщо такі є).\par
\vspace{2mm}

( 1 ) Операція над вільною від блокувань структурою даних завжди виконається за обмежений час.%bad

( 2 ) Одночасно використовувати атомарних змінних більше, ніж регістрів процесора, у програмі не можна.%bad

( 3 ) Структура даних без блокувань гарантує завершення виконання довільної операції над цією структурою даних за обмежений час.%bad

( 4 ) Структура даних без блокувань може не бути вільною від блокувань.%good
\newpage

%enditem


\vspace{5mm}
%beginitem
6. Виберіть усі правильні твердження (якщо такі є).\par
\vspace{2mm}

( 1 ) Потік має доступ тільки до своїх власних даних та до глобальних статичних даних.%bad

( 2 ) Потік може мати доступ до автоматичних змінних іншого потоку і це безпечно.%bad

( 3 ) Кожен потік програми може мати доступ до динамічних даних програми.%good

( 4 ) Автоматичні змінні не можуть розділятися між потоками.%bad
%enditem


\vspace{5mm}
%beginitem
7. У цьому фрагменті коду
\begin{verbatim}
template <class T> void f(T a1, T a2){}

void g(){
    int i=5;
    f(unique_ptr<int>(new int(i++)), unique_ptr<int>(new int(++i)));
\end{verbatim}

( 1 ) є стан гонки%bad

( 2 ) можливий виток ресурсів%bad

( 3 ) є невизначена поведінка%good

( 4 ) є гонка даних%bad
%enditem
}
%end
\end {large}

\vspace{4mm}
%end
\newpage
%begin
%begin
\begin {large}

\textbf{22.11.2024 Контрольна робота, варіант  \No { 76 }}\par


{
%beginitem
1. По яких днях тижня відбувались лекції з предмету?

%enditem


\vspace{5mm}
%beginitem
2. Як зовуть (ім'я, по батькові) викладача, що вів практичні заняття?

%enditem


\vspace{5mm}
%beginitem
3. Нехай int var=13; Один потік збільшив змінну var на 3, а інший збільшив її в три рази.

гонка даних (data race), стан гонки (race condition)

Виберіть усі правильні твердження (якщо такі є).\par
\vspace{2mm}
( 1 ) Тут гарантовано нема ані гонки даних, ані стану гонки.%bad

( 2 ) Тут гарантовано нема гонки даних, але може бути стан гонки.%bad

( 3 ) Тут гарантовано є стан гонки, але гонки даних може не бути.%good

( 4 ) Тут може бути тільки гонка даних.%bad
%enditem


\vspace{5mm}
%beginitem
4. Виберіть усі правильні твердження (якщо такі є).\par
\vspace{2mm}

( 1 ) Досвічений програміст за однопотоковою програмою (текстом мовою С++) завжди може вказати послідовність виконуваних обчислень на конкретних вхідних даних.%bad

( 2 ) Значення, що є результатом виконання функції, запущеної через async, зберігається безпосередньо в об’єкті проміза (promise).%bad

( 3 ) У програмі може бути не більше однієї критичної секції.%bad

( 4 ) За різних запусків програми на одних тих самих вхідних даних може виникати хронологічно різна послідовність обчислень.%good
%enditem


\vspace{5mm}
%beginitem
5. Виберіть усі правильні твердження (якщо такі є).\par
\vspace{2mm}

( 1 ) Реалізації різних операцій над структурою даних повинні використовувати різні м’ютекси.%bad

( 2 ) Структура даних є структурою даних без блокувань, якщо її реалізація не використовує м’ютекси.%bad

( 3 ) Структура даних без блокувань та вільна від блокувань структура даних --- це одне те саме.%bad

( 4 ) Операція над вільною від блокувань структурою даних може виконуватись скільки завгодно довго.%good
\newpage

%enditem


\vspace{5mm}
%beginitem
6. Виберіть усі правильні твердження (якщо такі є).\par
\vspace{2mm}

( 1 ) Потік мови С++ не може отримати доступ до автоматичних змінних іншого потоку.%bad

( 2 ) Потік може мати доступ до автоматичних змінних іншого потоку, але це небезпечно.%good

( 3 ) Динамічні змінні не можуть розділятися між потоками.%bad

( 4 ) Потік мови С++ не може отримати доступ до даних іншого потоку.%bad
%enditem


\vspace{5mm}
%beginitem
7. У цьому фрагменті коду
\begin{verbatim}
template <class T> void f(T a1, T a2){}

void g(){
    f(unique_ptr<int>(new int(3)), unique_ptr<int>(new int(35)));
\end{verbatim}

( 1 ) є гонка даних%bad

( 2 ) можливий виток ресурсів%bad

( 3 ) є стан гонки%bad

( 4 ) є невизначена поведінка%bad
%enditem
}
%end
\end {large}

\vspace{4mm}
%end
\newpage
%begin
%begin
\begin {large}

\textbf{22.11.2024 Контрольна робота, варіант  \No { 77 }}\par


{
%beginitem
1. По яких днях тижня відбувались лекції з предмету?

%enditem


\vspace{5mm}
%beginitem
2. Як зовуть (ім'я, по батькові) викладача, що вів практичні заняття?

%enditem


\vspace{5mm}
%beginitem
3. Нехай int var=13; Один потік збільшив змінну var на 3, а інший збільшив її в три рази.

гонка даних (data race), стан гонки (race condition)

Виберіть усі правильні твердження (якщо такі є).\par
\vspace{2mm}
( 1 ) Тут може не бути ані гонки даних, ані стану гонки.%bad

( 2 ) Тут гарантовано є стан гонки.%good

( 3 ) Тут може бути тільки стан гонки.%bad

( 4 ) Тут гарантовано є і стан гонки, і гонка даних.%bad
%enditem


\vspace{5mm}
%beginitem
4. Виберіть усі правильні твердження (якщо такі є).\par
\vspace{2mm}

( 1 ) Об’єкт future можна скопіювати та роздати копії різним потокам, щоб вони змогли побачити результат.%bad

( 2 ) Текст однопотокової програми мовою С++ однозначно визначає послідовність виконуваних обчислень.%bad

( 3 ) Об’єкти класів promise та future можна використовувати тільки в багатопотокових програмах.%bad

( 4 ) Коли потік спить, він не споживає ресурс процесора.%good
%enditem


\vspace{5mm}
%beginitem
5. Виберіть усі правильні твердження (якщо такі є).\par
\vspace{2mm}

( 1 ) Структура даних без блокувань гарантує, що жоден потік не буде очікувати доступу до виконання операцій над цією структурою даних.%bad

( 2 ) Якщо весь код структури даних утворює одну критичну секцію, то така структура даних може використовуватись в конкурентній програмі, але по-справжньо\-му конкурентною не є.%good

( 3 ) Стандартні контейнери stl безпечно використовувати з кількох потоків.%bad

( 4 ) Вільних від очікувань структур даних не існує.%bad
\newpage

%enditem


\vspace{5mm}
%beginitem
6. Виберіть усі правильні твердження (якщо такі є).\par
\vspace{2mm}

( 1 ) Потік має доступ тільки до своїх власних даних та до глобальних статичних даних.%bad

( 2 ) Кожен потік програми може мати доступ до динамічних даних програми.%good

( 3 ) Потік мови С++ не може отримати доступ до автоматичних змінних іншого потоку.%bad

( 4 ) Потік мови С++ не може отримати доступ до даних іншого потоку.%bad
%enditem


\vspace{5mm}
%beginitem
7. У цьому фрагменті коду
\begin{verbatim}
template <class T> void f(T a){}

void g(){
    f(unique_ptr<int>(new int(3)))+(unique_ptr<int>(new int(35)));
\end{verbatim}

( 1 ) є стан гонки%bad

( 2 ) є гонка даних%bad

( 3 ) можливий виток ресурсів%good

( 4 ) є невизначена поведінка%bad
%enditem
}
%end
\end {large}

\vspace{4mm}
%end
\newpage
%begin
%begin
\begin {large}

\textbf{22.11.2024 Контрольна робота, варіант  \No { 78 }}\par


{
%beginitem
1. По яких днях тижня відбувались лекції з предмету?

%enditem


\vspace{5mm}
%beginitem
2. Як зовуть (ім'я, по батькові) викладача, що вів практичні заняття?

%enditem


\vspace{5mm}
%beginitem
3. Нехай int var=13; Один потік збільшив змінну var на 3, а інший збільшив її в три рази.

гонка даних (data race), стан гонки (race condition)

Виберіть усі правильні твердження (якщо такі є).\par
\vspace{2mm}
( 1 ) Тут гарантовано нема стану гонки, але може бути гонка даних.%bad

( 2 ) Тут гарантовано є гонка даних.%bad

( 3 ) Тут може не бути стану гонки.%bad

( 4 ) Тут може не бути гонки даних.%good
%enditem


\vspace{5mm}
%beginitem
4. Виберіть усі правильні твердження (якщо такі є).\par
\vspace{2mm}

( 1 ) Розблокувати м’ютекс може тільки той потік, що його заблокував.%good

( 2 ) Клас thread має засоби, щоб повернути результат виконання запущеного потоку.%bad

( 3 ) У стані гонки програма відчайдушно намагається захопити ресурси процесора.%bad

( 4 ) Виклик \code{this\_thread::sleep\_for(5s)} змушує потік проспати рівно 5 секунд.%bad
%enditem


\vspace{5mm}
%beginitem
5. Виберіть усі правильні твердження (якщо такі є).\par
\vspace{2mm}

( 1 ) Структура даних без блокувань зобов’язана використовувати атомарні змінні.%bad

( 2 ) За використання атомарних змінних структура даних є структурою даних без блокувань.%bad

( 3 ) Операція над структурою даних без блокувань завжди виконається за обмежений час.%bad

( 4 ) За використання структурою даних рівно одного м’ютексу, взаємоблокування на її операціях неможливі.%good
\newpage

%enditem


\vspace{5mm}
%beginitem
6. Виберіть усі правильні твердження (якщо такі є).\par
\vspace{2mm}

( 1 ) Потік може мати доступ до автоматичних змінних іншого потоку і це безпечно.%bad

( 2 ) Потік може мати доступ до автоматичних змінних іншого потоку, але це небезпечно.%good

( 3 ) Автоматичні змінні не можуть розділятися між потоками.%bad

( 4 ) Динамічні змінні не можуть розділятися між потоками.%bad
%enditem


\vspace{5mm}
%beginitem
7. У цьому фрагменті коду
\begin{verbatim}
template <class T> void f(T a){}

void g(){
    f(unique_ptr<int>(new int(3)))+(unique_ptr<int>(new int(35)));
\end{verbatim}

( 1 ) є стан гонки%bad

( 2 ) є невизначена поведінка%bad

( 3 ) є гонка даних%bad

( 4 ) можливий виток ресурсів%good
%enditem
}
%end
\end {large}

\vspace{4mm}
%end
\newpage
%begin
%begin
\begin {large}

\textbf{22.11.2024 Контрольна робота, варіант  \No { 79 }}\par


{
%beginitem
1. По яких днях тижня відбувались лекції з предмету?

%enditem


\vspace{5mm}
%beginitem
2. Як зовуть (ім'я, по батькові) викладача, що вів практичні заняття?

%enditem


\vspace{5mm}
%beginitem
3. Нехай int var=13; Один потік збільшив змінну var на 3, а інший збільшив її в три рази.

гонка даних (data race), стан гонки (race condition)

Виберіть усі правильні твердження (якщо такі є).\par
\vspace{2mm}
( 1 ) Тут може не бути стану гонки.%bad

( 2 ) Тут гарантовано є стан гонки.%good

( 3 ) Тут може не бути ані гонки даних, ані стану гонки.%bad

( 4 ) Тут гарантовано нема ані гонки даних, ані стану гонки.%bad
%enditem


\vspace{5mm}
%beginitem
4. Виберіть усі правильні твердження (якщо такі є).\par
\vspace{2mm}

( 1 ) Функція \code{sleep\_for} може змусити інший потік заснути.%bad

( 2 ) Секція є критичною, якщо під час її виконання не можна генерувати винятки.%bad

( 3 ) Чи запустить функція  async виконання в новому потоці можна задавати політикою.%good

( 4 ) Об’єкт promise можна скопіювати та роздати копії різним потокам, щоб кожен з них міг передати результат своєї роботи.%bad
%enditem


\vspace{5mm}
%beginitem
5. Виберіть усі правильні твердження (якщо такі є).\par
\vspace{2mm}

( 1 ) Операція над структурою даних без блокувань завжди виконається за обмежений час.%bad

( 2 ) Одночасно використовувати атомарних змінних більше, ніж регістрів процесора, у програмі не можна.%bad

( 3 ) Кожна вільна від очікувань структура даних є структурою даних, вільною від блокувань.%good

( 4 ) Вільна від блокувань структура даних гарантує, що жоден потік не буде очікувати доступу до виконання операцій над цією структурою даних.%bad
\newpage

%enditem


\vspace{5mm}
%beginitem
6. Виберіть усі правильні твердження (якщо такі є).\par
\vspace{2mm}

( 1 ) Потік може мати доступ до автоматичних змінних іншого потоку і це безпечно.%bad

( 2 ) Динамічні змінні не можуть розділятися між потоками.%bad

( 3 ) Автоматичні змінні не можуть розділятися між потоками.%bad

( 4 ) Потік може мати доступ до автоматичних змінних іншого потоку, але це небезпечно.%good
%enditem


\vspace{5mm}
%beginitem
7. У цьому фрагменті коду
\begin{verbatim}
template <class T> void f(T a1, T a2){}

void g(){
    int i=5;
    f(unique_ptr<int>(new int(++i)), unique_ptr<int>(new int(++i)));
\end{verbatim}

( 1 ) є гонка даних%bad

( 2 ) є стан гонки%bad

( 3 ) можливий виток ресурсів%bad

( 4 ) є невизначена поведінка%good
%enditem
}
%end
\end {large}

\vspace{4mm}
%end
\newpage
%begin
%begin
\begin {large}

\textbf{22.11.2024 Контрольна робота, варіант  \No { 80 }}\par


{
%beginitem
1. По яких днях тижня відбувались лекції з предмету?

%enditem


\vspace{5mm}
%beginitem
2. Як зовуть (ім'я, по батькові) викладача, що вів практичні заняття?

%enditem


\vspace{5mm}
%beginitem
3. Нехай int var=13; Один потік збільшив змінну var на 3, а інший збільшив її в три рази.

гонка даних (data race), стан гонки (race condition)

Виберіть усі правильні твердження (якщо такі є).\par
\vspace{2mm}
( 1 ) Тут гарантовано є гонка даних.%bad

( 2 ) Тут гарантовано нема стану гонки, але може бути гонка даних.%bad

( 3 ) Тут може не бути гонки даних.%good

( 4 ) Тут гарантовано нема гонки даних, але може бути стан гонки.%bad
%enditem


\vspace{5mm}
%beginitem
4. Виберіть усі правильні твердження (якщо такі є).\par
\vspace{2mm}

( 1 ) Виклик функції \code{this\_thread::sleep\_for(5s)} змушує потік, з якого його здійснили, заснути якнайменш на 5 секунд.%good

( 2 ) Коли програма спить, то вона не виконується.%bad

( 3 ) Щоб розблокувати заблокований м’ютекс \code{recursive\_mutex}, його завжди достатньо розблокувати один раз.%bad

( 4 ) Тільки той потік, який обнулив значення семафора, може його збільшити.%bad
%enditem


\vspace{5mm}
%beginitem
5. Виберіть усі правильні твердження (якщо такі є).\par
\vspace{2mm}

( 1 ) За використання атомарних змінних структура даних є структурою даних без блокувань.%bad

( 2 ) Структура даних з блокуваннями може бути такою, що використовує атомарні змінні.%good

( 3 ) Структура даних без блокувань та вільна від блокувань структура даних --- це одне те саме.%bad

( 4 ) Для уникнення гонки даних структура даних може використовувати не більше одного м’ютекса.%bad
\newpage

%enditem


\vspace{5mm}
%beginitem
6. Виберіть усі правильні твердження (якщо такі є).\par
\vspace{2mm}

( 1 ) Потік мови С++ не може отримати доступ до даних іншого потоку.%bad

( 2 ) Потік мови С++ не може отримати доступ до автоматичних змінних іншого потоку.%bad

( 3 ) Кожен потік програми може мати доступ до динамічних даних програми.%good

( 4 ) Потік має доступ тільки до своїх власних даних та до глобальних статичних даних.%bad
%enditem


\vspace{5mm}
%beginitem
7. У цьому фрагменті коду
\begin{verbatim}
template <class T> void f(T a){}

void g(){
    int i=5;
    f(unique_ptr<int>(new int(i)))+(unique_ptr<int>(new int(++i)));
\end{verbatim}
( 1 ) можливий виток ресурсів%good

( 2 ) є гонка даних%bad

( 3 ) є стан гонки%bad

( 4 ) є невизначена поведінка%good
%enditem
}
%end
\end {large}

\vspace{4mm}
%end
\newpage
%begin
%begin
\begin {large}

\textbf{22.11.2024 Контрольна робота, варіант  \No { 81 }}\par


{
%beginitem
1. По яких днях тижня відбувались лекції з предмету?

%enditem


\vspace{5mm}
%beginitem
2. Як зовуть (ім'я, по батькові) викладача, що вів практичні заняття?

%enditem


\vspace{5mm}
%beginitem
3. Нехай int var=13; Один потік збільшив змінну var на 3, а інший збільшив її в три рази.

гонка даних (data race), стан гонки (race condition)

Виберіть усі правильні твердження (якщо такі є).\par
\vspace{2mm}
( 1 ) Тут гарантовано є і стан гонки, і гонка даних.%bad

( 2 ) Тут може бути тільки гонка даних.%bad

( 3 ) Тут може бути тільки стан гонки.%bad

( 4 ) Тут гарантовано є стан гонки, але гонки даних може не бути.%good
%enditem


\vspace{5mm}
%beginitem
4. Виберіть усі правильні твердження (якщо такі є).\par
\vspace{2mm}

( 1 ) Текст однопотокової програми мовою С++ однозначно визначає послідовність виконуваних обчислень.%bad

( 2 ) У програмі може бути не більше однієї критичної секції.%bad

( 3 ) Об’єкт future можна скопіювати та роздати копії різним потокам, щоб вони змогли побачити результат.%bad

( 4 ) Клас thread не має засобів, щоб повернути результат виконання запущеного потоку.%good
%enditem


\vspace{5mm}
%beginitem
5. Виберіть усі правильні твердження (якщо такі є).\par
\vspace{2mm}

( 1 ) Взаємоблокування потоків може виникнути тільки за використання м’ютексів.%bad

( 2 ) Структура даних може використовувати м’ютекси та атомарні змінні одночасно.%good

( 3 ) Структура даних є структурою даних без блокувань, якщо її реалізація не використовує м’ютекси.%bad

( 4 ) Мова С++20 має стандартну бібліотеку конкурентних структур даних.%bad
\newpage

%enditem


\vspace{5mm}
%beginitem
6. Виберіть усі правильні твердження (якщо такі є).\par
\vspace{2mm}

( 1 ) Автоматичні змінні не можуть розділятися між потоками.%bad

( 2 ) Потік мови С++ не може отримати доступ до даних іншого потоку.%bad

( 3 ) Кожен потік програми може мати доступ до динамічних даних програми.%good

( 4 ) Потік може мати доступ до автоматичних змінних іншого потоку і це безпечно.%bad
%enditem


\vspace{5mm}
%beginitem
7. У цьому фрагменті коду
\begin{verbatim}
template <class T> void f(T a1, T a2){}

void g(){
    f(unique_ptr<int>(new int(3)), unique_ptr<int>(new int(35)));
\end{verbatim}

( 1 ) є гонка даних%bad

( 2 ) можливий виток ресурсів%bad

( 3 ) є стан гонки%bad

( 4 ) є невизначена поведінка%bad
%enditem
}
%end
\end {large}

\vspace{4mm}
%end
\newpage
%begin
%begin
\begin {large}

\textbf{22.11.2024 Контрольна робота, варіант  \No { 82 }}\par


{
%beginitem
1. По яких днях тижня відбувались лекції з предмету?

%enditem


\vspace{5mm}
%beginitem
2. Як зовуть (ім'я, по батькові) викладача, що вів практичні заняття?

%enditem


\vspace{5mm}
%beginitem
3. Нехай int var=13; Один потік збільшив змінну var на 3, а інший збільшив її в три рази.

гонка даних (data race), стан гонки (race condition)

Виберіть усі правильні твердження (якщо такі є).\par
\vspace{2mm}
( 1 ) Тут може бути тільки стан гонки.%bad

( 2 ) Тут може бути тільки гонка даних.%bad

( 3 ) Тут може не бути ані гонки даних, ані стану гонки.%bad

( 4 ) Тут гарантовано є стан гонки, але гонки даних може не бути.%good
%enditem


\vspace{5mm}
%beginitem
4. Виберіть усі правильні твердження (якщо такі є).\par
\vspace{2mm}

( 1 ) У стані гонки програма відчайдушно намагається захопити ресурси процесора.%bad

( 2 ) Тільки від компілятора залежить, чи запустить функція async виконання в новому потоці.%bad

( 3 ) Секція є критичною, якщо під час її виконання не можна генерувати винятки.%bad

( 4 ) Щоб розблокувати заблокований \code{shared\_mutex}, його завжди достатньо розблокувати один раз.%good
%enditem


\vspace{5mm}
%beginitem
5. Виберіть усі правильні твердження (якщо такі є).\par
\vspace{2mm}

( 1 ) Структура даних без блокувань гарантує, що жоден потік не буде очікувати доступу до виконання операцій над цією структурою даних.%bad

( 2 ) Операція над вільною від блокувань структурою даних завжди виконається за обмежений час.%bad

( 3 ) Вільних від очікувань структур даних не існує.%bad

( 4 ) Реалізація структури даних з блокуваннями може використовувати атомарні змінні.%good
\newpage

%enditem


\vspace{5mm}
%beginitem
6. Виберіть усі правильні твердження (якщо такі є).\par
\vspace{2mm}

( 1 ) Динамічні змінні не можуть розділятися між потоками.%bad

( 2 ) Потік може мати доступ до автоматичних змінних іншого потоку, але це небезпечно.%good

( 3 ) Потік мови С++ не може отримати доступ до автоматичних змінних іншого потоку.%bad

( 4 ) Потік має доступ тільки до своїх власних даних та до глобальних статичних даних.%bad
%enditem


\vspace{5mm}
%beginitem
7. У цьому фрагменті коду
\begin{verbatim}
template <class T> void f(T a1, T a2){}

void g(){
    int i=5;
    f(unique_ptr<int>(new int(++i)), unique_ptr<int>(new int(++i)));
\end{verbatim}

( 1 ) є невизначена поведінка%good

( 2 ) можливий виток ресурсів%bad

( 3 ) є гонка даних%bad

( 4 ) є стан гонки%bad
%enditem
}
%end
\end {large}

\vspace{4mm}
%end
\newpage
%begin
%begin
\begin {large}

\textbf{22.11.2024 Контрольна робота, варіант  \No { 83 }}\par


{
%beginitem
1. По яких днях тижня відбувались лекції з предмету?

%enditem


\vspace{5mm}
%beginitem
2. Як зовуть (ім'я, по батькові) викладача, що вів практичні заняття?

%enditem


\vspace{5mm}
%beginitem
3. Нехай int var=13; Один потік збільшив змінну var на 3, а інший збільшив її в три рази.

гонка даних (data race), стан гонки (race condition)

Виберіть усі правильні твердження (якщо такі є).\par
\vspace{2mm}
( 1 ) Тут може не бути стану гонки.%bad

( 2 ) Тут гарантовано нема стану гонки, але може бути гонка даних.%bad

( 3 ) Тут гарантовано є і стан гонки, і гонка даних.%bad

( 4 ) Тут може не бути гонки даних.%good
%enditem


\vspace{5mm}
%beginitem
4. Виберіть усі правильні твердження (якщо такі є).\par
\vspace{2mm}

( 1 ) Функція async запускає виконання в новому потоці.%bad

( 2 ) Критичну секцію можна реалізувати за допомогою м’ютексу.%good

( 3 ) Не можна блокувати заблокований м’ютекс mutex.%bad

( 4 ) Об’єкт promise можна скопіювати та роздати копії різним потокам, щоб кожен з них міг передати результат своєї роботи.%bad
%enditem


\vspace{5mm}
%beginitem
5. Виберіть усі правильні твердження (якщо такі є).\par
\vspace{2mm}

( 1 ) Структура даних без блокувань зобов’язана використовувати атомарні змінні.%bad

( 2 ) На операціях структури даних без блокувань взаємоблокувань потоків виникнути не може.%bad

( 3 ) Структура даних без блокувань гарантує завершення виконання довільної операції над цією структурою даних за обмежений час.%bad

( 4 ) Реалізація структури даних без блокувань не може використовувати м’ютекси.%good
\newpage

%enditem


\vspace{5mm}
%beginitem
6. Виберіть усі правильні твердження (якщо такі є).\par
\vspace{2mm}

( 1 ) Потік мови С++ не може отримати доступ до автоматичних змінних іншого потоку.%bad

( 2 ) Кожен потік програми може мати доступ до динамічних даних програми.%good

( 3 ) Динамічні змінні не можуть розділятися між потоками.%bad

( 4 ) Потік має доступ тільки до своїх власних даних та до глобальних статичних даних.%bad
%enditem


\vspace{5mm}
%beginitem
7. У цьому фрагменті коду
\begin{verbatim}
template <class T> void f(T a){}

void g(){
    int i=5;
    f(unique_ptr<int>(new int(i)))+(unique_ptr<int>(new int(35)));
\end{verbatim}
( 1 ) є стан гонки%bad

( 2 ) можливий виток ресурсів%good

( 3 ) є гонка даних%bad

( 4 ) є невизначена поведінка%bad
%enditem
}
%end
\end {large}

\vspace{4mm}
%end
\newpage
%begin
%begin
\begin {large}

\textbf{22.11.2024 Контрольна робота, варіант  \No { 84 }}\par


{
%beginitem
1. По яких днях тижня відбувались лекції з предмету?

%enditem


\vspace{5mm}
%beginitem
2. Як зовуть (ім'я, по батькові) викладача, що вів практичні заняття?

%enditem


\vspace{5mm}
%beginitem
3. Нехай int var=13; Один потік збільшив змінну var на 3, а інший збільшив її в три рази.

гонка даних (data race), стан гонки (race condition)

Виберіть усі правильні твердження (якщо такі є).\par
\vspace{2mm}
( 1 ) Тут гарантовано нема гонки даних, але може бути стан гонки.%bad

( 2 ) Тут гарантовано є гонка даних.%bad

( 3 ) Тут гарантовано нема ані гонки даних, ані стану гонки.%bad

( 4 ) Тут гарантовано є стан гонки.%good
%enditem


\vspace{5mm}
%beginitem
4. Виберіть усі правильні твердження (якщо такі є).\par
\vspace{2mm}

( 1 ) Критична секція може складатися з кількох ділянок коду, які можуть розташовуватись де завгодно (не обов’язково зв’язний фрагмент коду).%good

( 2 ) Щоб розблокувати заблокований м’ютекс \code{recursive\_mutex}, його завжди достатньо розблокувати один раз.%bad

( 3 ) Критична секція це неперервна ділянка коду програми.%bad

( 4 ) М’ютекс має бути записаний на початку критичної секції, яку він забезпечує.%bad
%enditem


\vspace{5mm}
%beginitem
5. Виберіть усі правильні твердження (якщо такі є).\par
\vspace{2mm}

( 1 ) Стандартні контейнери stl безпечно використовувати з кількох потоків.%bad

( 2 ) Мова С++20 має стандартну бібліотеку вільних від блокувань структур даних.%bad

( 3 ) Реалізації різних операцій над структурою даних повинні використовувати різні м’ютекси.%bad

( 4 ) Структура даних з блокуваннями може бути такою, що використовує атомарні змінні.%good
\newpage

%enditem


\vspace{5mm}
%beginitem
6. Виберіть усі правильні твердження (якщо такі є).\par
\vspace{2mm}

( 1 ) Потік може мати доступ до автоматичних змінних іншого потоку і це безпечно.%bad

( 2 ) Потік мови С++ не може отримати доступ до даних іншого потоку.%bad

( 3 ) Потік може мати доступ до автоматичних змінних іншого потоку, але це небезпечно.%good

( 4 ) Автоматичні змінні не можуть розділятися між потоками.%bad
%enditem


\vspace{5mm}
%beginitem
7. У цьому фрагменті коду
\begin{verbatim}
template <class T> void f(T a1, T a2){}

void g(){
    int i=5;
    f(unique_ptr<int>(new int(i++)), unique_ptr<int>(new int(++i)));
\end{verbatim}

( 1 ) можливий виток ресурсів%bad

( 2 ) є стан гонки%bad

( 3 ) є гонка даних%bad

( 4 ) є невизначена поведінка%good
%enditem
}
%end
\end {large}

\vspace{4mm}
%end
\newpage
%begin
%begin
\begin {large}

\textbf{22.11.2024 Контрольна робота, варіант  \No { 85 }}\par


{
%beginitem
1. По яких днях тижня відбувались лекції з предмету?

%enditem


\vspace{5mm}
%beginitem
2. Як зовуть (ім'я, по батькові) викладача, що вів практичні заняття?

%enditem


\vspace{5mm}
%beginitem
3. Нехай int var=13; Один потік збільшив змінну var на 3, а інший збільшив її в три рази.

гонка даних (data race), стан гонки (race condition)

Виберіть усі правильні твердження (якщо такі є).\par
\vspace{2mm}
( 1 ) Тут може не бути стану гонки.%bad

( 2 ) Тут гарантовано нема стану гонки, але може бути гонка даних.%bad

( 3 ) Тут може не бути ані гонки даних, ані стану гонки.%bad

( 4 ) Тут гарантовано є стан гонки, але гонки даних може не бути.%good
%enditem


\vspace{5mm}
%beginitem
4. Виберіть усі правильні твердження (якщо такі є).\par
\vspace{2mm}

( 1 ) Довільний потік може розблокувати заблокований м’ютекс.%bad

( 2 ) Блокувати заблокований м’ютекс  \code{mutex} іноді можна.%good

( 3 ) Об’єкти класів promise та future можна використовувати тільки в багатопотокових програмах.%bad

( 4 ) Тільки той потік, який обнулив значення семафора, може його збільшити.%bad
%enditem


\vspace{5mm}
%beginitem
5. Виберіть усі правильні твердження (якщо такі є).\par
\vspace{2mm}

( 1 ) Структура даних без блокувань гарантує, що жоден потік не буде очікувати доступу до виконання операцій над цією структурою даних.%bad

( 2 ) Якщо весь код структури даних утворює одну критичну секцію, то така структура даних може використовуватись в конкурентній програмі, але по-справжньо\-му конкурентною не є.%good

( 3 ) Одночасно використовувати атомарних змінних більше, ніж регістрів процесора, у програмі не можна.%bad

( 4 ) Операція над структурою даних без блокувань завжди виконається за обмежений час.%bad
\newpage

%enditem


\vspace{5mm}
%beginitem
6. Виберіть усі правильні твердження (якщо такі є).\par
\vspace{2mm}

( 1 ) Потік мови С++ не може отримати доступ до автоматичних змінних іншого потоку.%bad

( 2 ) Потік може мати доступ до автоматичних змінних іншого потоку і це безпечно.%bad

( 3 ) Потік може мати доступ до автоматичних змінних іншого потоку, але це небезпечно.%good

( 4 ) Динамічні змінні не можуть розділятися між потоками.%bad
%enditem


\vspace{5mm}
%beginitem
7. У цьому фрагменті коду
\begin{verbatim}
template <class T> void f(T a1, T a2){}

void g(){
    f(unique_ptr<int>(new int(3)), unique_ptr<int>(new int(35)));
\end{verbatim}

( 1 ) можливий виток ресурсів%bad

( 2 ) є гонка даних%bad

( 3 ) є невизначена поведінка%bad

( 4 ) є стан гонки%bad
%enditem
}
%end
\end {large}

\vspace{4mm}
%end
\newpage
%begin
%begin
\begin {large}

\textbf{22.11.2024 Контрольна робота, варіант  \No { 86 }}\par


{
%beginitem
1. По яких днях тижня відбувались лекції з предмету?

%enditem


\vspace{5mm}
%beginitem
2. Як зовуть (ім'я, по батькові) викладача, що вів практичні заняття?

%enditem


\vspace{5mm}
%beginitem
3. Нехай int var=13; Один потік збільшив змінну var на 3, а інший збільшив її в три рази.

гонка даних (data race), стан гонки (race condition)

Виберіть усі правильні твердження (якщо такі є).\par
\vspace{2mm}
( 1 ) Тут гарантовано є стан гонки.%good

( 2 ) Тут може бути тільки гонка даних.%bad

( 3 ) Тут гарантовано є і стан гонки, і гонка даних.%bad

( 4 ) Тут гарантовано є гонка даних.%bad
%enditem


\vspace{5mm}
%beginitem
4. Виберіть усі правильні твердження (якщо такі є).\par
\vspace{2mm}

( 1 ) Виклик \code{this\_thread::sleep\_for(5s)} змушує потік проспати рівно 5 секунд.%bad

( 2 ) У синтаксично коректній програмі мовою С++ стан гонки (race condition) може виникати.%good

( 3 ) Функція \code{sleep\_for} може змусити інший потік заснути.%bad

( 4 ) Розблоковувати незаблокований м’ютекс можна.%bad
%enditem


\vspace{5mm}
%beginitem
5. Виберіть усі правильні твердження (якщо такі є).\par
\vspace{2mm}

( 1 ) Взаємоблокування потоків може виникнути тільки за використання м’ютексів.%bad

( 2 ) Вільна від блокувань структура даних гарантує, що жоден потік не буде очікувати доступу до виконання операцій над цією структурою даних.%bad

( 3 ) Структура даних без блокувань зобов’язана використовувати атомарні змінні.%bad

( 4 ) Кожна вільна від очікувань структура даних є структурою даних, вільною від блокувань.%good
\newpage

%enditem


\vspace{5mm}
%beginitem
6. Виберіть усі правильні твердження (якщо такі є).\par
\vspace{2mm}

( 1 ) Потік мови С++ не може отримати доступ до даних іншого потоку.%bad

( 2 ) Автоматичні змінні не можуть розділятися між потоками.%bad

( 3 ) Потік має доступ тільки до своїх власних даних та до глобальних статичних даних.%bad

( 4 ) Кожен потік програми може мати доступ до динамічних даних програми.%good
%enditem


\vspace{5mm}
%beginitem
7. У цьому фрагменті коду
\begin{verbatim}
template <class T> void f(T a1, T a2){}

void g(){
    int i=5;
    f(unique_ptr<int>(new int(++i)), unique_ptr<int>(new int(++i)));
\end{verbatim}

( 1 ) є гонка даних%bad

( 2 ) є стан гонки%bad

( 3 ) можливий виток ресурсів%bad

( 4 ) є невизначена поведінка%good
%enditem
}
%end
\end {large}

\vspace{4mm}
%end
\newpage
%begin
%begin
\begin {large}

\textbf{22.11.2024 Контрольна робота, варіант  \No { 87 }}\par


{
%beginitem
1. По яких днях тижня відбувались лекції з предмету?

%enditem


\vspace{5mm}
%beginitem
2. Як зовуть (ім'я, по батькові) викладача, що вів практичні заняття?

%enditem


\vspace{5mm}
%beginitem
3. Нехай int var=13; Один потік збільшив змінну var на 3, а інший збільшив її в три рази.

гонка даних (data race), стан гонки (race condition)

Виберіть усі правильні твердження (якщо такі є).\par
\vspace{2mm}
( 1 ) Тут може бути тільки стан гонки.%bad

( 2 ) Тут гарантовано нема ані гонки даних, ані стану гонки.%bad

( 3 ) Тут може не бути гонки даних.%good

( 4 ) Тут гарантовано нема гонки даних, але може бути стан гонки.%bad
%enditem


\vspace{5mm}
%beginitem
4. Виберіть усі правильні твердження (якщо такі є).\par
\vspace{2mm}

( 1 ) Клас thread має засоби, щоб повернути результат виконання запущеного потоку.%bad

( 2 ) Потік в активному очікуванні споживає ресурс процесора.%good

( 3 ) Значення, що є результатом виконання функції, запущеної через async, зберігається безпосередньо в об’єкті проміза (promise).%bad

( 4 ) Значення, що є результатом виконання функції, запущеної через async, зберігається безпосередньо в об’єкті ф’ючерса (future).%bad
%enditem


\vspace{5mm}
%beginitem
5. Виберіть усі правильні твердження (якщо такі є).\par
\vspace{2mm}

( 1 ) На операціях структури даних без блокувань взаємоблокувань потоків виникнути не може.%bad

( 2 ) Структура даних без блокувань може не бути вільною від блокувань.%good

( 3 ) Структура даних без блокувань гарантує завершення виконання довільної операції над цією структурою даних за обмежений час.%bad

( 4 ) За використання атомарних змінних структура даних є структурою даних без блокувань.%bad
\newpage

%enditem


\vspace{5mm}
%beginitem
6. Виберіть усі правильні твердження (якщо такі є).\par
\vspace{2mm}

( 1 ) Потік мови С++ не може отримати доступ до автоматичних змінних іншого потоку.%bad

( 2 ) Кожен потік програми може мати доступ до динамічних даних програми.%good

( 3 ) Динамічні змінні не можуть розділятися між потоками.%bad

( 4 ) Потік може мати доступ до автоматичних змінних іншого потоку і це безпечно.%bad
%enditem


\vspace{5mm}
%beginitem
7. У цьому фрагменті коду
\begin{verbatim}
template <class T> void f(T a){}

void g(){
    int i=5;
    f(unique_ptr<int>(new int(i)))+(unique_ptr<int>(new int(++i)));
\end{verbatim}
( 1 ) є стан гонки%bad

( 2 ) можливий виток ресурсів%good

( 3 ) є невизначена поведінка%good

( 4 ) є гонка даних%bad
%enditem
}
%end
\end {large}

\vspace{4mm}
%end
\newpage
%begin
%begin
\begin {large}

\textbf{22.11.2024 Контрольна робота, варіант  \No { 88 }}\par


{
%beginitem
1. По яких днях тижня відбувались лекції з предмету?

%enditem


\vspace{5mm}
%beginitem
2. Як зовуть (ім'я, по батькові) викладача, що вів практичні заняття?

%enditem


\vspace{5mm}
%beginitem
3. Нехай int var=13; Один потік збільшив змінну var на 3, а інший збільшив її в три рази.

гонка даних (data race), стан гонки (race condition)

Виберіть усі правильні твердження (якщо такі є).\par
\vspace{2mm}
( 1 ) Тут гарантовано є і стан гонки, і гонка даних.%bad

( 2 ) Тут гарантовано є стан гонки.%good

( 3 ) Тут гарантовано є гонка даних.%bad

( 4 ) Тут може не бути стану гонки.%bad
%enditem


\vspace{5mm}
%beginitem
4. Виберіть усі правильні твердження (якщо такі є).\par
\vspace{2mm}

( 1 ) Потік, запущений конструктором класу jthread, має обов’язково зупинитися, якщо пов’язаний об’єкт надішле команду зупинитися.%bad

( 2 ) Досвічений програміст за однопотоковою програмою (текстом мовою С++) завжди може вказати послідовність виконуваних обчислень на конкретних вхідних даних.%bad

( 3 ) Коли програма спить, то вона не виконується.%bad

( 4 ) Якщо два обчислення програми на С++ тільки читають одну ту саму ділянку пам'яті, то на них гонка даних  виникнути не може.%good
%enditem


\vspace{5mm}
%beginitem
5. Виберіть усі правильні твердження (якщо такі є).\par
\vspace{2mm}

( 1 ) Мова С++20 має стандартну бібліотеку конкурентних структур даних.%bad

( 2 ) Для уникнення гонки даних структура даних може використовувати не більше одного м’ютекса.%bad

( 3 ) За використання структурою даних рівно одного м’ютексу, взаємоблокування на її операціях неможливі.%good

( 4 ) Реалізації різних операцій над структурою даних повинні використовувати різні м’ютекси.%bad
\newpage

%enditem


\vspace{5mm}
%beginitem
6. Виберіть усі правильні твердження (якщо такі є).\par
\vspace{2mm}

( 1 ) Потік має доступ тільки до своїх власних даних та до глобальних статичних даних.%bad

( 2 ) Потік може мати доступ до автоматичних змінних іншого потоку, але це небезпечно.%good

( 3 ) Автоматичні змінні не можуть розділятися між потоками.%bad

( 4 ) Потік мови С++ не може отримати доступ до даних іншого потоку.%bad
%enditem


\vspace{5mm}
%beginitem
7. У цьому фрагменті коду
\begin{verbatim}
template <class T> void f(T a1, T a2){}

void g(){
    int i=5;
    f(unique_ptr<int>(new int(++i)), unique_ptr<int>(new int(++i)));
\end{verbatim}

( 1 ) є гонка даних%bad

( 2 ) можливий виток ресурсів%bad

( 3 ) є стан гонки%bad

( 4 ) є невизначена поведінка%good
%enditem
}
%end
\end {large}

\vspace{4mm}
%end
\newpage
%begin
%begin
\begin {large}

\textbf{22.11.2024 Контрольна робота, варіант  \No { 89 }}\par


{
%beginitem
1. По яких днях тижня відбувались лекції з предмету?

%enditem


\vspace{5mm}
%beginitem
2. Як зовуть (ім'я, по батькові) викладача, що вів практичні заняття?

%enditem


\vspace{5mm}
%beginitem
3. Нехай int var=13; Один потік збільшив змінну var на 3, а інший збільшив її в три рази.

гонка даних (data race), стан гонки (race condition)

Виберіть усі правильні твердження (якщо такі є).\par
\vspace{2mm}
( 1 ) Тут може бути тільки стан гонки.%bad

( 2 ) Тут може не бути гонки даних.%good

( 3 ) Тут може не бути ані гонки даних, ані стану гонки.%bad

( 4 ) Тут гарантовано нема стану гонки, але може бути гонка даних.%bad
%enditem


\vspace{5mm}
%beginitem
4. Виберіть усі правильні твердження (якщо такі є).\par
\vspace{2mm}

( 1 ) У програмі може бути не більше однієї критичної секції.%bad

( 2 ) Функція async запускає виконання в новому потоці.%bad

( 3 ) Порядок обчислень у потоці для одних тих самих вхідних даних може залежати від використовуваного компілятора.%good

( 4 ) Значення, що є результатом виконання функції, запущеної через async, зберігається безпосередньо в об’єкті проміза (promise).%bad
%enditem


\vspace{5mm}
%beginitem
5. Виберіть усі правильні твердження (якщо такі є).\par
\vspace{2mm}

( 1 ) Структура даних може використовувати м’ютекси та атомарні змінні одночасно.%good

( 2 ) Операція над вільною від блокувань структурою даних завжди виконається за обмежений час.%bad

( 3 ) Структура даних без блокувань та вільна від блокувань структура даних --- це одне те саме.%bad

( 4 ) Мова С++20 має стандартну бібліотеку вільних від блокувань структур даних.%bad
\newpage

%enditem


\vspace{5mm}
%beginitem
6. Виберіть усі правильні твердження (якщо такі є).\par
\vspace{2mm}

( 1 ) Потік мови С++ не може отримати доступ до даних іншого потоку.%bad

( 2 ) Кожен потік програми може мати доступ до динамічних даних програми.%good

( 3 ) Потік може мати доступ до автоматичних змінних іншого потоку і це безпечно.%bad

( 4 ) Потік має доступ тільки до своїх власних даних та до глобальних статичних даних.%bad
%enditem


\vspace{5mm}
%beginitem
7. У цьому фрагменті коду
\begin{verbatim}
template <class T> void f(T a1, T a2){}

void g(){
    int i=5;
    f(unique_ptr<int>(new int(i++)), unique_ptr<int>(new int(++i)));
\end{verbatim}

( 1 ) є гонка даних%bad

( 2 ) можливий виток ресурсів%bad

( 3 ) є стан гонки%bad

( 4 ) є невизначена поведінка%good
%enditem
}
%end
\end {large}

\vspace{4mm}
%end
\newpage
%begin
%begin
\begin {large}

\textbf{22.11.2024 Контрольна робота, варіант  \No { 90 }}\par


{
%beginitem
1. По яких днях тижня відбувались лекції з предмету?

%enditem


\vspace{5mm}
%beginitem
2. Як зовуть (ім'я, по батькові) викладача, що вів практичні заняття?

%enditem


\vspace{5mm}
%beginitem
3. Нехай int var=13; Один потік збільшив змінну var на 3, а інший збільшив її в три рази.

гонка даних (data race), стан гонки (race condition)

Виберіть усі правильні твердження (якщо такі є).\par
\vspace{2mm}
( 1 ) Тут гарантовано нема гонки даних, але може бути стан гонки.%bad

( 2 ) Тут гарантовано нема ані гонки даних, ані стану гонки.%bad

( 3 ) Тут гарантовано є стан гонки, але гонки даних може не бути.%good

( 4 ) Тут може бути тільки гонка даних.%bad
%enditem


\vspace{5mm}
%beginitem
4. Виберіть усі правильні твердження (якщо такі є).\par
\vspace{2mm}

( 1 ) Щоб розблокувати заблокований м’ютекс \code{recursive\_mutex}, його завжди достатньо розблокувати один раз.%bad

( 2 ) Не можна блокувати заблокований м’ютекс mutex.%bad

( 3 ) Об’єкт promise можна скопіювати та роздати копії різним потокам, щоб кожен з них міг передати результат своєї роботи.%bad

( 4 ) Розблоковувати незаблокований м’ютекс не можна.%good
%enditem


\vspace{5mm}
%beginitem
5. Виберіть усі правильні твердження (якщо такі є).\par
\vspace{2mm}

( 1 ) Вільних від очікувань структур даних не існує.%bad

( 2 ) Операція над вільною від блокувань структурою даних може виконуватись скільки завгодно довго.%good

( 3 ) Стандартні контейнери stl безпечно використовувати з кількох потоків.%bad

( 4 ) Структура даних є структурою даних без блокувань, якщо її реалізація не використовує м’ютекси.%bad
\newpage

%enditem


\vspace{5mm}
%beginitem
6. Виберіть усі правильні твердження (якщо такі є).\par
\vspace{2mm}

( 1 ) Потік може мати доступ до автоматичних змінних іншого потоку, але це небезпечно.%good

( 2 ) Автоматичні змінні не можуть розділятися між потоками.%bad

( 3 ) Динамічні змінні не можуть розділятися між потоками.%bad

( 4 ) Потік мови С++ не може отримати доступ до автоматичних змінних іншого потоку.%bad
%enditem


\vspace{5mm}
%beginitem
7. У цьому фрагменті коду
\begin{verbatim}
template <class T> void f(T a){}

void g(){
    f(unique_ptr<int>(new int(3)))+(unique_ptr<int>(new int(35)));
\end{verbatim}

( 1 ) є стан гонки%bad

( 2 ) є гонка даних%bad

( 3 ) можливий виток ресурсів%good

( 4 ) є невизначена поведінка%bad
%enditem
}
%end
\end {large}

\vspace{4mm}
%end
\newpage
%begin
%begin
\begin {large}

\textbf{22.11.2024 Контрольна робота, варіант  \No { 91 }}\par


{
%beginitem
1. По яких днях тижня відбувались лекції з предмету?

%enditem


\vspace{5mm}
%beginitem
2. Як зовуть (ім'я, по батькові) викладача, що вів практичні заняття?

%enditem


\vspace{5mm}
%beginitem
3. Нехай int var=13; Один потік збільшив змінну var на 3, а інший збільшив її в три рази.

гонка даних (data race), стан гонки (race condition)

Виберіть усі правильні твердження (якщо такі є).\par
\vspace{2mm}
( 1 ) Тут гарантовано нема ані гонки даних, ані стану гонки.%bad

( 2 ) Тут гарантовано нема стану гонки, але може бути гонка даних.%bad

( 3 ) Тут гарантовано є стан гонки, але гонки даних може не бути.%good

( 4 ) Тут може бути тільки стан гонки.%bad
%enditem


\vspace{5mm}
%beginitem
4. Виберіть усі правильні твердження (якщо такі є).\par
\vspace{2mm}

( 1 ) Критична секція це неперервна ділянка коду програми.%bad

( 2 ) М’ютекс має бути записаний на початку критичної секції, яку він забезпечує.%bad

( 3 ) Потік, запущений конструктором класу jthread, має обов’язково зупинитися, якщо пов’язаний об’єкт надішле команду зупинитися.%bad

( 4 ) Щоб розблокувати заблокований \code{shared\_mutex}, його завжди достатньо розблокувати один раз.%good
%enditem


\vspace{5mm}
%beginitem
5. Виберіть усі правильні твердження (якщо такі є).\par
\vspace{2mm}

( 1 ) Мова С++20 має стандартну бібліотеку вільних від блокувань структур даних.%bad

( 2 ) Реалізації різних операцій над структурою даних повинні використовувати різні м’ютекси.%bad

( 3 ) Вільна від блокувань структура даних гарантує, що жоден потік не буде очікувати доступу до виконання операцій над цією структурою даних.%bad

( 4 ) За використання структурою даних рівно одного м’ютексу, взаємоблокування на її операціях неможливі.%good
\newpage

%enditem


\vspace{5mm}
%beginitem
6. Виберіть усі правильні твердження (якщо такі є).\par
\vspace{2mm}

( 1 ) Потік може мати доступ до автоматичних змінних іншого потоку і це безпечно.%bad

( 2 ) Потік мови С++ не може отримати доступ до автоматичних змінних іншого потоку.%bad

( 3 ) Потік може мати доступ до автоматичних змінних іншого потоку, але це небезпечно.%good

( 4 ) Потік має доступ тільки до своїх власних даних та до глобальних статичних даних.%bad
%enditem


\vspace{5mm}
%beginitem
7. У цьому фрагменті коду
\begin{verbatim}
template <class T> void f(T a){}

void g(){
    int i=5;
    f(unique_ptr<int>(new int(i)))+(unique_ptr<int>(new int(35)));
\end{verbatim}
( 1 ) є стан гонки%bad

( 2 ) є невизначена поведінка%bad

( 3 ) є гонка даних%bad

( 4 ) можливий виток ресурсів%good
%enditem
}
%end
\end {large}

\vspace{4mm}
%end
\newpage
%begin
%begin
\begin {large}

\textbf{22.11.2024 Контрольна робота, варіант  \No { 92 }}\par


{
%beginitem
1. По яких днях тижня відбувались лекції з предмету?

%enditem


\vspace{5mm}
%beginitem
2. Як зовуть (ім'я, по батькові) викладача, що вів практичні заняття?

%enditem


\vspace{5mm}
%beginitem
3. Нехай int var=13; Один потік збільшив змінну var на 3, а інший збільшив її в три рази.

гонка даних (data race), стан гонки (race condition)

Виберіть усі правильні твердження (якщо такі є).\par
\vspace{2mm}
( 1 ) Тут може не бути стану гонки.%bad

( 2 ) Тут може бути тільки гонка даних.%bad

( 3 ) Тут гарантовано є і стан гонки, і гонка даних.%bad

( 4 ) Тут гарантовано є стан гонки.%good
%enditem


\vspace{5mm}
%beginitem
4. Виберіть усі правильні твердження (якщо такі є).\par
\vspace{2mm}

( 1 ) У синтаксично коректній програмі мовою С++ стан гонки (race condition) може виникати.%good

( 2 ) Тільки від компілятора залежить, чи запустить функція async виконання в новому потоці.%bad

( 3 ) Об’єкт future можна скопіювати та роздати копії різним потокам, щоб вони змогли побачити результат.%bad

( 4 ) У стані гонки програма відчайдушно намагається захопити ресурси процесора.%bad
%enditem


\vspace{5mm}
%beginitem
5. Виберіть усі правильні твердження (якщо такі є).\par
\vspace{2mm}

( 1 ) Структура даних без блокувань та вільна від блокувань структура даних --- це одне те саме.%bad

( 2 ) Структура даних без блокувань гарантує завершення виконання довільної операції над цією структурою даних за обмежений час.%bad

( 3 ) На операціях структури даних без блокувань взаємоблокувань потоків виникнути не може.%bad

( 4 ) Реалізація структури даних без блокувань не може використовувати м’ютекси.%good
\newpage

%enditem


\vspace{5mm}
%beginitem
6. Виберіть усі правильні твердження (якщо такі є).\par
\vspace{2mm}

( 1 ) Автоматичні змінні не можуть розділятися між потоками.%bad

( 2 ) Потік мови С++ не може отримати доступ до даних іншого потоку.%bad

( 3 ) Динамічні змінні не можуть розділятися між потоками.%bad

( 4 ) Кожен потік програми може мати доступ до динамічних даних програми.%good
%enditem


\vspace{5mm}
%beginitem
7. У цьому фрагменті коду
\begin{verbatim}
template <class T> void f(T a){}

void g(){
    int i=5;
    f(unique_ptr<int>(new int(i)))+(unique_ptr<int>(new int(35)));
\end{verbatim}
( 1 ) є гонка даних%bad

( 2 ) можливий виток ресурсів%good

( 3 ) є невизначена поведінка%bad

( 4 ) є стан гонки%bad
%enditem
}
%end
\end {large}

\vspace{4mm}
%end
\newpage
%begin
%begin
\begin {large}

\textbf{22.11.2024 Контрольна робота, варіант  \No { 93 }}\par


{
%beginitem
1. По яких днях тижня відбувались лекції з предмету?

%enditem


\vspace{5mm}
%beginitem
2. Як зовуть (ім'я, по батькові) викладача, що вів практичні заняття?

%enditem


\vspace{5mm}
%beginitem
3. Нехай int var=13; Один потік збільшив змінну var на 3, а інший збільшив її в три рази.

гонка даних (data race), стан гонки (race condition)

Виберіть усі правильні твердження (якщо такі є).\par
\vspace{2mm}
( 1 ) Тут гарантовано є гонка даних.%bad

( 2 ) Тут може не бути ані гонки даних, ані стану гонки.%bad

( 3 ) Тут може не бути гонки даних.%good

( 4 ) Тут гарантовано нема гонки даних, але може бути стан гонки.%bad
%enditem


\vspace{5mm}
%beginitem
4. Виберіть усі правильні твердження (якщо такі є).\par
\vspace{2mm}

( 1 ) Секція є критичною, якщо під час її виконання не можна генерувати винятки.%bad

( 2 ) За різних запусків програми на одних тих самих вхідних даних може виникати хронологічно різна послідовність обчислень.%good

( 3 ) Значення, що є результатом виконання функції, запущеної через async, зберігається безпосередньо в об’єкті ф’ючерса (future).%bad

( 4 ) Тільки той потік, який обнулив значення семафора, може його збільшити.%bad
%enditem


\vspace{5mm}
%beginitem
5. Виберіть усі правильні твердження (якщо такі є).\par
\vspace{2mm}

( 1 ) Операція над структурою даних без блокувань завжди виконається за обмежений час.%bad

( 2 ) Якщо весь код структури даних утворює одну критичну секцію, то така структура даних може використовуватись в конкурентній програмі, але по-справжньо\-му конкурентною не є.%good

( 3 ) Одночасно використовувати атомарних змінних більше, ніж регістрів процесора, у програмі не можна.%bad

( 4 ) Операція над вільною від блокувань структурою даних завжди виконається за обмежений час.%bad
\newpage

%enditem


\vspace{5mm}
%beginitem
6. Виберіть усі правильні твердження (якщо такі є).\par
\vspace{2mm}

( 1 ) Потік мови С++ не може отримати доступ до даних іншого потоку.%bad

( 2 ) Потік мови С++ не може отримати доступ до автоматичних змінних іншого потоку.%bad

( 3 ) Динамічні змінні не можуть розділятися між потоками.%bad

( 4 ) Потік може мати доступ до автоматичних змінних іншого потоку, але це небезпечно.%good
%enditem


\vspace{5mm}
%beginitem
7. У цьому фрагменті коду
\begin{verbatim}
template <class T> void f(T a1, T a2){}

void g(){
    int i=5;
    f(unique_ptr<int>(new int(++i)), unique_ptr<int>(new int(++i)));
\end{verbatim}

( 1 ) є невизначена поведінка%good

( 2 ) є стан гонки%bad

( 3 ) можливий виток ресурсів%bad

( 4 ) є гонка даних%bad
%enditem
}
%end
\end {large}

\vspace{4mm}
%end
\newpage
%begin
%begin
\begin {large}

\textbf{22.11.2024 Контрольна робота, варіант  \No { 94 }}\par


{
%beginitem
1. По яких днях тижня відбувались лекції з предмету?

%enditem


\vspace{5mm}
%beginitem
2. Як зовуть (ім'я, по батькові) викладача, що вів практичні заняття?

%enditem


\vspace{5mm}
%beginitem
3. Нехай int var=13; Один потік збільшив змінну var на 3, а інший збільшив її в три рази.

гонка даних (data race), стан гонки (race condition)

Виберіть усі правильні твердження (якщо такі є).\par
\vspace{2mm}
( 1 ) Тут може не бути стану гонки.%bad

( 2 ) Тут гарантовано є і стан гонки, і гонка даних.%bad

( 3 ) Тут може бути тільки гонка даних.%bad

( 4 ) Тут гарантовано є стан гонки.%good
%enditem


\vspace{5mm}
%beginitem
4. Виберіть усі правильні твердження (якщо такі є).\par
\vspace{2mm}

( 1 ) Чи запустить функція  async виконання в новому потоці можна задавати політикою.%good

( 2 ) Клас thread має засоби, щоб повернути результат виконання запущеного потоку.%bad

( 3 ) Розблоковувати незаблокований м’ютекс можна.%bad

( 4 ) Досвічений програміст за однопотоковою програмою (текстом мовою С++) завжди може вказати послідовність виконуваних обчислень на конкретних вхідних даних.%bad
%enditem


\vspace{5mm}
%beginitem
5. Виберіть усі правильні твердження (якщо такі є).\par
\vspace{2mm}

( 1 ) Структура даних є структурою даних без блокувань, якщо її реалізація не використовує м’ютекси.%bad

( 2 ) Кожна вільна від очікувань структура даних є структурою даних, вільною від блокувань.%good

( 3 ) Мова С++20 має стандартну бібліотеку конкурентних структур даних.%bad

( 4 ) Структура даних без блокувань гарантує, що жоден потік не буде очікувати доступу до виконання операцій над цією структурою даних.%bad
\newpage

%enditem


\vspace{5mm}
%beginitem
6. Виберіть усі правильні твердження (якщо такі є).\par
\vspace{2mm}

( 1 ) Потік має доступ тільки до своїх власних даних та до глобальних статичних даних.%bad

( 2 ) Автоматичні змінні не можуть розділятися між потоками.%bad

( 3 ) Потік може мати доступ до автоматичних змінних іншого потоку і це безпечно.%bad

( 4 ) Кожен потік програми може мати доступ до динамічних даних програми.%good
%enditem


\vspace{5mm}
%beginitem
7. У цьому фрагменті коду
\begin{verbatim}
template <class T> void f(T a1, T a2){}

void g(){
    int i=5;
    f(unique_ptr<int>(new int(i++)), unique_ptr<int>(new int(++i)));
\end{verbatim}

( 1 ) є гонка даних%bad

( 2 ) є невизначена поведінка%good

( 3 ) є стан гонки%bad

( 4 ) можливий виток ресурсів%bad
%enditem
}
%end
\end {large}

\vspace{4mm}
%end
\newpage
%begin
%begin
\begin {large}

\textbf{22.11.2024 Контрольна робота, варіант  \No { 95 }}\par


{
%beginitem
1. По яких днях тижня відбувались лекції з предмету?

%enditem


\vspace{5mm}
%beginitem
2. Як зовуть (ім'я, по батькові) викладача, що вів практичні заняття?

%enditem


\vspace{5mm}
%beginitem
3. Нехай int var=13; Один потік збільшив змінну var на 3, а інший збільшив її в три рази.

гонка даних (data race), стан гонки (race condition)

Виберіть усі правильні твердження (якщо такі є).\par
\vspace{2mm}
( 1 ) Тут гарантовано нема ані гонки даних, ані стану гонки.%bad

( 2 ) Тут гарантовано є гонка даних.%bad

( 3 ) Тут може не бути гонки даних.%good

( 4 ) Тут гарантовано нема стану гонки, але може бути гонка даних.%bad
%enditem


\vspace{5mm}
%beginitem
4. Виберіть усі правильні твердження (якщо такі є).\par
\vspace{2mm}

( 1 ) Коли потік спить, він не споживає ресурс процесора.%good

( 2 ) Текст однопотокової програми мовою С++ однозначно визначає послідовність виконуваних обчислень.%bad

( 3 ) Функція \code{sleep\_for} може змусити інший потік заснути.%bad

( 4 ) Об’єкти класів promise та future можна використовувати тільки в багатопотокових програмах.%bad
%enditem


\vspace{5mm}
%beginitem
5. Виберіть усі правильні твердження (якщо такі є).\par
\vspace{2mm}

( 1 ) Взаємоблокування потоків може виникнути тільки за використання м’ютексів.%bad

( 2 ) Реалізація структури даних з блокуваннями може використовувати атомарні змінні.%good

( 3 ) Для уникнення гонки даних структура даних може використовувати не більше одного м’ютекса.%bad

( 4 ) За використання атомарних змінних структура даних є структурою даних без блокувань.%bad
\newpage

%enditem


\vspace{5mm}
%beginitem
6. Виберіть усі правильні твердження (якщо такі є).\par
\vspace{2mm}

( 1 ) Динамічні змінні не можуть розділятися між потоками.%bad

( 2 ) Автоматичні змінні не можуть розділятися між потоками.%bad

( 3 ) Потік мови С++ не може отримати доступ до даних іншого потоку.%bad

( 4 ) Потік може мати доступ до автоматичних змінних іншого потоку, але це небезпечно.%good
%enditem


\vspace{5mm}
%beginitem
7. У цьому фрагменті коду
\begin{verbatim}
template <class T> void f(T a1, T a2){}

void g(){
    int i=5;
    f(unique_ptr<int>(new int(++i)), unique_ptr<int>(new int(++i)));
\end{verbatim}

( 1 ) є невизначена поведінка%good

( 2 ) можливий виток ресурсів%bad

( 3 ) є стан гонки%bad

( 4 ) є гонка даних%bad
%enditem
}
%end
\end {large}

\vspace{4mm}
%end
\newpage
%begin
%begin
\begin {large}

\textbf{22.11.2024 Контрольна робота, варіант  \No { 96 }}\par


{
%beginitem
1. По яких днях тижня відбувались лекції з предмету?

%enditem


\vspace{5mm}
%beginitem
2. Як зовуть (ім'я, по батькові) викладача, що вів практичні заняття?

%enditem


\vspace{5mm}
%beginitem
3. Нехай int var=13; Один потік збільшив змінну var на 3, а інший збільшив її в три рази.

гонка даних (data race), стан гонки (race condition)

Виберіть усі правильні твердження (якщо такі є).\par
\vspace{2mm}
( 1 ) Тут гарантовано нема гонки даних, але може бути стан гонки.%bad

( 2 ) Тут може не бути ані гонки даних, ані стану гонки.%bad

( 3 ) Тут може бути тільки стан гонки.%bad

( 4 ) Тут гарантовано є стан гонки, але гонки даних може не бути.%good
%enditem


\vspace{5mm}
%beginitem
4. Виберіть усі правильні твердження (якщо такі є).\par
\vspace{2mm}

( 1 ) Коли програма спить, то вона не виконується.%bad

( 2 ) Довільний потік може розблокувати заблокований м’ютекс.%bad

( 3 ) Якщо два обчислення програми на С++ тільки читають одну ту саму ділянку пам'яті, то на них гонка даних  виникнути не може.%good

( 4 ) Виклик \code{this\_thread::sleep\_for(5s)} змушує потік проспати рівно 5 секунд.%bad
%enditem


\vspace{5mm}
%beginitem
5. Виберіть усі правильні твердження (якщо такі є).\par
\vspace{2mm}

( 1 ) Вільних від очікувань структур даних не існує.%bad

( 2 ) Структура даних може використовувати м’ютекси та атомарні змінні одночасно.%good

( 3 ) Структура даних без блокувань зобов’язана використовувати атомарні змінні.%bad

( 4 ) Стандартні контейнери stl безпечно використовувати з кількох потоків.%bad
\newpage

%enditem


\vspace{5mm}
%beginitem
6. Виберіть усі правильні твердження (якщо такі є).\par
\vspace{2mm}

( 1 ) Потік мови С++ не може отримати доступ до автоматичних змінних іншого потоку.%bad

( 2 ) Потік може мати доступ до автоматичних змінних іншого потоку і це безпечно.%bad

( 3 ) Потік має доступ тільки до своїх власних даних та до глобальних статичних даних.%bad

( 4 ) Кожен потік програми може мати доступ до динамічних даних програми.%good
%enditem


\vspace{5mm}
%beginitem
7. У цьому фрагменті коду
\begin{verbatim}
template <class T> void f(T a1, T a2){}

void g(){
    f(unique_ptr<int>(new int(3)), unique_ptr<int>(new int(35)));
\end{verbatim}

( 1 ) можливий виток ресурсів%bad

( 2 ) є невизначена поведінка%bad

( 3 ) є гонка даних%bad

( 4 ) є стан гонки%bad
%enditem
}
%end
\end {large}

\vspace{4mm}
%end
\newpage
%begin
%begin
\begin {large}

\textbf{22.11.2024 Контрольна робота, варіант  \No { 97 }}\par


{
%beginitem
1. По яких днях тижня відбувались лекції з предмету?

%enditem


\vspace{5mm}
%beginitem
2. Як зовуть (ім'я, по батькові) викладача, що вів практичні заняття?

%enditem


\vspace{5mm}
%beginitem
3. Нехай int var=13; Один потік збільшив змінну var на 3, а інший збільшив її в три рази.

гонка даних (data race), стан гонки (race condition)

Виберіть усі правильні твердження (якщо такі є).\par
\vspace{2mm}
( 1 ) Тут гарантовано нема гонки даних, але може бути стан гонки.%bad

( 2 ) Тут гарантовано нема стану гонки, але може бути гонка даних.%bad

( 3 ) Тут гарантовано є стан гонки, але гонки даних може не бути.%good

( 4 ) Тут може не бути ані гонки даних, ані стану гонки.%bad
%enditem


\vspace{5mm}
%beginitem
4. Виберіть усі правильні твердження (якщо такі є).\par
\vspace{2mm}

( 1 ) Розблокувати м’ютекс може тільки той потік, що його заблокував.%good

( 2 ) Досвічений програміст за однопотоковою програмою (текстом мовою С++) завжди може вказати послідовність виконуваних обчислень на конкретних вхідних даних.%bad

( 3 ) Секція є критичною, якщо під час її виконання не можна генерувати винятки.%bad

( 4 ) Текст однопотокової програми мовою С++ однозначно визначає послідовність виконуваних обчислень.%bad
%enditem


\vspace{5mm}
%beginitem
5. Виберіть усі правильні твердження (якщо такі є).\par
\vspace{2mm}

( 1 ) Структура даних з блокуваннями може бути такою, що використовує атомарні змінні.%good

( 2 ) Структура даних є структурою даних без блокувань, якщо її реалізація не використовує м’ютекси.%bad

( 3 ) Взаємоблокування потоків може виникнути тільки за використання м’ютексів.%bad

( 4 ) Структура даних без блокувань гарантує завершення виконання довільної операції над цією структурою даних за обмежений час.%bad
\newpage

%enditem


\vspace{5mm}
%beginitem
6. Виберіть усі правильні твердження (якщо такі є).\par
\vspace{2mm}

( 1 ) Потік мови С++ не може отримати доступ до автоматичних змінних іншого потоку.%bad

( 2 ) Автоматичні змінні не можуть розділятися між потоками.%bad

( 3 ) Потік може мати доступ до автоматичних змінних іншого потоку, але це небезпечно.%good

( 4 ) Потік може мати доступ до автоматичних змінних іншого потоку і це безпечно.%bad
%enditem


\vspace{5mm}
%beginitem
7. У цьому фрагменті коду
\begin{verbatim}
template <class T> void f(T a){}

void g(){
    f(unique_ptr<int>(new int(3)))+(unique_ptr<int>(new int(35)));
\end{verbatim}

( 1 ) є гонка даних%bad

( 2 ) можливий виток ресурсів%good

( 3 ) є стан гонки%bad

( 4 ) є невизначена поведінка%bad
%enditem
}
%end
\end {large}

\vspace{4mm}
%end
\newpage
%begin
%begin
\begin {large}

\textbf{22.11.2024 Контрольна робота, варіант  \No { 98 }}\par


{
%beginitem
1. По яких днях тижня відбувались лекції з предмету?

%enditem


\vspace{5mm}
%beginitem
2. Як зовуть (ім'я, по батькові) викладача, що вів практичні заняття?

%enditem


\vspace{5mm}
%beginitem
3. Нехай int var=13; Один потік збільшив змінну var на 3, а інший збільшив її в три рази.

гонка даних (data race), стан гонки (race condition)

Виберіть усі правильні твердження (якщо такі є).\par
\vspace{2mm}
( 1 ) Тут може не бути гонки даних.%good

( 2 ) Тут може бути тільки стан гонки.%bad

( 3 ) Тут може не бути стану гонки.%bad

( 4 ) Тут гарантовано є гонка даних.%bad
%enditem


\vspace{5mm}
%beginitem
4. Виберіть усі правильні твердження (якщо такі є).\par
\vspace{2mm}

( 1 ) М’ютекс має бути записаний на початку критичної секції, яку він забезпечує.%bad

( 2 ) Критична секція це неперервна ділянка коду програми.%bad

( 3 ) Об’єкт future можна скопіювати та роздати копії різним потокам, щоб вони змогли побачити результат.%bad

( 4 ) Виклик функції \code{this\_thread::sleep\_for(5s)} змушує потік, з якого його здійснили, заснути якнайменш на 5 секунд.%good
%enditem


\vspace{5mm}
%beginitem
5. Виберіть усі правильні твердження (якщо такі є).\par
\vspace{2mm}

( 1 ) Структура даних без блокувань може не бути вільною від блокувань.%good

( 2 ) Структура даних без блокувань гарантує, що жоден потік не буде очікувати доступу до виконання операцій над цією структурою даних.%bad

( 3 ) Вільна від блокувань структура даних гарантує, що жоден потік не буде очікувати доступу до виконання операцій над цією структурою даних.%bad

( 4 ) За використання атомарних змінних структура даних є структурою даних без блокувань.%bad
\newpage

%enditem


\vspace{5mm}
%beginitem
6. Виберіть усі правильні твердження (якщо такі є).\par
\vspace{2mm}

( 1 ) Потік мови С++ не може отримати доступ до даних іншого потоку.%bad

( 2 ) Потік має доступ тільки до своїх власних даних та до глобальних статичних даних.%bad

( 3 ) Кожен потік програми може мати доступ до динамічних даних програми.%good

( 4 ) Динамічні змінні не можуть розділятися між потоками.%bad
%enditem


\vspace{5mm}
%beginitem
7. У цьому фрагменті коду
\begin{verbatim}
template <class T> void f(T a){}

void g(){
    int i=5;
    f(unique_ptr<int>(new int(i)))+(unique_ptr<int>(new int(++i)));
\end{verbatim}
( 1 ) є гонка даних%bad

( 2 ) можливий виток ресурсів%good

( 3 ) є стан гонки%bad

( 4 ) є невизначена поведінка%good
%enditem
}
%end
\end {large}

\vspace{4mm}
%end
\newpage
%begin
%begin
\begin {large}

\textbf{22.11.2024 Контрольна робота, варіант  \No { 99 }}\par


{
%beginitem
1. По яких днях тижня відбувались лекції з предмету?

%enditem


\vspace{5mm}
%beginitem
2. Як зовуть (ім'я, по батькові) викладача, що вів практичні заняття?

%enditem


\vspace{5mm}
%beginitem
3. Нехай int var=13; Один потік збільшив змінну var на 3, а інший збільшив її в три рази.

гонка даних (data race), стан гонки (race condition)

Виберіть усі правильні твердження (якщо такі є).\par
\vspace{2mm}
( 1 ) Тут може бути тільки гонка даних.%bad

( 2 ) Тут гарантовано є стан гонки.%good

( 3 ) Тут гарантовано нема ані гонки даних, ані стану гонки.%bad

( 4 ) Тут гарантовано є і стан гонки, і гонка даних.%bad
%enditem


\vspace{5mm}
%beginitem
4. Виберіть усі правильні твердження (якщо такі є).\par
\vspace{2mm}

( 1 ) Функція \code{sleep\_for} може змусити інший потік заснути.%bad

( 2 ) У стані гонки програма відчайдушно намагається захопити ресурси процесора.%bad

( 3 ) Клас thread не має засобів, щоб повернути результат виконання запущеного потоку.%good

( 4 ) Тільки від компілятора залежить, чи запустить функція async виконання в новому потоці.%bad
%enditem


\vspace{5mm}
%beginitem
5. Виберіть усі правильні твердження (якщо такі є).\par
\vspace{2mm}

( 1 ) Операція над вільною від блокувань структурою даних завжди виконається за обмежений час.%bad

( 2 ) Операція над вільною від блокувань структурою даних може виконуватись скільки завгодно довго.%good

( 3 ) Для уникнення гонки даних структура даних може використовувати не більше одного м’ютекса.%bad

( 4 ) Реалізації різних операцій над структурою даних повинні використовувати різні м’ютекси.%bad
\newpage

%enditem


\vspace{5mm}
%beginitem
6. Виберіть усі правильні твердження (якщо такі є).\par
\vspace{2mm}

( 1 ) Потік має доступ тільки до своїх власних даних та до глобальних статичних даних.%bad

( 2 ) Кожен потік програми може мати доступ до динамічних даних програми.%good

( 3 ) Автоматичні змінні не можуть розділятися між потоками.%bad

( 4 ) Потік може мати доступ до автоматичних змінних іншого потоку і це безпечно.%bad
%enditem


\vspace{5mm}
%beginitem
7. У цьому фрагменті коду
\begin{verbatim}
template <class T> void f(T a1, T a2){}

void g(){
    int i=5;
    f(unique_ptr<int>(new int(++i)), unique_ptr<int>(new int(++i)));
\end{verbatim}

( 1 ) є стан гонки%bad

( 2 ) можливий виток ресурсів%bad

( 3 ) є гонка даних%bad

( 4 ) є невизначена поведінка%good
%enditem
}
%end
\end {large}

\vspace{4mm}
%end
\newpage
%begin
%begin
\begin {large}

\textbf{22.11.2024 Контрольна робота, варіант  \No { 100 }}\par


{
%beginitem
1. По яких днях тижня відбувались лекції з предмету?

%enditem


\vspace{5mm}
%beginitem
2. Як зовуть (ім'я, по батькові) викладача, що вів практичні заняття?

%enditem


\vspace{5mm}
%beginitem
3. Нехай int var=13; Один потік збільшив змінну var на 3, а інший збільшив її в три рази.

гонка даних (data race), стан гонки (race condition)

Виберіть усі правильні твердження (якщо такі є).\par
\vspace{2mm}
( 1 ) Тут гарантовано є гонка даних.%bad

( 2 ) Тут може не бути ані гонки даних, ані стану гонки.%bad

( 3 ) Тут може не бути гонки даних.%good

( 4 ) Тут може не бути стану гонки.%bad
%enditem


\vspace{5mm}
%beginitem
4. Виберіть усі правильні твердження (якщо такі є).\par
\vspace{2mm}

( 1 ) У програмі може бути не більше однієї критичної секції.%bad

( 2 ) Об’єкти класів promise та future можна використовувати тільки в багатопотокових програмах.%bad

( 3 ) Значення, що є результатом виконання функції, запущеної через async, зберігається безпосередньо в об’єкті проміза (promise).%bad

( 4 ) Порядок обчислень у потоці для одних тих самих вхідних даних може залежати від використовуваного компілятора.%good
%enditem


\vspace{5mm}
%beginitem
5. Виберіть усі правильні твердження (якщо такі є).\par
\vspace{2mm}

( 1 ) Реалізація структури даних без блокувань не може використовувати м’ютекси.%good

( 2 ) Операція над структурою даних без блокувань завжди виконається за обмежений час.%bad

( 3 ) Мова С++20 має стандартну бібліотеку вільних від блокувань структур даних.%bad

( 4 ) Одночасно використовувати атомарних змінних більше, ніж регістрів процесора, у програмі не можна.%bad
\newpage

%enditem


\vspace{5mm}
%beginitem
6. Виберіть усі правильні твердження (якщо такі є).\par
\vspace{2mm}

( 1 ) Потік мови С++ не може отримати доступ до даних іншого потоку.%bad

( 2 ) Динамічні змінні не можуть розділятися між потоками.%bad

( 3 ) Потік мови С++ не може отримати доступ до автоматичних змінних іншого потоку.%bad

( 4 ) Потік може мати доступ до автоматичних змінних іншого потоку, але це небезпечно.%good
%enditem


\vspace{5mm}
%beginitem
7. У цьому фрагменті коду
\begin{verbatim}
template <class T> void f(T a1, T a2){}

void g(){
    int i=5;
    f(unique_ptr<int>(new int(++i)), unique_ptr<int>(new int(++i)));
\end{verbatim}

( 1 ) є невизначена поведінка%good

( 2 ) є стан гонки%bad

( 3 ) є гонка даних%bad

( 4 ) можливий виток ресурсів%bad
%enditem
}
%end
\end {large}

\vspace{4mm}
%end
\newpage
\end{document}

